\documentclass[12pt, a4paper]{article}

% ========== بسته‌های ضروری برای فارسی و ریاضی ==========
\usepackage[top=3cm, bottom=3cm, left=3cm, right=2.5cm]{geometry}
\usepackage{amsmath, amssymb}
\usepackage{graphicx}
\usepackage{multirow}
\usepackage{tabularx}
\usepackage{booktabs}
\usepackage{caption}
\usepackage{setspace}
\usepackage{xcolor}
\usepackage{hyperref}
\usepackage{fancyhdr}
\usepackage{titlesec}
\usepackage[stable]{footmisc}
\usepackage{pifont}

\usepackage{xepersian}


% ========== تنظیمات فونت برای فارسی و لاتین ==========

\settextfont[
  Path=fonts/,
  Extension=.ttf,
  UprightFont=*,
  BoldFont=*Bd,
  ItalicFont=*It,
  BoldItalicFont=*BdIt
]{XB-Niloofar}

\setdigitfont[
  Path=fonts/,
  Extension=.ttf,
  UprightFont=*,
  BoldFont=*Bd,
  ItalicFont=*It,
  BoldItalicFont=*BdIt
]{XB-Niloofar}


\setlatintextfont{TeX Gyre Termes} % فونت متن لاتین
\DefaultMathsDigits % استفاده از اعداد لاتین در محیط ریاضی
% ========== تنظیمات عنوان‌بندی بخش‌ها ==========
\titleformat{\section}
{\normalfont\large\bfseries\centering}{\thesection}{1em}{}
\titleformat{\subsection}
{\normalfont\normalsize\bfseries\centering}{\thesubsection}{1em}{}
\titleformat{\subsubsection}
{\normalfont\normalsize\bfseries}{\thesubsubsection}{1em}{}

% ========== تنظیمات سرصفحه و پاصفحه ==========
\pagestyle{fancy}
\fancyhf{}
\fancyhead[L]{\footnotesize \leftmark}
\fancyhead[R]{\footnotesize یک چارچوب مدل‌سازی کاهش‌یافته مبتنی بر توابع خاص آگاه از فیزیک ... }
\fancyfoot[C]{\thepage}
\renewcommand{\headrulewidth}{0.4pt}
\renewcommand{\footrulewidth}{0pt}

% ========== تنظیمات کلی ==========
\setlength{\parindent}{0pt}
\setlength{\parskip}{6pt}
\onehalfspacing
\captionsetup{labelsep=period, font=small}

\begin{document}
\title{یک چارچوب مدل‌سازی کاهش‌یافته مبتنی بر توابع خاص آگاه از فیزیک برای سیستم‌های تبدیلی گرماشیمیایی: اعتبارسنجی بر روی گازی‌سازی زیست‌توده}
\author{محسن درباز}
\date{ژانویه ۲۰۲۶}

\maketitle

\begin{abstract}
چارچوب مدل‌سازی کاهش‌یافته پارامتریشده با توابع خاص و آگاه از فیزیک ارائه می‌دهیم که تئوری کلاسیک توابع خاص را با کاهش ابعاد داده‌محور برای سیستم‌های با ابعاد بالا که توسط قوانین پایستگی اداره می‌شوند، پیوند می‌زند. این روش‌شناسی انتگرال‌های بیضوی کامل نوع اول — که از طریق الگوریتم میانگین حسابی-هندسی (AGM) محاسبه می‌شوند — را با استخراج ویژگی‌های هندسی، کدگذاری سلسله‌مراتبی توالی‌ها و آزمایشگاه‌های محاسباتی شبکه عصبی ادغام می‌کند تا بدون حل صریح معادلات دیفرانسیل جزئی (PDE)، به بسته‌بندی پایستگی دست یابد.

این چارچوب که بر روی گازی‌سازی زیست‌توده در راکتورهای بستر سیال توسعه داده شده و به شدت اعتبارسنجی شده، کارایی محاسباتی را همراه با دقت فیزیکی حفظ می‌کند. تحلیل بیش از ۲۰۰ مجموعه داده تجربی دربرگیرنده انواع مختلف سوخت (پسماندهای کشاورزی، گونه‌های چوبی، گیاهان انرژی) و شرایط عملیاتی (دما ۹۰۰-۶۰۰ درجه سلسیوس، نسبت هم‌ارزی ۰.۵-۰.۲، نسبت بخار به زیست‌توده ۱.۵-۰) عملکرد پیش‌بینی قوی‌ای ارائه می‌دهد: بازده تبدیل کربن (R2>0.85)، بازده گاز محصول (R2>0.82) و غلظت گونه‌های گازی مجزا (H2، CO، CO2، CH4؛ R2>0.78). نکته حائز اهمیت این است که چارچوب به روابط پیش‌بینی یکنوا در تمام گستره پارامترها دست می‌یابد – که بهینه‌های موضعی مشکل‌آفرین برای بهینه‌سازی مبتنی بر گرادیان را حذف می‌کند – در حالی که برای هر ارزیابی تنها به حدود ۵ تکرار AGM نیاز دارد که نشان‌دهنده شتابی در حدود 105−106 برابر نسبت به دینامیک سیالات محاسباتی است.

این کار جایگاهی منحصربه‌فرد در طیف تحلیلی-عددی اشغال می‌کند. برخلاف حل‌کننده‌های عددی سنتی PDE که معادلات حاکم را گسسته‌سازی می‌کنند یا مدل‌های جایگزین صرفاً داده‌محور که ساختار فیزیکی را نادیده می‌گیرند، روش ما هزینه محاسباتی را از انتگرال‌گیری تکراری PDE به ارزیابی کارای توابع خاص منتقل می‌کند. قوانین پایستگی به‌طور ساختاری از طریق فرمول‌بندی تابع انتقال H(s)=k∏(s−pi)∏(s−zi) تحمیل می‌شوند، جایی که انفعال، پایداری و علیت، تضمین‌کننده سازگاری فیزیکی در ساختار هستند. کاهش ابعاد از طریق تحلیل همبستگی متعارف و پارامترسازی هندسی، فضای ترکیب با ابعاد بالا (حدود ۱۰ متغیر) را به مجموعه‌ای فشرده از ناورداهای هندسی تبدیل می‌کند که – در ترکیب با معیار نوین «زمان محاسباتی ویژه» (sct) که صلبیت محاسباتی را ثبت می‌کند – به پارامترسازی انتگرال بیضوی تغذیه می‌کنند.

این روش‌شناسی نشان می‌دهد که توابع خاص لزوماً نباید مستقیماً از معادلات حاکم نشأت بگیرند تا مفید واقع شوند. در عوض، زمانی که کاهش ابعاد و تحمیل پایستگی از اهمیت بالایی برخوردارند، آن‌ها به عنوان «هسته‌های فیزیکی» کارا برای مدل‌های کاهش‌یافته پارامترشده عمل می‌کنند. چارچوب دو ناوردای جهانی مستقل را از توپولوژی تابع انتقال استخراج می‌کند: ξ (انحراف لگاریتمی پیوسته از خط پایه) و S (برابری ساختاری گسسته از پیکربندی قطب-صفر)، که هر دو همبستگی یکنوایی با کمیت‌های قابل مشاهده تجربی نشان می‌دهند. قابل توجه است که پارامتر بهره k در میان موارد مختلف از 10−47 تا 10−60 در نوسان است – میدانی بدون مقیاس که به صورت عددی، هنگام نرمالیزه شدن به صورت ξ=ln(ΛS/ΛG)، با مقیاس‌های ثابت کیهانی هم‌تراز است. این امر نشان می‌دهد که سیستم‌های حاکم از قانون پایستگی در مقیاس‌های مختلف ممکن است ساختارهای ابعادی بنیادین مشترکی داشته باشند که در پارامترسازی‌های تابع خاص کدگذاری شده‌اند.

اگرچه این چارچوب برای تبدیل گرماشیمیایی نشان داده شده است، ساختار ریاضی آن به طور طبیعی به سیستم‌هایی با قیود جمع به یک (کسرهای مولی، احتمالات، تخصیص منابع)، قوانین پایستگی (جرم، انرژی، آنتروپی) و ابعاد بالا که نیازمند استنتاج سریع هستند، تعمیم می‌یابد. با این حال، برای هر حوزه جدیدی اعتبارسنجی تجربی لازم است؛ شباهت‌های ساختاری برای تعمیم ضروری هستند اما کافی نیستند. این کار از طریق اعتبارسنجی گرماشیمیایی دقیق، قابلیت اجرا را اثبات می‌کند و نقشه‌ای برای تعمیم به شبکه‌های زیست‌شیمیایی، سیستم‌های بوم‌شناختی یا مدل‌سازی مواد ارائه می‌دهد که در آن‌ها ساختارهای ریاضی مشابهی بر پویایی‌ها حاکم هستند.
\end{abstract}

\keywords{روش‌های نیمه‌تحلیلی، انتگرال‌های بیضوی کامل، میانگین حسابی-هندسی، مدل‌سازی مرتبه کاهش‌یافته، قوانین پایستگی، توابع انتقال، گازی‌سازی زیست‌توده، بستر سیال، تبدیل گرماشیمیایی، تحلیل همبستگی متعارف، پارامترسازی هندسی، یکپارچه‌سازی داده‌های با ابعاد بالا}



\section{مقدمه}

\subsection{شکاف تحلیلی-عددی به‌عنوان انگیزه‌ای برای سیستم‌های قانون پایستگی}

سیستم‌های پیچیده‌ای که توسط قوانین پایستگی اداره می‌شوند—چه در جریان‌های چندفازی واکنشی، شبکه‌های بیوشیمیایی یا دینامیک بوم‌شناختی—یک چالش ریاضی بنیادین ارائه می‌دهند که دهه‌ها علوم محاسباتی را شکل داده است. معادلات حاکم معمولاً معادلات دیفرانسیل جزئی (PDE) غیرخطی هستند که هیچ‌گونه حل تحلیلی بسته‌شکلی ندارند و محققان را در یک دوگانگی مداوم قرار می‌دهند. از یک سو، حل‌کننده‌های عددی PDE، فضا و زمان را گسسته می‌کنند تا معادلات حاکم را به‌صورت تکراری حل کنند. اگرچه این رویکرد دقیق است و قادر به ثبت فیزیک جزئیات است، اما برای سناریوهایی که نیاز به کاوش فضای پارامتری با ابعاد بالا (هزاران تا میلیون‌ها متغیر)، ارزیابی‌های بلادرنگ یا تکرارشونده (بهینه‌سازی فرآیند، سیستم‌های کنترل، کمّی‌سازی عدم قطعیت)، مسائل معکوس با داده‌های پراکنده و پُرنویز (برآورد پارامتر، شناسایی سیستم) یا دینامیک چندمقیاسی با مقیاس‌های زمانی نامتشابه (ثانیه تا ساعت در یک شبیه‌سازی واحد) دارند، از نظر محاسباتی غیرممکن می‌شود. از سوی دیگر، مدل‌های کاملاً داده‌محور، سیستم را به‌عنوان یک جعبه سیاه در نظر می‌گیرند و نگاشت‌های ورودی-خروجی را از طریق یادگیری ماشین فرا می‌گیرند. اگرچه از نظر محاسباتی در طول استنتاج کارآمد هستند، اما این امر به بهای قربانی کردن تفسیرپذیری فیزیکی (ضرایب فاقد معنای مکانیکی)، تضمین اعمال قانون پایستگی (جرم، انرژی، تکانه ممکن است نقض شوند)، قابلیت برونیابی فراتر از توزیع‌های داده آموزشی و درک نظری که امکان طراحی منطقی یا عیب‌یابی را فراهم می‌کند، تمام می‌شود.

برای مورد خاص گازی‌سازی زیست‌توده در رآکتورهای بستر سیال—حوزه کاربردی اعتبارسنجی‌شده این کار—چالش حاد است. شبیه‌سازی‌های کامل دینامیک سیالات محاسباتی (CFD) که تعاملات گاز-جامد، واکنش‌های شیمیایی، انتقال حرارت و جرم در مقیاس‌های نامتشابه را ردیابی می‌کنند، به $10^{7}$ تا $10^{9}$ عمل برای هر ارزیابی مجموعه پارامتری نیاز دارند که کنترل بلادرنگ یا بهینه‌سازی جامع را غیرممکن می‌سازند. در مقابل، همبستگی‌های کاملاً تجربی، قادر به ثبت جفت‌شدگی غیرخطی بین ترکیب سوخت، شرایط عملیاتی و توزیع محصولات نیستند، به‌ویژه هنگام برونیابی به خوراک‌های جدید یا پیکربندی‌های رآکتور.

\subsection{پارامترسازی تابع خاص فیزیک‌آگاه به‌عنوان راه سوم}

این کار یک راه میانه را نشان می‌دهد که نقاط قوت هر دو پارادایم را ترکیب می‌کند: استفاده از توابع خاص شناخته‌شده به‌عنوان "هسته‌های فیزیکی" کارآمد در یک چارچوب پارامترسازی داده‌محور. به‌طور خاص، ما از انتگرال‌های بیضوی کامل نوع اول K(m)\_K\_(\_m\_) استفاده می‌کنیم که از طریق میانگین حسابی-هندسی (AGM) محاسبه می‌شوند، تا یک مدل کاهش‌یافته بسازیم که به‌طور همزمان قوانین پایستگی را از طریق فرمول‌بندی تابع انتقال (جایی که انفعال، پایداری و علیت، سازگاری فیزیکی را از طریق ساختار تضمین می‌کنند) رعایت می‌کند، کارایی محاسباتی را از طریق AGM با همگرایی درجه‌دوم (∼5∼5 تکرار برای دقت ماشین) به دست می‌آورد.

\subsection{مکانیک سماوی به‌عنوان الهام تاریخی}

رویکرد ما به‌طور صریح از یک پارادایم کلاسیک در مکانیک سماوی الهام می‌گیرد که نقش توابع خاص را در سیستم‌های فیزیکی روشن می‌سازد. مسئله دو-جسم یک راه‌حل تحلیلی دقیق قابل بیان برحسب انتگرال‌ها و توابع بیضوی می‌پذیرد و محاسبات دقیق مدار سیاره‌ای را ممکن می‌سازد. در مقابل، مسئله \en(N)-جسم (N≥3\_N\_≥3) هیچ راه‌حل تحلیلی کلی ندارد و نیاز به انتگرال‌گیری عددی معادلات نیوتن دارد. با این حال، حتی برای مسئله \en(N)-جسم، وقتی تقارن‌ها یا قیود وجود دارند (مانند مقیاس‌های جرمی سلسله‌مراتبی یا مدارهای تقریباً دایره‌ای)، بسط‌های اغتشاش شامل توابع خاص، توصیف‌های کاهش‌یافته ارزشمندی ارائه می‌دهند که امکان تحلیل پایداری بلندمدت و طراحی مسیر را فراهم می‌کنند. مشاهده کلیدی این است که توابع خاص، ساختار اساسی تحمیل‌شده توسط قوانین پایستگی (انرژی، تکانه زاویه‌ای) را ثبت می‌کنند حتی زمانی که راه‌حل‌های تحلیلی کامل در دسترس نیستند.

به‌طور مشابه، در حالی که دینامیک سیالات محاسباتی چندفازی واکنشی کامل برای گازی‌سازی از نظر تحلیلی غیرقابل حل باقی می‌ماند، ما نشان می‌دهیم که انتگرال‌های بیضوی می‌توانند رفتار اساسی سیستم را ثبت کنند وقتی مهندسی ویژگی داده‌محور، ابعاد را کاهش می‌دهد در حالی که معنای فیزیکی را حفظ می‌کند، قیود پایستگی به‌طور ساختاری (نه الگوریتمی پس‌ین) اعمال می‌شوند و توابع خاص، دینامیک کاهش‌یافته را با ضرایب تفسیرپذیر پارامترسازی می‌کنند. این تشابه بیش از استعاری است. در هر دو مکانیک سماوی و تبدیل ترموشیمیایی، قوانین پایستگی، حالت‌های قابل دسترس را به منیفولدهای با ابعاد پایین‌تر در فضای فاز محدود می‌کنند. توابع خاص به‌طور طبیعی به‌عنوان سیستم‌های مختصات روی این منیفولدها ظهور می‌کنند و مسیرهایی را که قیود را رعایت می‌کنند، پارامترسازی می‌کنند. سهم ما نشان دادن این اصل است که فراتر از سیستم‌هایی با هامیلتونی‌های شناخته‌شده به سیستم‌های قانون پایستگی مشخص‌شده تجربی گسترش می‌یابد.

\subsection{گازی‌سازی زیست‌توده به‌عنوان بستر آزمایش اعتبارسنجی}

گازی‌سازی زیست‌توده، خوراک کربن‌دار جامد را از طریق اکسایش جزئی در دماهای بالا (۹۰۰–۶۰۰ درجه سلسیوس) به مخلوط‌های گاز قابل اشتعال (عمدتاً \ce(H2)، \ce(CO)، \ce(CO2)، \ce(CH4)) تبدیل می‌کند. رآکتورهای بستر سیال—جایی که ذرات زیست‌توده در عامل گازی‌ساز جریان بالارونده (هوا، اکسیژن، بخار یا مخلوط‌ها) معلق می‌شوند—مزایایی از جمله انعطاف‌پذیری سوخت، توزیع دمای یکنواخت و نرخ‌های انتقال حرارت بالا ارائه می‌دهند. با این حال، فیزیک حاکم، هیدرودینامیک چندفازی (تبادل تکانه گاز-جامد، تشکیل و پیوست‌شدن حباب، به‌دام‌افتادن ذرات)، انتقال حرارت (تابش، جابه‌جایی، رسانش در فازهای با خواص حرارتی نامتشابه)، انتقال جرم (نفوذ گاز از طریق منافذ ذرات، انتقال گونه‌های بین‌فازی) و سینتیک شیمیایی (خشک‌شدن، پیرولیز شامل صدها واکنش موازی، گازی‌سازی از طریق واکنش‌های زغال-گاز، شکستن قطران و تعادل‌های تبدیل آب-گاز) را به هم پیوند می‌دهد. مدل‌سازی دقیق اصول اولیه نیاز به حل معادلات پیوستگی، تکانه، انرژی و گونه برای هر دو فاز دارد که روی شبکه‌های فضایی گسسته‌شده‌اند تا پدیده‌های در مقیاس ذره را حل کنند—از نظر محاسباتی برای بهینه‌سازی طراحی یا کاربردهای کنترل غیرممکن است.

ما بیش از ۲۰۰ مجموعه داده تجربی از ۱۵ منبع مستقل ادبیات را جمع‌آوری کردیم که انواع رآکتور بستر سیال جوشان (BFB) و بستر سیال در گردش (CFB)، ۳۸ نوع زیست‌توده متمایز از جمله بقایای کشاورزی (پوست برنج، کاه گندم، ساقه ذرت)، گونه‌های چوبی (کاج، راش، اکالیپتوس) و گیاهان انرژی‌زا (سویچ‌گراس، ميسکانتوس) را پوشش می‌دهند. پارامترهای عملیاتی شامل دما (\en(T))، نسبت هم‌ارزی (\en(ER))، نسبت بخار به زیست‌توده (\en(SBR)) و ترکیب سوخت (عنصری، تقریبی، ساختاری) و همچنین مشاهده‌پذیرهایی از جمله ترکیب گاز محصول (H2,CO,CO2,CH4,N2\_H\_2,\_CO\_,\_CO\_2,\_CH\_4,\_N\_2)، بازده (\si(\cubic\metre\per\kilogram) زیست‌توده) و بازده تبدیل کربن (\en(CCE)\\%) است. این غنای مجموعه داده—که سوخت‌ها، شرایط و پیکربندی‌های رآکتور متنوع را پوشش می‌دهد—اعتبارسنجی دقیقی برای هر چارچوب مدل‌سازی پیشنهادی فراهم می‌کند.

\subsection{قابلیت گسترش از طریق تشابه‌های ساختاری در حوزه‌های وسیع‌تر}

اگرچه این مقاله به‌طور انحصاری بر نتایج اعتبارسنجی‌شده گازی‌سازی تمرکز دارد، اما ساختار ریاضی چارچوب، قابلیت گسترش به حوزه‌های دیگر با ویژگی‌های کلیدی مشترک را نشان می‌دهد. سیستم‌هایی که متغیرهای آن‌ها باید قیود جمع‌به-یک (∑yi=1∑\_y\_\_i\_=1) را برآورده کنند (قابل اعمال به کسرهای مولی، احتمالات، تخصیص‌های پرتفوی)، قوانین پایستگی (جرم، انرژی، آنتروپی یا سایر کمیت‌های پایستار که دینامیک را محدود می‌کنند)، ابعاد بالا (هزاران تا میلیون‌ها متغیر نیازمند کاهش ابعاد)، دینامیک چندمقیاسی (مقیاس‌های زمانی یا فضایی نامتشابه در یک سیستم واحد) یا نیازمندی‌های بلادرنگ (بهینه‌سازی، کنترل یا استنتاج تحت محدودیت‌های زمانی سخت) ممکن است از این رویکرد بهره‌مند شوند. حوزه‌های کاربردی بالقوه شامل شبکه‌های بیوشیمیایی (پایستگی متابولیت، تعادل ATP/ردوکس)، سیستم‌های بوم‌شناختی (چرخه مواد مغذی، دینامیک جمعیت)، مدل‌های اقتصادی (جریان‌های مالی، تعادل بازار) و مدل‌سازی مواد (چندمقیاسی از اتمی تا پیوسته) است. با این حال، تأکید می‌کنیم که برای هر حوزه جدید، اعتبارسنجی تجربی لازم است؛ تشابه‌های ساختاری ریاضی برای تعمیم ضروری هستند اما کافی نیستند. این کار امکان‌پذیری چارچوب را از طریق اعتبارسنجی ترموشیمیایی دقیق اثبات می‌کند.

\subsection{سازمان‌دهی سند و سهم‌های کلیدی}

باقی‌مانده این مقاله به شرح زیر ارائه می‌شود. بخش ۲ روش‌شناسی ما را در طیف تحلیلی-عددی جای‌دهی می‌کند و ادبیات مرتبط را مرور می‌کند و زمینه را برای رویکرد ترکیبی ما ایجاد می‌کند. بخش ۳ روش‌شناسی عمومی را ارائه می‌دهد و تحلیل همبستگی متعارف، تبدیل هندسی، آزمایشگاه‌های محاسباتی شبکه عصبی، تشکیل ماتریس مشخصه و پارامترسازی انتگرال بیضوی را به تفصیل شرح می‌دهد. بخش ۴ این روش‌ها را روی گازی‌سازی زیست‌توده اعمال می‌کند و مجموعه داده، جزئیات مهندسی ویژگی، جزئیات پیاده‌سازی و معادلات پایستگی را توصیف می‌کند. بخش ۵ نتایج اعتبارسنجی و بحث از جمله معیارهای عملکرد کمّی، تفسیر فیزیکی ناورداها و مقایسه با مدل‌های موجود را ارائه می‌دهد. بخش ۶ به اختصار در مورد قابلیت گسترش و پیامدهای وسیع‌تر بحث می‌کند. بخش ۷ با محدودیت‌ها و جهت‌های آینده به نتیجه می‌رسد.

سهم‌های کلیدی ما شامل پنج پیشرفت است. اول، نشان می‌دهیم که توابع خاص می‌توانند به‌عنوان پارامترسازی‌های کارآمد مدل‌های کاهش‌یافته یادگرفته‌شده از داده عمل کنند، حتی زمانی که مستقیماً از معادلات حاکم نشأت نمی‌گیرند. دوم، اعتبارسنجی دقیقی روی بیش از ۲۰۰ مجموعه داده گازی‌سازی ارائه می‌دهیم که به $R^{2}>0.85$ برای همه خروجی‌ها دست می‌یابد در حالی که بست پایستگی و کارایی محاسباتی (شتاب $10^{5}$ تا $10^{6}$ برابری در مقابل CFD) حفظ می‌شود. سوم، "زمان محاسباتی ویژه" (\en(sct)) را به‌عنوان یک ویژگی دینامیک جدید معرفی می‌کنیم که صلبیت محاسباتی غیرتعادلی را ثبت می‌کند. چهارم، کشف می‌کنیم که بهره تابع انتقال از $10^{-47}$ تا $10^{-60}$ متغیر است و یک میدان بی‌بعد ایجاد می‌کند که از نظر عددی با مقیاس‌های ثابت کیهانی هم‌تراز است. پنجم، یک نقشه روش‌شناختی برای گسترش چارچوب به سایر سیستم‌های اداره‌شده توسط قانون پایستگی ایجاد می‌کنیم.

\section{مرور ادبیات و جایگاه‌یابی روش‌شناختی}

\subsection{توصیف طیف رویکردهای مدل‌سازی}

برای قراردادن صحیح مشارکت ما در زمینه، ابتدا منظره رویکردها برای سیستم‌های پیچیده قانون پایستگی را توصیف می‌کنیم، با این شناسایی که هر روش‌شناسی باید تعادل بین خواسته‌های رقابتی دقت، کارایی، تفسیرپذیری و قابلیت تعمیم را برقرار کند. جدول ۱ این جایگاه‌یابی را در پنج پارادایم تثبیت‌شده خلاصه می‌کند و کار ما را درون این طیف قرار می‌دهد.

\begin{table}[t]
\centering
\small
\caption{جایگاه‌یابی روش‌شناختی در طیف تحلیلی-عددی}
\begin{tabularx}{\linewidth}{@{}*{5}{>{\raggedleft\arraybackslash}X}@{}}
\toprule
رویکرد & پایستگی & کارایی & تفسیرپذیری & برون‌یابی \\
\midrule
\textbf(تحلیلی (فرم بسته)) & عالی & عالی & عالی & عالی \\
\textbf(روش‌های اغتشاشی) & خوب & خوب & متوسط & ضعیف \\
\textbf(PDE عددی (CFD)) & عالی & ضعیف & خوب & خوب \\
\textbf(مدل‌های مرتبه کاهش‌یافته) & خوب & خوب & متوسط & متوسط \\
\textbf(صرفاً داده‌محور (ML)) & ضعیف & خوب & ضعیف & ضعیف \\
\textbf(این کار) & عالی & خوب & خوب & خوب \\
\bottomrule
\end{tabularx}
\end{table}


\textbf(راهنما:) عالی (عالی)، خوب (خوب)، متوسط (متوسط)، ضعیف (ضعیف)

رویکردهای تحلیلی خالص – حل مستقیم معادلات دیفرانسیل جزئی حاکم به فرم بسته – آرمانی را نشان می‌دهند و دقت، تفسیرپذیری، ارزیابی کارا و اعتبار برای تمام مقادیر پارامتر درون دامنه مدل را ارائه می‌دهند. مثال‌های کلاسیک شامل معادله حرارت با شرایط مرزی ساده، جریان استوکس یا جریان پتانسیل حول هندسه‌های ساده هستند. با این حال، برای سیستم‌های چندفاز واکنشی غیرخطی مانند گازی‌سازی، جواب‌های تحلیلی وجود ندارند. معادلات حاکم، دینامیک سیالات آشفته، سینتیک شیمیایی ناهمگن، پدیده‌های انتقال بین فازی و دینامیک ذرات را در طول و مقیاس‌های زمانی متفاوت به هم پیوند می‌زنند – حل‌ناپذیری ریاضی اساسی است، نه صرفاً یک مانع فنی.

هنگامی که معادلات حاکم تقریباً با یک حالت پایه قابل حل مطابقت دارند، بسط‌های اغتشاشی (منظم، تکین، چندمقیاسی) می‌توانند جواب‌های تحلیلی تقریبی ارائه دهند. با این حال، این‌ها نیازمند وجود یک پارامتر کوچک کنترل‌کننده غیرخطی، همواری کافی برای تضمین همگرایی و یکنواختی اعتبار در گستره پارامترها هستند. برای گازی‌سازی، چنین پارامتر کوچک جهانی وجود ندارد. رژیم‌های عملیاتی از رفتار کنترل‌شده با سینتیک (دمای پایین) تا رفتار محدودشده با انتقال (دمای بالا) را در برمی‌گیرند بدون نقطه بسط طبیعی. در نتیجه، روش‌های اغتشاشی فایده محدودی ارائه می‌دهند مگر در همسایگی‌های پارامتری باریک.

CFD با وفاداری بالا معادلات حاکم (معمولاً چارچوب‌های اولری-اولری یا اولری-لاگرانژی) را گسسته‌سازی می‌کند و به صورت تکراری حل می‌کند. در حالی که قادر به ثبت فیزیک جزئیات است، هزینه محاسباتی از طریق گسسته‌سازی فضایی بد مقیاس می‌شود. گسسته‌سازی فضایی نیازمند O(Nx×Ny×Nz) سلول شبکه برای حل مقیاس ذرات (حدود 105 تا 106 سلول) است. انتگرال‌گیری زمانی محدودیت‌های پایداری کوران-فریدریش-لوی (CFL) را برای روش‌های صریح تحمیل می‌کند یا نیازمند حل سیستم خطی تکراری برای روش‌های ضمنی است. بسته شدن آشفته از طریق شبیه‌سازی گردابه بزرگ (LES) یا معادلات ناویر-استوکس میانگین‌گیری رینولدز (RANS) تجربی‌گرایی را معرفی می‌کند با وجود هزینه محاسباتی بالا. سینتیک شیمیایی معادلات دیفرانسیل معمولی سخت‌افزاری را ارائه می‌دهد که نیازمند حل‌کننده‌های تخصصی با مکانیسم‌های دقیق شامل صدها گونه است که همچنان ممنوع‌کننده است. یک شبیه‌سازی CFD منفرد برای یک مجموعه پارامتر نیازمند ساعتها تا روزها روی خوشه‌های محاسباتی با کارایی بالاست. برای بهینه‌سازی طراحی که هزار تا ده هزار ترکیب پارامتر را کاوش می‌کند، یا کنترل بلادرنگ که نیازمند پاسخ میلی‌ثانیه است، CFD غیرممکن است.

رویکردهای سنتی مدل مرتبه کاهش‌یافته (ROM) – تجزیه متعامد مناسب (POD)، تجزیه حالت دینامیکی (DMD)، قطع متوازن – جواب‌های PDE با ابعاد بالا را بر روی زیرفضاهای با ابعاد پایین‌تر پوشیده شده با مدهای غالب فرافکنی می‌کنند. در حالی که برای سیستم‌های خطی یا ضعیفاً غیرخطی مؤثرند، غیرخطی‌های قوی گازی‌سازی (سینتیک آرنیوس، تغییر فاز) و مرزهای متحرک (انقباض زغال) قابلیت کاربرد ROM را محدود می‌کنند. علاوه بر این، ساخت ROMها نیازمند عکس‌های با وفاداری بالا از CFD است که گلوگاه محاسباتی را در طول آموزش به ارث می‌برد.

یادگیری ماشین – شبکه‌های عصبی، فرآیندهای گاوسی، رگرسیون نمادین – سیستم را به عنوان یک جعبه سیاه در نظر می‌گیرد و نگاشت‌های ورودی-خروجی را از داده می‌آموزد. یادگیری عمیق مدرن موفقیت تجربی قابل توجهی کسب می‌کند، اما محدودیت‌های بنیادینی باقی می‌مانند: پیش‌بینی‌ها ممکن است موازنه جرم یا انرژی را نقض کنند مگر اینکه صریحاً محدود شوند (برای معماری‌های پیچیده دشوار است)، عملکرد به شدت خارج از توزیع داده‌های آموزشی تخریب می‌شود (برون‌یابی ضعیف)، میلیون‌ها پارامتر فاقد معنی فیزیکی هستند و عیب‌یابی یا طراحی منطقی را مختل می‌کنند (تفسیرپذیری کدر) و نیازهای قابل توجه داده برچسب‌دار برای سیستم‌های تجربی پرهزینه مشکل‌زا هستند.

\subsection{جایگاه‌یابی مدل مرتبه کاهش‌یافته پارامترشده با تابع خاص آگاه از فیزیک}

این کار عناصری از پارادایم‌های چندگانه را ترکیب می‌کند تا یک جایگاه روش‌شناختی منحصربه‌فرد اشغال کند. از روش‌های تحلیلی، از توابع خاص (انتگرال‌های بیضوی) با خواص به‌خوبی درک‌شده استفاده می‌کنیم که از طریق الگوریتم‌های عددی پایدار (AGM) محاسبه می‌شوند. از تئوری تابع انتقال، پایستگی را به‌طور ساختاری از طریق انفعال، پایداری و علیت فرمول‌بندی H(s) تحمیل می‌کنیم. از مدل‌سازی داده‌محور، کاهش ابعاد (CCA) و مهندسی ویژگی را از داده تجربی بدون فرض شکل‌های تابعی می‌آموزیم. از فلسفه ROM، مشاهدات با ابعاد بالا را بر روی ناورداهای با ابعاد پایین ثبت‌کننده پویایی‌های ضروری فرافکنی می‌کنیم.

به طور بحرانی، ما نیازی نداریم که توابع خاص از معادلات حاکم نشأت بگیرند. در عوض، آن‌ها به عنوان پارامترسازی‌های کارای پویایی‌های کاهش‌یافته یادگرفته‌شده عمل می‌کنند. این تغییر فلسفی – توابع خاص به عنوان «هسته‌های فیزیکی» برای مدل‌های داده‌محور – کاربرد در سیستم‌هایی را ممکن می‌سازد که مشتق اصول اول در آن‌ها حل‌ناپذیر است اما قیود پایستگی شناخته شده هستند.

\subsection{کارهای مرتبط در مدل‌سازی گازی‌سازی}

ادبیات مدل‌سازی گازی‌سازی زیست‌توده دهه‌ها را در برمی‌گیرد و رویکردها در سه دسته قرار می‌گیرند که هر کدام جنبه‌های متفاوتی از فضای مسئله را روشن می‌کنند. مدل‌های تعادلی نرخ واکنش بی‌نهایت را فرض می‌کنند و ترکیب محصول را با کمینه‌سازی انرژی آزاد گیبس تحت موازنه‌های عنصری پیش‌بینی می‌کنند. اگرچه ارزیابی سریعی دارند، اما تبدیل را بیش‌برآورد می‌کنند و قادر به ثبت محدودیت‌های سینتیکی یا گرادیان‌های دما نیستند. انواع آن‌ها دماهای شبه‌تعادلی معرفی می‌کنند یا تعادل را به واکنش‌های خاصی محدود می‌کنند و دقت را به بهای معرفی پارامترهای تجربی با معنی فیزیکی مبهم بهبود می‌بخشند.

مدل‌های سینتیکی شبکه‌های واکنش شیمیایی (ده‌ها تا صدها واکنش) را با پدیده‌های انتقال پیوند می‌زنند. مدل‌های مبتنی بر اصول اول نیازمند پارامترهای سینتیکی دقیقی هستند که اغلب برای ماتریس‌های پیچیده زیست‌توده در دسترس نیستند، که ساختار ناهمگن آن‌ها از ثابت‌های نرخ تک‌مقداری سر باز می‌زند. مدل‌های سینتیکی تجربی از واکنش‌های کلی با پارامترهای آرنیوس برازش‌شده استفاده می‌کنند که قابلیت محاسباتی را بهبود می‌بخشند اما تعمیم‌پذیری بین خوراک‌ها و شرایط عملیاتی را قربانی می‌کنند – مدلی که برای خاک‌اره کاج تنظیم شده ممکن است برای پوسته برنج به دلیل شیمی خاکستر و خواص ساختاری متفاوت به شدت شکست بخورد.

مدل‌های شبکه عصبی نگاشت‌ها را مستقیماً از داده تجربی می‌آموزند. در حالی که درون توزیع‌های آموزشی به درون‌یابی خوبی دست می‌یابند، بیشتر آن‌ها فاقد ساختار فیزیکی هستند – پیش‌بینی‌ها ممکن است استوکیومتری یا موازنه انرژی را نقض کنند و عملکرد برون‌یابی ضعیف است. برخی کارهای اخیر پایستگی را به عنوان قیود نرم در تابع زیان گنجانده‌اند که سازگاری فیزیکی را بهبود می‌بخشد اما آن را به طور دقیق تضمین نمی‌کند. هنگامی که موازنه جرم به عنوان یک جمله جریمه در برابر دقت پیش‌بینی وزن‌دهی می‌شود، بهینه‌سازی ممکن است سازگاری فیزیکی را با برازش تجربی معامله کند، به ویژه در مناطق با داده اندک.

یک مطالعه تطبیقی توسط مشتقی و همکاران سه مدل برتر گازی‌سازی (شبه‌تعادلی، سینتیکی تجربی، سینتیکی دقیق) را در برابر مجموعه‌های داده مستقل ارزیابی کرد و خطای متوسط ۳۰-۱۵٪ برای ترکیب گاز محصول در شرایط دیده‌نشده یافت. چارچوب ما، که بر روی مجموعه‌های آزمایشی مشابه اعتبارسنجی شده، به خطای متوسط <10\% دست می‌یابد در حالی که بسته بودن پایستگی و ارزیابی 105× سریع‌تر را حفظ می‌کند – ترکیبی که کاربردهای قبلاً غیرممکن را ممکن می‌سازد.

\subsection{تحلیل همبستگی متعارف برای یکپارچه‌سازی چندگروهی}

مشخصه‌سازی خوراک زیست‌توده به سه گروه متمایز تقسیم می‌شود که هر کدام چشم‌اندازهای مکمل اما غیرقابل قیاسی ارائه می‌دهند که باید برای مدل‌سازی جامع آشتی داده شوند. ترکیب عنصری (کربن، هیدروژن، اکسیژن، نیتروژن، گوگرد) به طور گسترده اندازه‌گیری شده و مستقیماً به استوکیومتری مرتبط است اما فاقد اطلاعات ساختاری درباره معماری مولکولی است. آنالیز تقریبی (مواد فرار، کربن ثابت، خاکستر، رطوبت) رفتار تجزیه گرمایی قابل مشاهده از طریق روش‌های استاندارد را ثبت می‌کند اما ساختار مولکولی و کارکرد شیمیایی را مخفی می‌کند. اجزای ساختاری (سلولز، همی‌سلولز، لیگنین) بینش مکانیکی به مسیرهای پیرولیز و مکانیسم‌های واکنش ارائه می‌دهند اما نیازمند روش‌های اندازه‌گیری‌فشرده هستند و کمتر گزارش شده‌اند.

این گروه‌ها همبستگی‌های درونی (مثلاً محتوای کربن بالا با کربن ثابت بالا همبستگی دارد) و روابط بین گروهی (محتوای لیگنین از طریق ساختار آروماتیک خود بر نسبت اکسیژن به کربن عنصری تأثیر می‌گذارد) نشان می‌دهند، اما الحاق مستقیم آن‌ها به یک بردار ویژگی واحد مشکل‌زا است. واحدها و مقیاس‌های مختلف (درصد جرم، درصد مول) ناخوش‌تعینی عددی در الگوریتم‌های پایین‌دستی ایجاد می‌کنند. افزونگی ابعاد را بدون افزودن اطلاعات افزایش می‌دهد و اصل پارسیمونی ضروری برای تعمیم را نقض می‌کند. معنی فیزیکی زمانی مخفی می‌شود که کمیت‌های ناهمگن به طور یکنواخت رفتار شوند – میانگین گرفتن یک کسر جرمی با یک نسبت مولی عددی تولید می‌کند که فاقد تفسیر فیزیکی است.

CCA این چالش‌ها را با شناسایی همبستگی‌های حداکثری بین گروه‌ها در حالی که ساختار درون گروهی را از طریق یک چارچوب آماری اصولی حفظ می‌کند، حل می‌کند. با توجه به دو مجموعه چندمتغیره اندازه‌گیری شده روی نمونه‌های یکسان:

\begin{equation}
\mathbf{X}=[X_{1},X_{2},\ldots,X_{p}],\qquad \mathbf{Y}=[Y_{1},Y_{2},\ldots,Y_{q}]
\end{equation}

CCA ترکیبات خطی زیر را جست‌جو می‌کند:

\begin{equation}
U_{1}V_{1}=a_{11}X_{1}+a_{12}X_{2}+\cdots+a_{1}pXp=a_{1}^{\top}X=b_{11}Y_{1}+b_{12}Y_{2}+\cdots+b_{1}qYq=b_{1}^{\top}Y\tag{1}
\end{equation}

که همبستگی متعارف ρ1∗=Corr(U1,V1) را تحت شرایط Var(U1)=Var(V1)=1 بیشینه می‌کند. جفت‌های بعدی (Uk,Vk) برای k=2,3,… به صورت تکراری یافت می‌شوند، هر کدام متعامد با جفت‌های قبلی، که بزرگی‌های کاهش‌یافته واریانس مشترک را در یک سلسله مراتب احترام‌گذار به ساختار کوواریانس ذاتی داده استخراج می‌کنند.

\subsubsection{معماری دو مرحله‌ای با پیش‌شرط‌سازی دو لگاریتمی}

ما از یک پیکربندی گراف مینیمال استفاده می‌کنیم که اتصال کامل بین هر سه گروه خصوصیت را تضمین می‌کند – توپولوژی لازم و کافی برای یکپارچه‌سازی کامل. مرحله ۱ گروه‌های عنصری و تقریبی را از طریق پایه کربن مشترکشان جفت می‌کند. نسبت‌های عنصری به کربن نرمال می‌شوند: rE=[C/O,C/H,C/N,C/S]⊤ و نسبت‌های تقریبی به کربن ثابت: rP=[FC/VM,FC/A]⊤. هر دو مجموعه تحت پیش‌شرط‌سازی دو لگاریتمی قرار می‌گیرند:

\begin{equation}
x_{E}=\ln(\ln r_{E}),\qquad x_{P}=\ln(\ln r_{P})
\end{equation}

CCA اعمال‌شده بر این مجموعه‌های پیش‌شرط شده، اولین جفت متعارف (E1,P1) را که همبستگی ρ(E1,P1) را بیشینه می‌کند، نتیجه می‌دهد.

مرحله ۲ گروه‌های عنصری و ساختاری را جفت می‌کند. نسبت‌های ساختاری به همی‌سلولز نرمال می‌شوند: rS=[Hem/Cel,Hem/Lig]⊤ و به طور یکسان پیش‌شرط می‌شوند:

\begin{equation}
x_{S}=\ln(\ln r_{S})
\end{equation}

CCA دومین جفت متعارف (E2,S1) را که همبستگی ρ(E2,S1) را بیشینه می‌کند، نتیجه می‌دهد. این معماری تضمین می‌کند که ترکیب عنصری به عنوان هسته مرکزی عمل می‌کند که اطلاعات تقریبی و ساختاری را به هم متصل می‌کند در حالی که استقلال آماری بین دو جفت متعارف حفظ می‌شود.

\subsubsection{تبدیل دو لگاریتمی به عنوان نگاشت دیفئومورفیسم}

تبدیل T:x↦ln[ln(x)] و معکوس آن T−1:y↦exp[exp(y)] یک نگاشت دوطرفه هموار تشکیل می‌دهند که یک دیفئومورفیسم بین فضای نسبت و خط حقیقی برقرار می‌کند:

R>1lnR>0lnRexpR>0expR>1

این ترکیب دورانی به طور همزمان به سه هدف آماری دست می‌یابد. اول، تثبیت توزیعی به چالش بنیادین می‌پردازد که نسبت‌های زیست‌توده در گستره‌ای از مرتبه‌های بزرگی قرار دارند – نسبت‌های کربن به اکسیژن از C/O∈[0.67,15] در حالی که نسبت‌های سلولز به لیگنین از Cel/Lig∈[0.3,6.2] در نوسان هستند. تبدیل T این گستره‌های شدید را به متغیرهای تقریباً گاوسی بر روی R فشرده می‌کند، فرضیات ضمنی نرمال بودن CCA را ارضا می‌کند و از سلطه داده‌های پرت در محاسبات کوواریانس جلوگیری می‌کند.

دوم، هم‌ریختی گروه-نظری روابط ضربی در فضای نسبت را به روابط جمعی در فضای تبدیل‌شده تبدیل می‌کند. تحت T، حاصلضرب دو نسبت تقریباً به صورت زیر تبدیل می‌شود:

\begin{equation}
T(r_{1}\cdot r_{2})=\ln\!\bigl(\ln(r_{1}r_{2})\bigr)=\ln\!\bigl(\ln(r_{1})+\ln(r_{2})\bigr)\approx T(r_{1})+T(r_{2})
\end{equation}

برای r1,r2≫1. این تقریب برقرار است زیرا ln(r)≫1 وقتی r به طور قابل توجهی بزرگتر از واحد است، و اجازه می‌دهد لگاریتم یک جمع با خود جمع تقریب زده شود در تبدیل بیرونی. تناظر جبری حاصل، مقایسه مستقیم بین انواع کمیت ناهمگن را ممکن می‌سازد که در غیر این صورت فاقد مقیاس مشترک خواهند بود.

سوم، همگن‌سازی واریانس مقادیر شدید را از طریق ژاکوبین تبدیل تضعیف می‌کند:

\begin{equation}
\frac{d x}{dT}=\frac{1}{x\ln(x)}
\end{equation}

این مشتق با افزایش x به سرعت کاهش می‌یابد، دنباله‌های توزیع‌های چولگون را فشرده می‌کند و الگوهای همبستگی ناوردای تبدیل را نتیجه می‌دهد که برای استخراج متغیرهای متعارف بدون اریب ضروری است. اثر ترکیبی، ناخوش‌تعینی عددی در مسئله مقدار ویژه CCA بعدی را تثبیت می‌کند در حالی که ترتیب‌های نسبی که روابط فیزیکی را کد می‌کنند حفظ می‌شوند.

خروجی شامل چهار متغیر متعارف است که نمایش فشرده زیر را تشکیل می‌دهند:

\begin{equation}
\mathrm{EPS}=[E_{1},E_{2},S_{1},P_{1}]\in\mathbb{R}^{4}
\end{equation}

این تعبیه چهاربعدی، تغییرپذیری ضروری زیست‌توده را ثبت می‌کند در حالی که اطلاعات عنصری، تقریبی و ساختاری را از طریق بیشینه‌سازی کوواریانس داده‌محور به جای وزن‌دهی‌های دلخواه خبره یکپارچه می‌کند. مختصات EPS به عنوان ورودی‌های اصولی آماری به معماری‌های یادگیری ماشین بعدی عمل می‌کنند، روابط بین حوزه‌ای را حفظ می‌کنند و در عین حال تفسیرپذیری را از طریق بارهای متعارف وابسته به گروه خاص حفظ می‌کنند.

\subsection{تبدیل هندسی: گشودن فضای تحلیل «آب‌های آزاد»}

داده‌های ترکیب – خواه کسرهای مولی، کسرهای جرمی یا احتمالات – از قیود جمع به یک بیان‌شده توسط قانون دالتون تبعیت می‌کنند:

\begin{equation}
\sum_{i=1}^{n} y_{i}=1
\end{equation}

در مختصات دکارتی، این قید یک ابرصفحه (n−1) بعدی (سیمپلکس) تعبیه‌شده در Rn تعریف می‌کند. اگرچه از نظر ریاضی معتبر است، این نمایش از سه نارسایی بحرانی رنج می‌برد. اول، مرزهای شکننده در لبه‌ها و رئوس سیمپلکس ظهور می‌یابند جایی که ژاکوبین‌های نزدیک به تکین، ناخوش‌تعینی عددی در الگوریتم‌های بهینه‌سازی ایجاد می‌کنند – روش‌های مبتنی بر گرادیان هنگام نزدیک شدن به محدودیت‌های ترکیب که در آن یک یا چند جزء ناپدید می‌شوند، مشکل دارند. دوم، جعبه ابزار هندسی محدود موجود در مختصات سیمپلکس به این معنی است که فواصل اقلیدسی فاقد تفسیرپذیری فیزیکی هستند، زیرا فاصله ظاهری بین دو ترکیب به انتخاب دلخواه دستگاه مختصات به جای ناهمگونی ذاتی ترکیب بستگی دارد. سوم، تورم همخطی‌گری از وابستگی ذاتی متغیر ناشی می‌شود، که عدد شرایط را در مسائل رگرسیون و وارون افزایش می‌دهد، جایی که فرضیات استاندارد استقلال پیش‌بین به شدت شکست می‌خورند.

\subsubsection{پارامترسازی مثلثاتی به عنوان هم‌ریختی ساختاری}

پارامترسازی مثلثاتی یک جایگزین ظریف ارائه می‌دهد که ریشه در هویت ساختاری بین قیود ترکیب و هویت فیثاغورث دارد:

\begin{equation}
\sin^{2}\theta+\cos^{2}\theta=1
\end{equation}

نگاشت ترکیب‌ها به مختصات زاویه‌ای:

\begin{equation}
y_{1}=\sin^{2}\theta,\qquad y_{2}=\cos^{2}\theta
\end{equation}

سیمپلکس ۱ بعدی (مشروط به y1,y2≥0 و y1+y2=1) را بر روی یک دایره واحد پارامتریشده توسط θ∈[0,2π] تعبیه می‌کند. تعمیم به ابعاد بالاتر از مختصات کروی یا ترکیبی از آن‌ها استفاده می‌کند و منیفولدهای همواری بدون لبه یا رأس ایجاد می‌کند که در آن روش‌های عددی مشکل دارند.

ما این فضای تبدیل‌شده را «آب‌های آزاد» می‌نامیم تا بر رهایی از قیود شکننده ترازشده با محور که رفتارهای دکارتی سنتی را زندانی می‌کنند، تأکید کنیم. این بازفرمول‌بندی هندسی چهار مزیت متمایز ارائه می‌دهد. توپولوژی ذاتی هندسه دایره/کره به طور طبیعی نرمال‌سازی را بدون قیود صریح تحمیل می‌کند – هویت فیثاغورثی رفتار جمع به یک را به طور خودکار تضمین می‌کند و نیاز به ضرب‌کننده‌های لاگرانژ یا توابع جریمه در بهینه‌سازی را حذف می‌کند. تناوب هموار تضمین می‌کند که مختصات زاویه‌ای در همه جا مشتق‌پذیر باقی می‌مانند بدون هیچ برخورد خاصی برای موارد مرزی، زیرا θ=0 و θ=2π یک حالت فیزیکی یکسان را بدون ناپیوستگی نشان می‌دهند. تحول نویز بهبود می‌یابد زیرا مراکز ثقل، طول ضلع‌ها و مساحت‌های مثلث محاسبه‌شده از تصویرهای زاویه‌ای، ناخوش‌تعینی عددی بهتری نسبت به محاسبات انجام‌شده مستقیم بر روی مختصات دکارتی مقید نشان می‌دهند. تفسیرپذیری فیزیکی باقی می‌ماند زیرا فواصل هندسی، ناهمگونی ترکیبی را به شیوه‌ای با معنی متریکی تقریب می‌زنند – جدایی زاویه‌ای روی دایره واحد با تفاوت ترکیبی واقعی به گونه‌ای همبستگی دارد که نسبت به چرخش مختصات ناوردا است.

این «آب‌های آزاد» – عدم وجود قیود شکننده ترازشده با محور – جعبه ابزار بزرگتری از هویت‌ها و ناورداهای هندسی که به اصول اول احترام می‌گذارند، ارائه می‌دهد. این مسئله یک مسئله بهینه‌سازی مقید به یک منیفولد با ابعاد پایین‌تر را به یک مسئله بدون قید روی یک دامنه تناوبی تبدیل می‌کند و پایه‌ای برای استخراج ویژگی پایدار که به پارامترسازی انتگرال بیضوی بعدی تغذیه می‌کند، ایجاد می‌کند.

\subsubsection{تصویر مثلثاتی: از متغیرهای متعارف به ویژگی‌های هندسی}

خروجی چهاربعدی CCA:

\begin{equation}
EPS=[E_{1},E_{2},S_{1},P_{1}]
\end{equation}

از طریق تصویرهای مثلثاتی که ساختار رابطه‌ای ضروری را ثبت می‌کنند، به دو ویژگی هندسی کاهش می‌یابد. اولین مثلث، تشکیل‌شده از رئوس در مختصات (E1,E2,S1)، رابطه ساختاری-عنصری را با مرکز ثقل زیر کد می‌کند:

\begin{equation}
Gs=(3E_{1}+E_{2}+S_{1},3yE_{1}+yE_{2}+yS_{1})
\end{equation}

این شیء هندسی نشان می‌دهد که چگونه معماری مولکولی زیست‌توده – بلورینگی سلولز که آرایش کربن را دیکته می‌کند، محتوای آروماتیک لیگنین که توزیع اکسیژن را تحت تأثیر قرار می‌دهد – در ترکیب عنصری کلی ظاهر می‌شود. دومین مثلث، تشکیل‌شده از رئوس در (E1,E2,P1)، رابطه تقریبی-عنصری را با مرکز ثقل زیر کد می‌کند:

\begin{equation}
Gp=(3E_{1}+E_{2}+P_{1},3yE_{1}+yE_{2}+yP_{1})
\end{equation}

این تصویر ثبت می‌کند که چگونه محصولات تجزیه گرمایی – مواد فرار آزادشده در طول گرمایش سریع، کربن ثابت باقی‌مانده به عنوان زغال – از طریق انتخاب مسیر پیرولیز به استوکیومتری عنصری مربوط می‌شوند.

از این ساختارهای مثلثی، دو ویژگی هندسی با تفسیر فیزیکی واضح ظهور می‌یابند. فاصله مرکز ثقل:

\begin{equation}
d=(Gpx-Gsx)2+(Gpy-Gsy)2
\end{equation}

ارزیابی می‌کند که چگونه اجزای ساختاری و تقریبی نسبت به ترکیب عنصری موقعیت‌یابی می‌کنند و معیاری از ناهمگونی ترکیبی چندمقیاسی ارائه می‌دهند. زیست‌توده با مراکز ثقل هم‌تراز (d≈0) سازگاری بین روش‌های مشخصه‌یابی نشان می‌دهد – ساختار مولکولی رفتار گرمایی را پیش‌بینی می‌کند که به نوبه خود با ترکیب عنصری هم‌تراز می‌شود. d بزرگ نشان‌دهنده ویژگی‌های ساختاری است که تنها از ترکیب عنصری کلی قابل پیش‌بینی نیستند، حالتی که در سوخت‌های به شدت پردازش‌شده یا مخلوط‌شده رایج است، جایی که تورفیکاسیون یا اختلاط، همبستگی‌های ترکیبی طبیعی موجود در زیست‌توده دست‌نخورده را مختل می‌کند.

مساحت مثلث K نیازمند ساخت یک نقطه کمکی است. نقطه میانی پاره اتصال‌دهنده مراکز ثقل را تعریف کنید:

\begin{equation}
M=(2Gpx+Gsx,2Gpy+Gsy)
\end{equation}

مساحت مثلث تشکیل‌شده از قاعده E1E2 و رأس M از فرمول دترمینان استاندارد پیروی می‌کند:

\begin{equation}
K=\tfrac{1}{2}|E_{1}(yE_{2}-yM)+E_{2}(yM-yE_{1})+Mx(yE_{1}-yE_{2})|
\end{equation}

مساحت K «گستره» تغییرپذیری زیست‌توده در گروه‌های خصوصیت را ثبت می‌کند. K بزرگتر نشان‌دهنده ناهمگونی بیشتر بین گروهی است و سوخت‌هایی را نشان می‌دهد که رفتار آن‌ها ممکن است به طور قابل توجهی از همبستگی‌های ساده منحرف شود – دقیقاً مواردی که مشخصه‌یابی جامع در ابعاد عنصری، تقریبی و ساختاری برای پیش‌بینی دقیق ارزشمندتر است.

\subsubsection{نمایش ناوردا برای مدل‌سازی مستحکم}

این دو پارامتر:

\begin{equation}
(d,K)\in\mathbb{R}_{\ge 0}^{2}
\end{equation}

۱۳ خصوصیت اولیه زیست‌توده (۵ عنصری + ۴ تقریبی + ۳ ساختاری + رطوبت) را به یک جفت با معنی هندسی و عددی پایدار تقطیر می‌کنند. به طور بحرانی، (d,K) ناورداهای تحت تبدیل‌های رایج از جمله انتقال (جابجایی تمام مختصات با یک بردار ثابت) و مقیاس یکنواخت (ضرب تمام مختصات در یک عامل ثابت) تشکیل می‌دهند. این ناوردایی ویژگی‌های مستحکمی برای مدل‌سازی پایین‌دستی ارائه می‌دهد که حتی زمانی که واحدهای اندازه‌گیری یا قراردادهای نرمال‌سازی در بین مجموعه‌های داده تغییر می‌کنند، به‌خوبی تعریف شده باقی می‌مانند – ملاحظه‌ای عملی با توجه به اینکه داده‌های مشخصه‌یابی زیست‌توده از آزمایشگاه‌های متنوعی با پروتکل‌های تحلیلی و استانداردهای گزارش‌دهی متفاوت سرچشمه می‌گیرند.

\subsection{کدگذاری سلسله‌مراتبی توالی‌ها از طریق جاسازی زمانی ساختاریافته فیبوناچی}

سیستم‌های فیزیکی پویایی در مقیاس‌های زمانی چندگانه نشان می‌دهند، ویژگی‌ای که به ویژه در فرآیندهای گازی‌سازی برجسته است، جایی که گرمایش ذره در طول ثانیه‌ها رخ می‌دهد، پیرولیز در طول ثانیه‌ها تا دقیقه‌ها گسترش می‌یابد، گازی‌سازی زغال از دقیقه‌ها تا ساعت‌ها ادامه می‌یابد و زمان اقامت کلی راکتور در طول دقیقه‌ها است. ثبت این ساختار چندمقیاسی نیازمند نمایش‌هایی است که روابط سلسله‌مراتبی را بدون فروریختن پیچیدگی زمانی به یک مقیاس زمانی مشخصه واحد کد می‌کنند. رویکردهای سنتی با استفاده از ثابت‌های واپاشی نمایی منفرد یا برازش‌های قانون توانی، زمانی ناکافی هستند که مقیاس‌های ناهمگون به طور غیرخطی برهم‌کنش می‌کنند، همان‌طور که اغلب در سیستم‌های حاکم از قوانین پایستگی رخ می‌دهد، جایی که فرآیندهای سریع و کند از طریق کمیت‌های پایستار مشترک جفت می‌شوند.

توالی‌های فیبوناچی – جایی که هر جمله برابر با مجموع دو جمله پیشین است – به طور طبیعی ساختار سلسله‌مراتبی و خودمتشابهی را تعبیه می‌کنند که به ویژه برای نمایش چندمقیاسی مناسب است. در حالی که در ابتدا به عنوان یک رابطه بازگشتی گسسته فرمول‌بندی شده:

\begin{equation}
f_{n}=f_{n-1}+f_{n-2}
\end{equation}

الگوهای رشد شبه فیبوناچی در سیستم‌های فیزیکی نشان‌دهنده رشد نمایی تعدیل‌شده با اثرات حافظه ظاهر می‌شوند، جایی که حالت‌های قبلی تکامل آینده را از طریق پایستگی انرژی، تکانه یا دیگر کمیت‌ها محدود می‌کنند. ساختار بازگشتی هم تاریخچه فوری (جمله n−1 نشان‌دهنده پویایی اخیر) و هم حافظه عمیق‌تر (جمله n−2 نشان‌دهنده تأثیرات مقیاس زمانی طولانی‌تر) را کد می‌کند و سلسله مراتبی طبیعی ایجاد می‌کند که آینه ساختار زمانی تودرتو فرآیندهای فیزیکی چندمقیاسی است.

برای گازی‌سازی زیست‌توده، ما یک دایره با مساحت واحد (مشابه تابع چگالی احتمال که به واحد انتگرال می‌گیرد) را به چهار بخش هم‌پوشان تقسیم می‌کنیم که نشان‌دهنده محرک‌های فرآیند متمایز است که اهمیت نسبی آن‌ها در رژیم‌های عملیاتی مختلف تغییر می‌کند. قید مساحت واحد:

\begin{equation}
\pi r^{2}=1\Rightarrow r=\frac{1}{\sqrt{\pi}}
\end{equation}

نرمال‌سازی هندسی را ایجاد می‌کند. بخش اول تأثیر نسبت بخار به زیست‌توده (SBR) را کد می‌کند و نقش بخار را هم به عنوان یک واکنش‌گر (در واکنش جابجایی آب-گاز و اصلاح بخار) و هم حامل گرما ثبت می‌کند. بخش دوم اثرات جفت‌شده دما (T) و نسبت هم‌ارزی (ER) را نشان می‌دهد و تشخیص می‌دهد که این پارامترها به طور غیرخطی برهم‌کنش می‌کنند – دما سینتیک واکنش را تعیین می‌کند در حالی که نسبت هم‌ارزی اکسیژن موجود را کنترل می‌کند و ترکیب آن‌ها تعریف می‌کند که آیا سیستم در رژیم کنترل‌شده با سینتیک یا رژیم محدودشده با انتقال عمل می‌کند. بخش سوم ویژگی‌های ترکیبی زیست‌توده را از طریق ویژگی‌های هندسی d و K استخراج‌شده از تحلیل CCA ثبت می‌کند و کد می‌کند که چگونه خصوصیات سوخت، مسیرهای واکنش قابل دسترسی را محدود می‌کنند. بخش چهارم نشان‌دهنده برهم‌کنش‌های تکمیلی و جفت‌شدگی‌های متقابل است که توسط سه بخش اول ثبت نشده‌اند و کامل بودن را حتی زمانی که مکانیسم‌های برهم‌کنش خاص به طور ناقص مشخصه‌یابی شده‌اند، تضمین می‌کند.

مساحت بخش‌ها اهمیت نسبی هر عامل را کد می‌کنند، نرمال‌شده به جمع واحد، در نتیجه قید احتمال را ارضا می‌کنند:

\begin{equation}
\sum_{i=1}^{4} s_{i}=1
\end{equation}

یک متغیر مرکب Nj سپس ساخته می‌شود که شامل نسبت‌های بخش (مانند s1/s2 و s3/s4 که تأثیر نسبی گروه‌های محرک مختلف را کمّی می‌کنند)، شرایط عملیاتی (T، ER، SBR که حالت ترمودینامیکی را فراهم می‌کنند)، ویژگی‌های هندسی (d، K که تغییرپذیری سوخت را کد می‌کنند) و ثابت‌های فیزیکی (ثابت گاز R و ثابت فارادی F) است که ابعاد را بین مقیاس‌های ترمودینامیکی و الکتروشیمیایی پیوند می‌دهند. این متغیر مرکب به عنوان بذر آغازین برای تولید توالی فیبوناچی عمل می‌کند و اطلاعات خاص سیستم را در ساختار زمانی سلسله‌مراتبی تعبیه می‌کند.

با توجه به پارامتر نرمال‌شده Nj مشتق‌شده از این ساخت بخش، ما توالی فیبوناچی را از طریق سه جمله آغازین با دقت انتخاب‌شده مقداردهی اولیه می‌کنیم. جمله اول:

\begin{equation}
f_{1}=N_{j}+2\exp(-N_{j})
\end{equation}

پارامتر خطی سیستم را با یک سهم میراشده نمایی ترکیب می‌کند که از رشد نامحدود جلوگیری می‌کند در حالی که حساسیت به تغییرات Nj را حفظ می‌کند. جمله دوم:

\begin{equation}
f_{2}=\exp(-N_{j})\cdot f_{1}
\end{equation}

جفت‌شدگی ضربی بین عامل میرایی و جمله اول را معرفی می‌کند و عدم تقارنی ایجاد می‌کند که تقارن ضمنی در مقداردهی اولیه استاندارد فیبوناچی را می‌شکند. جمله سوم:

\begin{equation}
f_{3}=f_{1}+f_{2}
\end{equation}

از بازگشت استاندارد فیبوناچی پیروی می‌کند و ساختار سلسله‌مراتبی خودمتشابه را آغاز می‌کند. جملات بعدی از بازگشت متعارف پیروی می‌کنند:

\begin{equation}
f_{n}=f_{n-1}+f_{n-2},\qquad n\ge 4
\end{equation}

و توالی‌هایی را تولید می‌کنند که ویژگی‌های رشد آن‌ها حالت سیستم را در شکلی قابل دسترس برای الگوریتم‌های تشخیص الگو کد می‌کند.

میرایی نمایی در مقداردهی اولیه اهداف چندگانه‌ای فراتر از جلوگیری از رشد نامحدود ارائه می‌دهد. یک مقیاس طول مشخصه 1/Nj را در پویایی توالی معرفی می‌کند و تضمین می‌کند که توالی‌های از نقاط عملیاتی مختلف، ویژگی‌های مقیاس‌گذاری نشان می‌دهند که مقایسه بین موارد را تسهیل می‌کند. یک اتصال هموار بین مقداردهی اولیه با انگیزه فیزیکی (مرتبط با شرایط عملیاتی و خصوصیات سوخت) و ساختار ریاضی انتزاعی بازگشت فیبوناچی ایجاد می‌کند. پایداری عددی را با تضمین این که جملات اولیه توالی حتی زمانی که Nj مقادیر شدید می‌گیرد، محدود باقی می‌مانند، ارائه می‌دهد و از سرریز در محاسبات بعدی جلوگیری می‌کند. این توالی‌ها هم اطلاعات محلی (جملات اخیر کدکننده پاسخ فوری سیستم) و هم اطلاعات سراسری (مجموع‌های تجمعی منعکس‌کننده رفتار انتگرالی) را ثبت می‌کنند، ویژگی‌های مقیاس‌گذاری نشان می‌دهند که تشخیص الگو توسط شبکه‌های عصبی را تسهیل می‌کنند (خودمتشابهی در بخش‌های مختلف توالی) و نمایش‌های بدون بعد ارائه می‌دهند که مقایسه‌های معنادار بین موارد با مقیاس‌های فیزیکی بسیار متفاوت را ممکن می‌سازند. به طور معمول، توالی‌هایی با طول ۱۰ تا ۱۵ برای ثبت پویایی‌های ضروری بدون تحمیل بار محاسباتی اضافی کافی هستند، یافته‌ای که به طور تجربی از طریق مطالعات تغییر سیستماتیک بر روی مجموعه داده آموزشی تعیین شده است.

\subsection{آزمایشگاه‌های محاسباتی شبکه عصبی به عنوان موتورهای استنتاج آماری}

به جای آموزش یک شبکه عصبی بزرگ منفرد برای پیش‌بینی مستقیم نتایج گازی‌سازی – رویکردی که خطر بیش‌برازش به داده آموزشی را دارد و تفسیرپذیری فیزیکی محدودی ارائه می‌دهد – ما از سه شبکه کوچک استفاده می‌کنیم که به عنوان «آزمایشگاه‌های محاسباتی» برای استخراج ساختار آماری از توالی‌های فیبوناچی عمل می‌کنند. هر شبکه از یک تابع فعال‌سازی متفاوت استفاده می‌کند و چشم‌اندازهای مکملی از پویایی‌های سیستم کدشده ایجاد می‌کند.

شبکه خطی:

\begin{equation}
\phi(z)=z
\end{equation}

روابط خطی و اطلاعات گرادیان را ثبت می‌کند و حساسیت به تغییرات نسبی در ویژگی‌های ورودی ارائه می‌دهد. شبکه تانژانت هذلولوی:

\begin{equation}
\phi(z)=\tanh(z)=\frac{e^{z}-e^{-z}}{e^{z}+e^{-z}}
\end{equation}

ناهخطی‌های هموار را ثبت می‌کند و خروجی محدود تولید می‌کند و آن را به ویژه مناسب برای شناسایی اثرات اشباع و انتقال‌های سیگموئیدی رایج در سیستم‌های گرماشیمیایی می‌سازد. شبکه نمایی/سیگموئید:

\begin{equation}
\phi(z)=\frac{1}{1+e^{-z}}
\end{equation}

اشباع نامتقارن و اثرات آستانه‌ای را ثبت می‌کند و نشان می‌دهد که رفتار سیستم کجا انتقال‌های تیز بین رژیم‌های کیفی متفاوت نشان می‌دهد.

هر شبکه توالی‌های فیبوناچی را به عنوان ورودی دریافت می‌کند و پس از آموزش بر روی یک زیرمجموعه با دقت انتخاب‌شده از داده گازی‌سازی نشان‌دهنده شرایط عملیاتی و انواع سوخت متنوع، یک ماتریس 5×5 از وزن‌ها و بایاس‌های داخلی تولید می‌کند. به طور بحرانی، این شبکه‌ها برای پیش‌بینی مستقیم ترکیب یا بازده استفاده نمی‌شوند؛ در عوض، به عنوان ابزارهای استنتاج آماری عمل می‌کنند که ساختار را از توالی‌ها استخراج می‌کنند و آن ساختار را در ماتریس‌های وزن یادگرفته‌شده کد می‌کنند. این استفاده غیرمستقیم از شبکه‌های عصبی – به عنوان استخراج‌کننده‌های ویژگی به جای پیش‌بین‌کننده‌ها – استحکام در برابر بیش‌برازش ارائه می‌دهد و یادگیری انتقالی در مجموعه‌های داده مختلف را بدون نیاز به آموزش کامل مجدد ممکن می‌سازد.

برای هر شبکه آموزش‌دیده و هر مورد داده، اثر (رد) ماتریس را محاسبه می‌کنیم:

\begin{equation}
\operatorname{tr}(M)=\sum_{i=1}^{5} M_{ii}
\end{equation}

جایی که M ماتریس داخلی 5×5 را نشان می‌دهد. اثر – که مجموع مقادیر ویژه است – یک ناوردا تحت تبدیل‌های شباهت تشکیل می‌دهد و رفتار تجمعی سیستم را بدون حساسیت به انتخاب پایه دلخواه ضمنی در نمایش وزن شبکه عصبی کد می‌کند. این ناوردایی برای استنتاج مستحکم ضروری است، زیرا سیستم‌های فیزیکی معادل ممکن است بسته به مقداردهی اولیه یا مسیر بهینه‌سازی، ماتریس‌های وزن متفاوتی نتیجه دهند، اما اثر آن‌ها ثابت باقی می‌ماند. اجرای توالی‌های فیبوناچی از تمام موارد از طریق ۳ شبکه ×2 جهت ×3 پیکربندی =۱۸ اثر برای هر مورد یک گسترش ویژگی قابل توجه که عدم تقارن جهت‌دار را ثبت می‌کند در حالی که از نظر محاسباتی قابل مدیریت باقی می‌ماند.

این ۱۸ اثر در دو بردار نمره ادغام می‌شوند:

\begin{equation}
\mathbf{S}^{AB}\in\mathbb{R}^{9},\qquad \mathbf{S}^{BA}\in\mathbb{R}^{9}
\end{equation}

به ترتیب نمره‌های جهت رو به جلو و نمره‌های جهت معکوس را ثبت می‌کنند. عدم تقارن بین رو به جلو و معکوس از نظر فیزیکی معنادار است و وابستگی مسیر را کد می‌کند که گازی‌سازی را به عنوان یک فرآیند ترمودینامیکی برگشت‌ناپذیر مشخصه‌یابی می‌کند. سیستم‌های شیمیایی نزدیک‌شونده به تعادل از حالت‌های اولیه مختلف، مسیرهای متفاوتی را از فضای ترکیب طی می‌کنند، حتی اگر در نهایت به یک نقطه تعادل یکسان برسند. این عدم تقارن جهت‌دار، پسماند (رفتار وابسته به تاریخچه که حالت فعلی به مسیر طی‌شده برای رسیدن به آن بستگی دارد)، محدودیت‌های نرخ (جایی که نرخ‌های واکنش رو به جلو و معکوس به طور قابل توجهی متفاوت هستند) و اثرات عدم تعادل (جایی که سیستم در طول تکامل گذرا از تعادل محلی آنی منحرف می‌شود) را ثبت می‌کند. با کدگذاری صریح جهت‌داری از طریق پردازش رو به جلو و معکوس، ما عدم تقارن ترمودینامیکی بنیادین را در نمایش ویژگی تعبیه می‌کنیم بدون نیاز به مدل‌سازی صریح پویایی وابسته به زمان.

\subsection{تشکیل ماتریس مشخصه از طریق ضرب‌های خارجی متقارن}

با توجه به بردارهای نمره SAB و SBA تولیدشده توسط آزمایشگاه‌های محاسباتی شبکه عصبی، ابتدا هر کدام را به طول واحد از طریق نرمال‌سازی اقلیدسی استاندارد نرمال می‌کنیم:

\begin{equation}
\widehat{\mathbf{S}}^{AB}=\frac{\mathbf{S}^{AB}}{\lVert\mathbf{S}^{AB}\rVert},\qquad \widehat{\mathbf{S}}^{BA}=\frac{\mathbf{S}^{BA}}{\lVert\mathbf{S}^{BA}\rVert}
\end{equation}

این نرمال‌سازی تضمین می‌کند که ساخت ضرب خارجی بعدی بر ساختار همبستگی تأکید می‌کند به جای اینکه توسط تفاوت‌های بزرگی که ممکن است بازتاب انتخاب‌های مقیاس‌گذاری دلخواه در فرآیند آموزش شبکه عصبی باشند، تسلط یابد.

ماتریس مشخصه سپس از طریق یک ضرب خارجی متقارن تشکیل می‌شود که چشم‌اندازهای رو به جلو و معکوس را به شیوه‌ای متوازن ترکیب می‌کند:

\begin{equation}
\mathbf{S}_{c}=\tfrac{1}{2}\Bigl(\widehat{\mathbf{S}}^{AB}\widehat{\mathbf{S}}^{BA\top}+\widehat{\mathbf{S}}^{BA}\widehat{\mathbf{S}}^{AB\top}\Bigr)
\end{equation}

این عمل همبستگی بین چشم‌اندازهای ورودی-اول و خروجی-اول در تکامل سیستم را ثبت می‌کند. تقارنسازی – میانگین گرفتن ضرب خارجی با ترانهاده آن – تضمین می‌کند که Sc متقارن و نیمه معین مثبت باشد، ویژگی ریاضی که تضمین می‌کند مقادیر ویژه حقیقی و غیرمنفی هستند. این ویژگی برای تفسیر فیزیکی بعدی ضروری است، زیرا مقادیر ویژه منفی نشان‌دهنده رفتار غیرفیزیکی (نقض مثبت بودن انرژی یا غیرمنفی بودن احتمال) ناسازگار با قیود ترمودینامیکی است.

این ماتریس 9×9 سپس از طریق بردارسازی و بازتاشیدن به یک ماتریس 3×3 تغییر شکل می‌یابد که اطلاعات ساختاری ضروری را حفظ می‌کند در حالی که ابعاد را فشرده می‌کند. عمل تغییر شکل دلخواه نیست بلکه از یک الگوی سیستماتیک پیروی می‌کند که درایه‌های مرتبط در ماتریس اصلی را به درایه‌های مرتبط در فرم فشرده نگاشت می‌کند و تضمین می‌کند که همبستگی‌های فیزیکی علیرغم کاهش ابعاد کدگذاری شده باقی می‌مانند. انتخاب ابعاد نهایی 3×3 تعادل بین خواسته‌های رقابتی برقرار می‌کند: ماتریس‌های بزرگتر اطلاعات بیشتری حفظ می‌کنند اما هزینه محاسباتی را افزایش می‌دهند و خطر بیش‌برازش دارند، در حالی که ماتریس‌های کوچکتر ساختار را بیش از حد ساده می‌کنند و همبستگی‌های ضروری را از دست می‌دهند.

به طور قابل توجه، در تمام بیش از ۲۰۰ مورد شامل سوخت‌ها، انواع راکتور و شرایط عملیاتی متنوع، ماتریس مشخصه Sc یک الگوی جهت‌دار سازگار نشان می‌دهد که فراتر از مقادیر پارامتر خاص است. در حالی که درایه‌های عددی ماتریس در بین موارد مختلف است – چنان که باید تا اطلاعات خاص مورد را کد کنند – ساختار علامت (کدام درایه مثبت در مقابل منفی هستند) و جهت بردار ویژه غالب به طور قابل توجهی ناوردا باقی می‌مانند. این جهان‌شمولی نشان می‌دهد که Sc توپولوژی ترمودینامیکی بنیادینی را ثبت می‌کند که مستقل از مقادیر پارامتر خاص است و هندسه منیفولد قید را کد می‌کند که تمام مسیرهای گازی‌سازی توسط قوانین پایستگی به آن محدود شده‌اند. پدیده شبیه به نظم فرومغناطیسی است که اسپین‌های محلی با یک میدان سراسری هم‌تراز می‌شوند: موارد منفرد (اسپین‌های محلی) ممکن است از نظر کمی متفاوت باشند، اما یک جهت‌گیری مشترک (ساختار علامت) دارند که توسط تقارن‌های زیرین (قوانین پایستگی) تعیین شده است.

این الگوی ساختاری جهانی جهانی، شکلی از تنظیم‌کنندگی فیزیکی ارائه می‌دهد که تکنیک‌های تنظیم‌کنندگی آماری رایج در یادگیری ماشین را تکمیل می‌کند. در حالی که تنظیم‌کنندگی استاندارد (جریمه‌های L1/L2، حذف تصادفی، توقف زودهنگام) با محدود کردن پیچیدگی مدل از بیش‌برازش جلوگیری می‌کند، تنظیم‌کنندگی فیزیکی از طریق الگوهای ساختاری سازگار با محدود کردن فضای جواب به پیکربندی‌های سازگار با تقارن‌ها و قوانین پایستگی شناخته‌شده، از پیش‌بینی‌های غیرفیزیکی جلوگیری می‌کند. کشف این الگوی جهانی از پیش انتظار نمی‌رفت بلکه از تحلیل سیستماتیک مدل‌های آموزش‌دیده در مجموعه داده کامل ظهور کرد، نمونه‌ای از اینکه چگونه رویکردهای داده‌محور می‌توانند ساختاری را آشکار کنند که بلافاصله از تحلیل اصول اول مشهود نیست.

\subsection{زمان محاسباتی ویژه به عنوان یک ویژگی پویای غیرتعادلی}

خصوصیات ترمودینامیکی سنتی – دما، فشار، ترکیب – حالت‌های تعادل ایستا را مشخصه‌یابی می‌کنند که سیستم در یک پیکربندی مستقل از زمان که انرژی آزاد را تحت قیود کمینه می‌کند، مستقر شده است. با این حال، سینتیک گازی‌سازی شامل فرآیندهای پویایی است که سیستم به طور پیوسته به سمت تعادل تکامل می‌یابد (اما ممکن است به آن نرسد): واکنش‌پذیری زغال، نرخ تبدیل کربن فاز جامد را کنترل می‌کند، نرخ‌های شکست قطران تعیین می‌کنند که آیا ترکیبات پیچیده هیدروکربنی شکسته می‌شوند یا خیر. برای ثبت این جنبه‌های پویا بدون مدل‌سازی صریح گذرا که نیاز به مجموعه داده‌های وابسته به زمان گسترده‌ای دارد، ما «زمان محاسباتی ویژه» (sct) را معرفی می‌کنیم.

sct از یک ویژگی به ظاهر عجیب شکل‌گیری ماتریس مشخصه Sc استفاده می‌کند – عمل بردارسازی و مرتب‌سازی مرتبه‌ای درایه‌های ماتریس – و زمان پردازش واقعی مورد نیاز برای تکمیل این عملیات در یک پلتفرم استاندارد محاسباتی را اندازه‌گیری می‌کند. علیرغم ماهیت ظاهراً تصادفی آن، این زمان پردازش همبستگی‌های قوی با ویژگی‌های فیزیکی عدم تعادل سیستم نشان می‌دهد. پیاده‌سازی دقیق شامل مرتب‌سازی درایه‌های بردارشده ماتریس با استفاده از الگوریتم مرتب‌سازی سریع (quicksort) و ثبت زمان دیواری سپری‌شده برای تکمیل در میکروثانیه است. برای اطمینان از تکرارپذیری، این عمل ۱۰۰ بار تکرار می‌شود و میانه زمان‌ها گزارش می‌شود، و اثرات پرت ناشی از مدیریت فرآیند پس‌زمینه یا نوسانات زمان‌بندی سیستم عامل را کاهش می‌دهد.

فرضیه زمینه‌ساز این است که ماتریس‌های مشخصه حاصل از سیستم‌های نزدیک به تعادل، ساختار منظم‌تری نشان می‌دهند که باعث می‌شود عملیات مرتب‌سازی سریع‌تر انجام شود (به دلیل الگوهای از پیش مرتب تقریباً)، در حالی که ماتریس‌های مربوط به حالت‌های به شدت غیرتعادلی، توزیع درایه‌های نامنظم‌تری دارند که الگوریتم‌های مرتب‌سازی را به چالش می‌کشند و زمان‌های پردازش طولانی‌تری ایجاد می‌کنند. sct، که در این دیدگاه یک معیار غیرمستقیم از «بینظمی» یا پیچیدگی ساختاری ماتریس است، سپس به عنوان یک نماینده کمّی برای درجه عدم تعادل سیستم عمل می‌کند.

اعتبارسنجی تجربی این فرضیه را تأیید می‌کند: موارد با مقادیر sct بالا به طور سیستماتیک با مشخصه‌های عدم تعادل از جمله بازده قطران افزایش‌یافته (نشان‌دهنده تجزیه ناقص)، واکنش‌پذیری زغال پایین‌تر و پراکندگی بیشتر در اندازه‌گیری‌های تجربی تکرارشده مطابقت دارند. برعکس، موارد با sct پایین به حالت‌های نزدیک به تعادل با توزیع محصول مطابق با پیش‌بینی‌های تعادل ترمودینامیکی و تکرارپذیری تجربی بالا گرایش دارند. بنابراین، sct اطلاعاتی در مورد پویایی سیستم فراتر از متغیرهای حالت ایستا ارائه می‌دهد – پلی بین ترمودینامیک تعادل و سینتیک غیرتعادلی – و همه اینها بدون نیاز به مجموعه داده‌های وابسته به زمان یا مدل‌سازی صریح مکانیسم‌های واکنش.

\subsection{پارامترسازی انتگرال بیضوی کامل از طریق همگرایی میانگین حسابی-هندسی}

انتگرال بیضوی کامل نوع اول K(m)، که برای پارامتر m به صورت زیر تعریف می‌شود:

\begin{equation}
K(m)=\int_{0}^{\pi/2}\frac{d\theta}{\sqrt{1-m\sin^{2}\theta}}
\end{equation}

یک تابع خاص کلاسیک است که در مکانیک سماوی، نظریه الاستیسیته و دیگر زمینه‌ها ظاهر می‌شود. محاسبه عددی مستقیم K(m) از طریق انتگرال‌گیری عددی ممکن است به ازای هر m جدید پرهزینه باشد. خوشبختانه، الگوریتم میانگین حسابی-هندسی (AGM) یک روش با همگرایی درجه دو برای محاسبه K(m) با کارایی بالا ارائه می‌دهد.

الگوریتم AGM دنباله‌هایی an و bn را با مقداردهی اولیه a0=1 و b0=1−m تعریف می‌کند و سپس تکرار می‌کند:

\begin{equation}
a_{n+1}=\frac{a_{n}+b_{n}}{2},\qquad b_{n+1}=\sqrt{a_{n}b_{n}}
\end{equation}

تا زمانی که ∣an−bn∣ به تحمل ماشین برسد. آنگاه انتگرال بیضوی از رابطه زیر پیروی می‌کند:

\begin{equation}
K(m)=\frac{\pi}{2\,\AGM\!\bigl(1,\sqrt{1-m}\bigr)}
\end{equation}

خطا تقریباً به صورت 2−2n واپاشی می‌یابد، به این معنی که هر تکرار تعداد ارقام صحیح را دو برابر می‌کند: تکرار ۱ تقریباً ۲ رقم صحیح، تکرار ۲ عدد ۴، تکرار ۳ عدد ۸، تکرار ۴ عدد ۱۶ (فراتر از نیازهای دقت مضاعف) و تکرار ۵ دقت کامل ماشین را برای حساب ممیز شناور دقت مضاعف IEEE 754 با مانتیس ۵۳ بیتی به دست می‌آورد. این همگرایی درجه دو – مشخصه AGM و مزیت محاسباتی که نسبت به روش‌های جایگزین ارائه می‌دهد – تضمین می‌کند که ارزیابی تابع خاص تنها اندکی بیش از عملیات حسابی ابتدایی هزینه دارد.

K(m) محاسبه‌شده – همراه با مقادیر تابع خاص مرتبط مشتق‌شده از تئوری انتگرال بیضوی – سپس از طریق یک ماژول "ellipap" که عملکرد آن مشابه تابع ellipap در جعبه ابزار سیستم کنترل MATLAB است اما برای قیود خاص قانون پایستگی ما تطبیق داده شده است، به پارامترهای تابع انتقال نگاشت می‌شود. این ماژول یک تابع انتقال تولید می‌کند:

\begin{equation}
H(s)=k\,\frac{\prod_{i=1}^{n_{z}}(s-z_{i})}{\prod_{j=1}^{n_{p}}(s-p_{j})}
\end{equation}

جایی که صفرهای zi (ریشه‌های چندجمله‌ای صورت، نشان‌دهنده فرکانس‌هایی که خروجی سیستم صفر می‌شود)، قطب‌های pj (ریشه‌های چندجمله‌ای مخرج مربوط به تشدیدها و فرکانس‌های طبیعی سیستم) و بهره k (عامل مقیاس‌دهی دامنه که بزرگی پاسخ کلی سیستم را تنظیم می‌کند) همگی به صورت جبری از مقادیر انتگرال بیضوی از طریق تبدیل‌های جبری کدکننده قیود فیزیکی ظهور می‌یابند.

ماژول ellipap سه نیاز بحرانی را تحمیل می‌کند که قابلیت تحقق فیزیکی را تضمین می‌کنند. اول، انفعال می‌طلبد که تمام قطب‌ها Re(pj)<0 را ارضا کنند، آن‌ها را در نیم صفحه چپ صفحه مختلط قرار می‌دهد – شرطی معادل با پایداری سیستم که در آن ورودی‌های محدود، خروجی‌های محدود تولید می‌کنند و انرژی تلف می‌شود به جای اینکه بدون حد رشد کند. دوم، نیازهای پایستگی، توپولوژی قطب-صفر و بزرگی بهره را محدود می‌کنند تا سازگاری با موازنه جرم و انرژی تضمین شود و از پیش‌بینی‌هایی که قوانین فیزیکی بنیادین مانند ایجاد ماده یا انرژی از هیچ را نقض می‌کنند، جلوگیری می‌کنند. سوم، علیت نیاز دارد که درجه چندجمله‌ای مخرج برابر یا بیشتر از صورت باشد، تضمین می‌کند که پاسخ سیستم نمی‌تواند مقدم بر ورودی باشد – قید بنیادین در تمام سیستم‌های فیزیکی که اثرات علل را با حداکثر پاسخ آنی (نشان‌داده‌شده توسط درجات برابر) دنبال می‌کنند اما هرگز ورودی‌های آینده را پیش‌بینی نمی‌کنند.

این تحمیل ساختاری – تعبیه قوانین پایستگی در فرم ریاضی تابع انتقال به جای تحمیل آن‌ها از طریق تصحیح‌های پس‌از وقوع – رویکرد ما را به طور بنیادین از شبکه‌های عصبی آموزش‌دیده با پایستگی به عنوان جملات جریمه نرم در توابع زیان متمایز می‌کند. در رویکرد شبکه عصبی، الگوریتم‌های بهینه‌سازی دقت پیش‌بینی را با نقض‌های پایستگی بر اساس وزن جریمه نسبی معامله می‌کنند، به طور بالقوه منجر به سازش‌هایی می‌شود که در آن پیش‌بینی‌ها ممکن است قوانین فیزیکی را نقض کنند، به ویژه در مناطق با داده اندک. در مقابل، چارچوب ما با تضمین اینکه ساختار ریاضی خود مانع از نقض‌ها می‌شود، از چنین سازش‌هایی جلوگیری می‌کند.

متر مکعب به ازای هر کیلوگرم زیست‌توده خشک و عاری از خاکستر: ξ بالاتر نشان‌دهنده بازده بیشتر تولید گاز است، با مقیاس لگاریتمی که تضمین می‌کند همبستگی تقریباً خطی باقی می‌ماند حتی زمانی که بازده مطلق در بین سوخت‌ها و شرایط عملیاتی مختلف با ضریب سه تا چهار تغییر می‌کند. ناوردای گسسته S با بازده تبدیل کربن بیان‌شده به عنوان درصد کربن ورودی ظاهر شده در محصولات فاز گاز همبستگی دارد: موارد S=1 معمولاً ۱۵-۱۰ واحد درصد تبدیل بالاتری نسبت به موارد S=0 در دما و نسبت هم‌ارزی معادل به دست می‌آورند، که منعکس‌کننده مسیرهای واکنش اضافی در دسترس زمانی است که عدم تقارن قطب-صفر شیمی برگشت‌ناپذیر قوی را مجاز می‌کند. به طور بحرانی، این همبستگی‌ها در سراسر فضای پارامتری تجربی یکنوا باقی می‌مانند – عاری از ماکزیمم‌های محلی، نقاط عطف یا مناطق غیریکنوا. این یکنوایی، یک نارسایی را که در بسیاری از مدل‌های گرماشیمیایی مبتنی بر گرادیان مشکل‌آفرین است، حذف می‌کند: بهینه‌های محلی که در آن الگوریتم‌های بهینه‌سازی در پیکربندی‌های پارامتری زیربهینه به دام می‌افتند زیرا تابع هدف چندین قله یا دره نشان می‌دهد. با روابط یکنوا، هر بهینه‌سازی که از حدس اولیه دلخواه شروع می‌شود، با قابلیت اطمینان به بهینه سراسری همگرا می‌شود و طراحی فرآیند مستحکم را بدون نیاز به مطالعات گسترده مقداردهی اولیه یا استراتژی‌های چندشروع ممکن می‌سازد.

\subsection{یکپارچه‌سازی: از داده خام تا پیش‌بینی‌ها از طریق تبدیل‌های احترام‌گذار به پایستگی}

پنج مؤلفه توصیف‌شده در بالا – یکپارچه‌سازی CCA، تبدیل هندسی، کدگذاری فیبوناچی، آزمایشگاه‌های عصبی و پارامترسازی انتگرال بیضوی – ماژول‌های مستقل نیستند بلکه یک خط لوله یکپارچه تشکیل می‌دهند که در آن هر مرحله مرحله بعدی را تغذیه می‌کند در حالی که به تدریج ابعاد را کاهش می‌دهد و ساختار با معنی فیزیکی استخراج می‌کند. داده‌های تجربی خام و مشخصه‌یابی سوخت، شامل ده‌ها اندازه‌گیری برای هر مورد (ترکیب عنصری، آنالیز تقریبی، اجزای ساختاری، شرایط عملیاتی، هندسه راکتور)، ابتدا تحت CCA قرار می‌گیرند تا نمایش فشرده چهاربعدی EPS=[E1,E2,S1,P1] را نتیجه دهند. تبدیل هندسی سپس ویژگی‌های فاصله و مساحت (d,K) را که ناهمگونی ترکیبی چندمقیاسی و تغییرپذیری بین گروهی را کد می‌کنند، استخراج می‌کند. این ویژگی‌های هندسی با شرایط عملیاتی ترکیب می‌شوند تا تولید توالی فیبوناچی را آغاز کنند و جاسازی‌های زمانی (fn) را تولید کنند که پویایی سلسله‌مراتبی در مقیاس‌های چندگانه را ثبت می‌کنند. آزمایشگاه‌های محاسباتی شبکه عصبی این توالی‌ها را از طریق گذرهای رو به جلو و معکوس با سه تابع فعال‌سازی پردازش می‌کنند و بردارهای نمره SAB و SBA را تولید می‌کنند که عدم تقارن جهت‌دار را کد می‌کنند. تشکیل ماتریس مشخصه از طریق ضرب‌های خارجی Sc را نتیجه می‌دهد، که زمان محاسباتی ویژه sct آن صلبیت محاسباتی را به عنوان نماینده عدم تعادل کمّی می‌کند. در نهایت، ارزیابی انتگرال بیضوی – که توسط sct، مجموع تجمعی فیبوناچی و دما آغاز می‌شود – پارامترهای تابع انتقال (k,(zi),(pj)) را از طریق تکرار AGM و تبدیل ellipap تولید می‌کند. از این پارامترها، ناورداهای جهانی (ξ,S) ظهور می‌یابند که به طور یکنوا با کمیت‌های قابل مشاهده تجربی همبستگی دارند: بازده، بازده تبدیل کربن و ترکیب گاز.

کاهش ابعاد حاصل شده از طریق این خط لوله قابل توجه است: از ابعاد اولیه حدود ۳۰ متغیر (خصوصیات سوخت، شرایط عملیاتی، هندسه راکتور) به ۴ متغیر CCA، سپس به ۲ ویژگی هندسی، سپس به ساختار ماتریس مشخصه کدشده در ۹ عنصر، و در نهایت همگرا به ۲ ناوردا (ξ,S) که رفتار ضروری سیستم را ثبت می‌کنند. هر مرحله افزونگی را حذف می‌کند در حالی که محتوای اطلاعاتی مرتبط با پیش‌بینی را حفظ می‌کند، با اثر تجمعی که پیش‌بینی‌های نهایی به یک پارامترسازی فشرده وابسته هستند که چندین مرتبه بزرگی کوچک‌تر از فضای حالت اصلی است اما برای دستیابی به عملکرد کمی قوی کافی است.

تحمیل پایستگی به طور ساختاری در مرحله انتگرال بیضوی رخ می‌دهد نه به صورت الگوریتمی پس‌از وقوع. رویکردهای سنتی – چه حل‌کننده‌های عددی PDE و چه مدل‌های یادگیری ماشین – اغلب پایستگی را از طریق تصحیح‌های اعمال‌شده پس از پیش‌بینی‌های اولیه تحمیل می‌کنند: اگر موازنه جرم نقض شود، ترکیب‌ها نرمال‌سازی مجدد می‌شوند؛ اگر موازنه انرژی شکست بخورد، دماها تنظیم می‌شوند. رویکرد ما چنین تصحیح‌هایی را با تضمین اینکه ساختار ریاضی خود از نقض‌های پایستگی جلوگیری می‌کند، غیرضروری می‌سازد. فرمول‌بندی تابع انتقال:

\begin{equation}
H(s)=k\prod(s-pj)\prod(s-zi)
\end{equation}

با قطب‌های محدود به نیم صفحه چپ، علیت ارضا شده از طریق ساختار گویا مناسب، و بهره محدود شده توسط نگاشت انتگرال بیضوی نمی‌تواند پیش‌بینی‌هایی تولید کند که موازنه جرم یا انرژی را نقض کنند – چنین پیش‌بینی‌هایی خارج از مجموعه ممکن توابع انتقال تولیدپذیر توسط ماژول ellipap قرار دارند. این تضمین به ویژه هنگام برون‌یابی فراتر از داده آموزشی ارزشمند است، جایی که تصحیح‌های الگوریتمی ممکن است شکست بخورند (ترکیب‌های منفی یا دماهای موهومی تولید کنند) اما قیود ساختاری همچنان نقض‌ناپذیر باقی می‌مانند.

کارایی محاسباتی حاصل شده از طریق این چارچوب شایسته تأکید است. یک پیش‌بینی کامل برای یک مورد جدید – با توجه به خصوصیات سوخت و شرایط عملیاتی – نیازمند حدود ۵۰ میلی‌ثانیه زمان دیواری روی یک ایستگاه کاری استاندارد (پردازنده Intel i7، حافظه ۱۶ گیگابایت) اجرای کد MATLAB بدون بهینه‌سازی یا موازی‌سازی خاص است. این زمان‌بندی شامل تمام عملیات است: محاسبه متغیرهای CCA، استخراج ویژگی هندسی، تولید فیبوناچی، گذرهای رو به جلو شبکه عصبی (مشارکت‌کننده غالب در زمان سپری‌شده)، تشکیل ماتریس مشخصه، اندازه‌گیری sct، تکرار AGM (نیازمند حدود ۵ تکرار همان‌طور که توسط تحلیل همگرایی درجه دو پیش‌بینی شده) و استخراج ناوردای نهایی. در مقایسه، یک شبیه‌سازی CFD با وفاداری بالا برای یک مورد معادل نیازمند ۱۲-۴ ساعت روی یک خوشه محاسباتی با کارایی بالا با ۱۲۸-۶۴ هسته است، که نشان‌دهنده شتابی حدود 106× در زمان دیواری سری یا 105× هنگام احتساب موازی‌سازی CFD است. این مزیت محاسباتی چشم‌گیر، کاربردهایی را که قبلاً با CFD غیرممکن بودند، ممکن می‌سازد: بهینه‌سازی فرآیند بلادرنگ که شرایط عملیاتی را هر چند ثانیه بر اساس بازخورد حسگر تنظیم می‌کند، کمّی‌سازی عدم قطعیت جامع از طریق نمونه‌گیری مونت‌کارلو نیازمند هزاران ارزیابی مدل، غربالگری توان‌پذیری بالا ترکیب‌های خوراک جدید در صدها ترکیب ترکیب، و برآورد پارامتر از طریق استنتاج بیزی نیازمند میلیون‌ها ارزیابی درست‌نمایی.\textbf(۳ کاربرد در گازی‌سازی زیست‌توده: مجموعه داده، پیاده‌سازی و چارچوب پایستگی)

\subsection{گردآوری مجموعه داده و پروتکل‌های تضمین کیفیت}

ما بیش از ۲۰۰ مجموعه داده تجربی از ۱۵ منبع مستقل ادبیات در طول دو دهه تحقیق گازی‌سازی گردآوری کردیم، داده‌ها را به طور سیستماتیک استخراج کردیم، سازگاری در بین اندازه‌گیری‌های گزارش‌شده را تأیید کردیم و نتایج را در چهار دسته اولیه سازماندهی کردیم که با هم مشخصه‌یابی جامع سیستم را ارائه می‌دهند. خصوصیات زیست‌توده شامل ترکیب عنصری (کسرهای جرمی کربن، هیدروژن، اکسیژن، نیتروژن، گوگرد)، آنالیز تقریبی (مواد فرار، کربن ثابت، محتوای خاکستر، درصد رطوبت) و اجزای ساختاری (کسرهای جرمی سلولز، همی‌سلولز، لیگنین) برای ۳۸ نوع سوخت متمایز از پسماندهای کشاورزی گرفته تا گونه‌های چوبی و گیاهان انرژی اختصاصی هستند. شرایط عملیاتی دما از ۶۰۰ تا 900∘C، نسبت هم‌ارزی از ۰.۲ تا ۰.۵ (اکسایش زیراستوکیومتری مشخصه گازی‌سازی به جای احتراق)، نسبت بخار به زیست‌توده از ۰ (بدون افزودن بخار) تا ۱.۵ (عملیات غنی از بخار) و ترکیب عامل گازی‌سازی شامل هوا، اکسیژن، بخار و مخلوط‌های مختلف را در بر می‌گیرد. جزئیات پیکربندی راکتور، طراحی بستر سیال حبابی (BFB) در مقابل بستر سیال در گردش (CFB)، نوع مواد بستر و توزیع اندازه ذرات و سرعت گاز سطحی تعیین‌کننده رژیم سیال‌سازی را مشخص می‌کند. مشاهدات تجربی ترکیب گاز محصول (H2، CO، CO2، CH4، N2 گزارش شده بر مبنای خشک و عاری از نیتروژن)، بازده کل گاز اندازه‌گیری‌شده در متر مکعب نرمال به ازای هر کیلوگرم زیست‌توده خشک و عاری از خاکستر و بازده تبدیل کربن کمّی‌کننده کسر کربن ورودی ظاهر شده در محصولات فاز گاز را ثبت می‌کنند.

ادبیات گازی‌سازی تغییرپذیری قابل توجهی در استانداردهای گزارش‌دهی نشان می‌دهد – برخی مطالبات مشخصه‌یابی جامع ارائه می‌دهند در حالی که برخی دیگر تنها زیرمجموعه‌ای از اندازه‌گیری‌ها را گزارش می‌کنند، واحدها ممکن است متفاوت باشند (درصد جرم در مقابل درصد مول، مبنای خشک در مقابل حالت دریافتی) و پروتکل‌های تجربی در بین آزمایشگاه‌ها متفاوت هستند (هندسه‌های راکتور مختلف، روش‌های نمونه‌برداری یا روش‌های تحلیل متفاوت). ما پروتکل‌های دقیق تضمین کیفیت را برای اطمینان از سازگاری و قابلیت اطمینان مجموعه داده پیاده‌سازی کردیم. بسته بودن موازنه جرم با بررسی اینکه موازنه‌های عنصری کربن، هیدروژن و اکسیژن درون عدم قطعیت تجربی بین ورودی (ترکیب سوخت و عامل گازی‌سازی) و خروجی (ترکیب گاز محصول، بازده و باقیمانده زغال/قطران) بسته می‌شوند، تأیید شد. موازنه نیتروژن در مواردی که هم ترکیب گاز و هم نیتروژن ورودی به طور مستقل گزارش شده بودند، اعتبارسنجی متقابل شد و ناسازگاری‌هایی را که ممکن است نشان‌دهنده خطاهای اندازه‌گیری یا واکنش‌های تثبیت/آزادسازی نیتروژن در نظر گرفته‌نشده در تحلیل استاندارد باشند، پرچم‌گذاری کرد. بازده تبدیل کربن زمانی که تنها داده جزئی در دسترس بود (مثلاً زمانی که مطالعات بازده و ترکیب را گزارش کردند اما CCE را به صراحت گزارش نکردند) با فرض بسته بودن ۱۰۰٪ کربن و محاسبه CCE از کربن فاز گاز محاسبه شد؛ مواردی که مقادیر غیرمنطقی نتیجه دادند (CCE کمتر از ۷۰٪ یا بیش از ۱۰۰٪) برای حذف پرچم‌گذاری شدند زیرا احتمالاً نشان‌دهنده ناسازگاری‌های اندازه‌گیری یا مشخصه‌یابی ناقص محصول هستند. تشخیص داده‌های پرت به طور سیستماتیک مواردی را شناسایی کرد با ناسازگاری‌های اندازه‌گیری – مانند بازده گزارش‌شده ناسازگار با ترکیب و استوکیومتری زمانی که موازنه‌های عنصری تحمیل می‌شوند – و این موارد را از توسعه مدل حذف کرد.

پس از فیلتر کیفیت که موارد با داده ناکافی، اندازه‌گیری‌های ناسازگار یا ناهنجاری‌های تجربی آشکار را حذف کرد، ۱۹۸ مجموعه داده با اطمینان بالا برای توسعه و اعتبارسنجی مدل باقی ماند. این مجموعه داده تنظیم‌شده فضای پارامتری قابل توجهی را پوشش می‌دهد: دما از ۶۵۰ تا 880∘C (پس از حذف چند مورد زیر 650∘C که پیرولیز ناقص نشان می‌دادند)، نسبت هم‌ارزی از ۰.۲۲ تا ۰.۴۸، نسبت بخار به زیست‌توده از ۰ تا ۱.۳، محتوای کربن زیست‌توده از ۴۲ تا ۵۱ درصد جرمی، محتوای اکسیژن زیست‌توده از ۳۸ تا ۴۸ درصد جرمی و غلظت هیدروژن گاز محصول حاصل از ۸ تا ۲۴ درصد حجمی (مبنای خشک و عاری از نیتروژن). این تنوع تضمین می‌کند که مدل‌های آموزش‌دیده بر روی مجموعه داده باید روابط فیزیکی واقعی را ثبت کنند به جای حفظ کردن همبستگی‌های پارامتری باریک، و اطمینان می‌دهد که عملکرد اعتبارسنجی‌شده، بازتاب توان پیش‌بینی واقعی است به جای بیش‌برازش به رژیم‌های عملیاتی محدود.

\subsection{استراتژی تقسیم آموزش-آزمایش برای تضمین اعتبارسنجی مستحکم}

برای اطمینان از اینکه معیارهای عملکرد گزارش‌شده بازتاب توان تعمیم واقعی هستند به جای حفظ کردن داده آموزشی، ما از یک استراتژی اعتبارسنجی محافظه‌کارانه سه‌لایه استفاده کردیم که داده‌های استفاده‌شده برای مراحل مختلف توسعه مدل را جدا می‌کند. مجموعه آموزشی شامل ۱۶۸ مجموعه داده است که برای تعیین ضرایب CCA (ترکیبات خطی تعریف‌کننده متغیرهای متعارف از گروه‌های خصوصیت سوخت)، آموزش آزمایشگاه‌های محاسباتی شبکه عصبی (یادگیری ماتریس‌های وزن 5×5 که به عنوان استخراج‌کننده‌های ویژگی عمل می‌کنند) و انجام رگرسیون غیرخطی برای تعیین نگاشت از (sct، ∑fn، T) به پارامتر انتگرال بیضوی m استفاده شدند. مجموعه اعتبارسنجی شامل ۳۰ مجموعه داده است – که به طور کامل در طول تمام عملیات آموزش نگهداشته شدند – و تنها برای تنظیم ابرپارامترها (مانند معماری شبکه عصبی، توابع فعال‌سازی، الگوریتم‌های آموزش) و انتخاب روش (مقایسه تعاریف جایگزین ویژگی هندسی یا طرح‌های مقداردهی اولیه فیبوناچی) استفاده شدند. مجموعه آزمایش شامل مجموعه‌های داده مستقل گزارش‌شده در مطالعات تطبیقی توسط نویسندگان دیگر (به ویژه مشتقی و همکاران) است که هرگز در هیچ نقطه‌ای از توسعه مدل استفاده نشدند و سخت‌ترین ارزیابی تعمیم به شرایط واقعاً دیده‌نشده و خوراک‌های جدید را ارائه می‌دهند.

علاوه بر این، برای پیش‌بینی‌های بازده و بازده تبدیل کربن – که به مجموعه‌های داده تا حدی متمایز از آن‌هایی که برای مدل‌سازی ترکیب استفاده شدند تکیه می‌کنند زیرا تمام مطالعات تمام کمیت‌های قابل مشاهده را گزارش نمی‌کنند – ما از مجموعه‌های اعتبارسنجی ثانویه استفاده می‌کنیم که استقلال را تضمین می‌کنند. برای پیش‌بینی بازده، ۱۱۴ مورد تجربی با مقادیر بازده گزارش‌شده قبلاً برای توسعه مدل ترکیب استفاده نشده بودند و ارزیابی مستقل ارائه دادند. برای پیش‌بینی CCE، ۹۶ مورد تجربی با مقادیر CCE گزارش‌شده به طور مشابه مستقل از داده آموزش ترکیب بودند. این استراتژی اعتبارسنجی چندلایه – جدا کردن داده برای آموزش، انتخاب ابرپارامتر و آزمایش نهایی – از بیش‌برازش از طریق چندین مکانیسم محافظت می‌کند: جلوگیری از نشت اطلاعات که در آن داده اعتبارسنجی از طریق پالایش تکراری مدل به طور ناخواسته بر آموزش تأثیر می‌گذارد، تضمین اینکه انتخاب‌های ابرپارامتر بازتاب عملکرد پیش‌بینی واقعی هستند به جای اتفاقی همراه با تنظیم‌های تصادفی.

\subsection{جزئیات پیاده‌سازی: اجرای CCA، آموزش شبکه عصبی و رگرسیون پارامترها}

پیاده‌سازی CCA از کتابخانه‌های استاندارد MATLAB استفاده می‌کند که ماتریس‌های کوواریانس متقاطع را محاسبه می‌کنند و مسئله مقدار ویژه تعمیم‌یافته را برای به دست آوردن بردارهای متعارف حل می‌کنند. ترکیب‌های خطی به دست آمده E1,E2,S1,P1 برای تمام موارد در مجموعه داده ذخیره می‌شوند و سپس برای محاسبه ویژگی‌های هندسی (d,K) با استفاده از فرمول‌های بخش ۳.۵.۲ استفاده می‌شوند. آموزش شبکه عصبی در پایتون با استفاده از TensorFlow انجام می‌شود، با معماری‌های ساده ۵ نورون در یک لایه پنهان و تابع فعال‌سازی مناسب (خطی، تانژانت هذلولوی یا سیگموئید). هر شبکه با استفاده از الگوریتم بهینه‌سازی Adam با نرخ یادگیری ۰.۰۰۱ برای ۵۰۰ دوره آموزش داده می‌شود، با تابع زیان میانگین مربعات خطا. پس از آموزش، ماتریس وزن 5×5 استخراج می‌شود و اثر آن برای هر مورد محاسبه می‌شود. رگرسیون غیرخطی برای نگاشت از (sct، ∑fn، T) به m با استفاده از تابع `lsqcurvefit` در MATLAB انجام می‌شود که روش‌های کمترین مربعات غیرخطی را پیاده‌سازی می‌کند. پارامترهای بهینه‌شده (α,β,γ,δ,ϵ,ζ) سپس برای پیش‌بینی m برای موارد جدید استفاده می‌شوند.

\subsection{گردش کار محاسباتی: از داده ورودی تا پیش‌بینی‌ها}

برای یک مورد گازی‌سازی جدید مشخص‌شده با خصوصیات سوخت داده‌شده (ترکیب عنصری، آنالیز تقریبی، اجزای ساختاری در صورت موجود بودن) و شرایط عملیاتی (دما، نسبت هم‌ارزی، نسبت بخار به زیست‌توده، نوع راکتور)، پیش‌بینی از طریق گردش کار نه‌مرحله‌ای توسعه‌یافته از طریق روش‌شناسی یکپارچه ارائه‌شده در بخش ۳ پیش می‌رود.

اول، متغیرهای CCA را محاسبه کنید:

\begin{equation}
EPS=[E_{1},E_{2},S_{1},P_{1}]
\end{equation}

با اعمال ضرایب همبستگی متعارف یادگرفته‌شده به اندازه‌گیری‌های خصوصیت سوخت – هنگامی که داده ساختاری در دسترس نیست، S1 از مقادیر عنصری و تقریبی با استفاده از همبستگی‌های ایجادشده در طول توسعه CCA برآورد می‌شود. دوم، ویژگی‌های هندسی (d,K) را از ساخت مثلثی شامل اجزای EPS استخراج کنید، موقعیت مراکز ثقل، فاصله بین مراکز ثقل و مساحت مثلث را طبق فرمول‌های مشخص‌شده در بخش ۳.۲ محاسبه کنید. سوم، توالی فیبوناچی را تولید کنید:

(f1,f2,…,f15)

با استفاده از مقداردهی اولیه مبتنی بر بخش از پارامتر مرکب نرمال‌شده Nj مشتق‌شده از ویژگی‌های هندسی، شرایط عملیاتی و ثابت‌های فیزیکی. چهارم، این توالی را از طریق سه آزمایشگاه محاسباتی شبکه عصبی آموزش‌دیده (توابع فعال‌سازی L، T، E) در هر دو جهت رو به جلو و معکوس تغذیه کنید و ۱۸ مقدار اثر کدکننده ساختار سیستم از چندین چشم‌انداز به دست آورید. پنجم، بردارهای نمره SAB و SBA را با سازماندهی این اثرها تشکیل دهید، سپس ماتریس مشخصه Sc را از طریق ضرب خارجی متقارن و کاهش ابعاد به فرم 3×3 بسازید. ششم، زمان محاسباتی ویژه sct را با اجرای عمل بردارسازی و مرتب‌سازی ماتریس و ثبت زمان سپری‌شده اندازه‌گیری کنید. هفتم، پارامتر انتگرال بیضوی m را از نگاشت غیرخطی یادگرفته‌شده ترکیب‌کننده sct، مجموع تجمعی فیبوناچی ∑fn و دما T ارزیابی کنید. هشتم، انتگرال بیضوی کامل K(m) را از طریق تکرار AGM (معمولاً در ۵ تکرار به دقت ماشین همگرا می‌شود) محاسبه کنید و پارامترهای تابع انتقال (k,(zi),(pj)) را از ماژول ellipap استخراج کنید. نهم، ناورداهای جهانی را محاسبه کنید:

\begin{equation}
\xi=\ln\!\Bigl(\frac{\Lambda_{S}}{\Lambda_{G}}\Bigr),\qquad S=\lvert p\rvert-\lvert z\rvert
\end{equation}

سپس کمیت‌های قابل مشاهده را پیش‌بینی کنید: بازده گاز محصول از ξ از طریق همبستگی تجربی ایجادشده بر داده آموزشی، بازده تبدیل کربن از S از طریق طبقه‌بندی باینری با پالایش رگرسیون و غلظت گونه‌های گازی مجزا از تابع انتقال کامل H(s) از طریق تبدیل معکوس لاپلاس یا تحلیل حوزه فرکانس معادل.

هزینه محاسباتی این گردش کار علیرغم ساختار چندمرحله‌ای آن به طور قابل توجهی کم است. محاسبه متغیر CCA تنها شامل ضرب ماتریس-بردار و عملیات حسابی ساده است.

\subsection{معادلات پایستگی و چارچوب وارسی}

\subsubsection{محدوده و زاویه وارسی}

چارچوب پایستگی، قیود فیزیکی بنیادین بر سیستم را از طریق موازنه جرم و انرژی تحمیل می‌کند. موازنه جرم عنصری (کربن، هیدروژن، اکسیژن، نیتروژن، گوگرد) و موازنه کلی انرژی (قانون اول ترمودینامیک) برای هر مورد اعمال می‌شوند. این موازنه‌ها هم به عنوان معیاری برای سازگاری فیزیکی مدل و هم به عنوان ابزاری برای تصحیح پیش‌بینی‌ها در صورت نیاز عمل می‌کنند.

\subsubsection{موازنه‌های عنصری (جرمی) بر مبنای ۱ کیلوگرم}

چارچوب موازنه جرم با ورودی‌ها، بازده گاز و مشخصات ترکیب مشتق‌شده از مشخصه‌یابی تجربی آغاز می‌شود. آنالیز نهایی زیست‌توده را بر مبنای خشک در نظر بگیرید که کسرهای جرمی (wC,wH,wO,wN,wS,…) را نشان‌دهنده ترکیب عنصری نتیجه می‌دهد. بر مبنای mb=1 کیلوگرم زیست‌توده خشک، مقادیر مولی هر عنصر واردشونده به گازی‌فایر از تقسیم کسرهای جرمی بر جرم اتمی مربوطه پیروی می‌کند:

\begin{equation}
nC,in=MCwCmb,nH,in=MHwHmb,nO,in=MOwOmb\tag{3}
\end{equation}

جایی که MC، MH و MO به ترتیب جرم مولی کربن، هیدروژن و اکسیژن را نشان می‌دهند. جریان‌های ورودی اضافی – اکسیدان و بخار در صورت وجود – از طریق جمله‌هایی مانند nO2,in (اکسیژن از هوا یا تغذیه اکسیژن خالص) و nH2O,in (افزودن بخار) سهم عنصری بیشتری دارند.

مدل بازده گاز محصول Y[Nm3/kg] و کسرهای مولی خشک yi برای گونه i∈(H2,CO,CO2,CH4,…) در شرایط نرمال را پیش‌بینی می‌کند. با نشان دادن حجم مولی در شرایط نرمال با VmN[Nm3/kmol] (معمولاً ۲۲.۴۱۴ برای رفتار گاز ایده‌ال در 0∘C و ۱ اتمسفر)، کل مول‌های محصول و مول‌های گونه مجزا از جفت‌شدگی بازده-ترکیب ظهور می‌یابند:

\begin{equation}
ngas,tot=VmNY(kmol),ni=yingas,tot,i\sumyi=1\tag{4}
\end{equation}

این قید نرمال‌سازی تضمین می‌کند که ترکیب‌های پیش‌بینی‌شده یک حسابداری کامل از محصولات فاز گاز تشکیل می‌دهند.

موازنه کربن، کربن ورودی را به کربن فاز گاز و کربن باقیمانده زغال/خاکستر پیوند می‌زند. با نشان دادن νC,i به عنوان تعداد اتم‌های کربن در گونه i (یک برای CO، CO2 و CH4 در گونه‌های گازی اولیه، صفر برای H2 و N2)، موازنه کربن نیاز دارد:

\begin{equation}
nC,in=i\sumνC,ini+nC,char/ash\tag{5}
\end{equation}

جایی که جمله کربن زغال/خاکستر، کربن جامد تبدیل‌نشده خارج‌شونده از راکتور را حساب می‌کند. بازده تبدیل کربن پیش‌بینی‌شده مستقیماً از این موازنه به صورت زیر به دست می‌آید:

\begin{equation}
CCE%=100nC,in\sumiνC,ini\tag{6}
\end{equation}

با ساختار، کربن تبدیل‌نشده nC,char/ash=nC,in(1−CCE/100) را ارضا می‌کند و حسابداری کامل کربن در فازهای گاز و جامد را تضمین می‌کند.

\subsubsection{موازنه کلی انرژی (قانون اول)}

چارچوب موازنه انرژی از فرمالیسم ترمودینامیکی استاندارد با حسابداری صریح اثرات گرمای محسوس استفاده می‌کند. تعریف Δhˉf∘(Tref) به عنوان آنتالپی مولی استاندارد تشکیل در دمای مرجع (معمولاً 25∘C یا ۲۹۸.۱۵ کلوین) و cˉp,j(T) به عنوان ظرفیت گرمایی مولی (احتمالاً وابسته به دما)، آنتالپی ویژه گونه j در دما T از رابطه زیر پیروی می‌کند:

\begin{equation}
hˉj(T)=Δhˉf,j∘(Tref)+∫TrefTcˉp,j(T′)dT′\tag{9}
\end{equation}

جایی که انتگرال نشان‌دهنده تجمع گرمای محسوس در طول گرمایش از شرایط مرجع به دمای عملیاتی است. برای سیستم‌هایی که ظرفیت گرمایی وابستگی ضعیفی به دما نشان می‌دهد، تقریب‌های cˉp,j ثابت انتگرال را به cˉp,j(T−Tref) ساده می‌کنند، اگرچه نمایش‌های چندجمله‌ای دقیق‌تر می‌توانند زمانی که داده پیچیدگی اضافی را توجیه می‌کند، به کار روند.

برای یک حجم کنترل تک، به‌خوبی مخلوط‌شده که در حالت پایا کار می‌کند با تغییرات انرژی جنبشی و پتانسیل ناچیز و بدون استخراج یا ورود کار شفت، قانون اول ترمودینامیک به یک موازنه آنتالپی بین جریان‌های ورودی و خروجی به اضافه انتقال حرارت خالص کاهش می‌یابد:

\begin{equation}
r\inreactants\sumnrhˉr(Tr)=p\inproducts\sumnphˉp(Tp)+Q\tag{10}
\end{equation}

جایی که Q نشان‌دهنده گرمای خالص منتقل‌شده به محیط اطراف (مثبت زمانی که راکتور گرما از دست می‌دهد، پیرو قرارداد علامت متداولی که گرمای خارج‌شونده از سیستم به عنوان Q مثبت حساب می‌شود) است. این فرمول‌بندی به طور ضمنی فرض می‌کند که تمام آنتالپی‌های واکنش شیمیایی در جمله‌های آنتالپی تشکیل ثبت شده‌اند و تغییرات گرمای محسوس به طور کامل اثرات دما را حساب می‌کنند، هر دو فرض استاندارد در موازنه‌های انرژی گرماشیمیایی.

استراتژی جواب دو شاخه‌ای، ساختار درجه آزادی بنیادین ذاتی در این موازنه انرژی را بهره می‌گیرد. شاخه A سناریوهایی را مورد توجه قرار می‌دهد که جمله گرمای خالص Q شناخته شده است یا می‌تواند از طریق ملاحظات کمکی برآورد شود (مانند عملیات بی‌دررو که در آن طبق تعریف Q=0 است، یا مدل‌های تلفات گرمای تجویزشده به فرم Q=UA(Tbed−T∞) که در آن UA نشان‌دهنده رسانایی کلی انتقال حرارت و T∞ دمای محیط را نشان می‌دهد). با Q مشخص و موازنه‌های عنصری که مقادیر محصول np را ثابت می‌کنند، معادله انرژی به یک معادله جبری غیرخطی در Tbed تبدیل می‌شود – دمای بستر از طریق انتگرال گرمای محسوس وارد آنتالپی محصولات hˉp(Tbed) می‌شود و وابستگی دمایی ایجاد می‌کند که باید به صورت تکراری یا از طریق الگوریتم‌های یافتن ریشه حل شود. حل این معادله، دمای بستر پیش‌بینی‌شده سازگار با مدل تلفات گرمای مشخص‌شده و شرایط ورودی را نتیجه می‌دهد.

شاخه B سناریوی مکمل را مورد توجه قرار می‌دهد که در آن دمای بستر Tbed به عنوان یک کمیت شناخته شده تحمیل می‌شود (معمولاً به طور تجربی از طریق ترموکوپل‌های تعبیه‌شده در بستر سیال اندازه‌گیری می‌شود). با تنظیم Tbed به مقدار اندازه‌گیری‌شده یا هدف، تمام آنتالپی محصولات hˉp(Tbed) به طور کامل تعیین می‌شوند و معادله انرژی می‌تواند به صورت جبری برای انتقال حرارت مورد نیاز حل شود:

\begin{equation}
Q=p\sumnphˉp(Tbed)-r\sumnrhˉr(Tr)\tag{11}
\end{equation}

این Q محاسبه‌شده، گرمای خالصی را که باید برای حفظ عملیات در دمای تحمیل‌شده به راکتور تأمین یا از آن حذف شود، آشکار می‌کند – اطلاعاتی اساسی برای طراحی گرمایی راکتور و مطالعات یکپارچه‌سازی گرما.

یک قانون تلفات گرمای اختیاری می‌تواند بسته‌بندی را برای پیاده‌سازی‌های شاخه A نیازمند مشخص‌سازی صریح Q فراهم کند. فرمول‌بندی‌های رایج شامل قانون خنک‌کنندگی نیوتن Q=UA(Tbed−T∞) است که تلفات گرما را به نیروی محرکه دما از طریق یک رسانایی کلی مرتبط می‌کند، یا هم‌بستگی‌های تجربی به فرم Q=ϵmˉbHHVbiomass که در آن ϵ نشان‌دهنده یک ضریب تلفات گرمای کسری کوچک (معمولاً ۰.۰۵-۰.۰۱ برای راکتورهای آزمایشگاهی به‌خوبی عایق) و HHVbiomass نشان‌دهنده ارزش گرمایی بالای خوراک زیست‌توده است. با این حال، چنین قوانینی باید با احتیاط اعمال شوند – ضریب ϵ تغییرپذیری قابل توجهی از راکتور به راکتور بسته به کیفیت عایق، مقیاس و شیوه‌های عملیاتی نشان می‌دهد و کاربرد جهانی را محدود می‌کند. به همین دلیل، شاخه B (تحمیل دمای اندازه‌گیری‌شده) اغلب برای اعتبارسنجی مدل در برابر داده تجربی قابل اطمینان‌تر است، زیرا عدم قطعیت در پارامترهای انتقال حرارت را حذف می‌کند.

\subsubsection{معادلات دیفرانسیل جزئی پایستگی خط پایه (دو فاز گاز/جامد)}

وارسی حجم کنترل صفر بعدی توصیف‌شده در بالا برای بسته‌بندی لازم است اما نشان‌دهنده ساده‌سازی قابل توجهی از پویایی فضایی-زمانی زیرین است. برای کامل بودن، معادلات پایستگی دو فاز خط پایه را ارائه می‌دهیم که گازی‌سازی بستر سیال را در مقیاس پیوسته حکم‌فرمایی می‌کنند، با این شناسایی که چارچوب کاهش‌یافته ما به طور مؤثر یک بسته‌بندی پیچیده از این معادلات را انجام می‌دهد به جای حل صریح آن‌ها. یک سیستم دو فازی با اندیس فاز k∈(g,s) به ترتیب نشان‌دهنده گاز و جامد، کسر حجمی αk (با ارضا αg+αs=1 توسط پایستگی حجم)، چگالی فاز ρk و میدان سرعت فاز uk را در نظر بگیرید. تبادلات جرم، تکانه و حرارت بین فازی به طور صریح از طریق جمله‌های منبع جفت‌کننده دو فاز ظاهر می‌شوند.

قید پیوستگی مخلوط نیاز دارد که کسرهای حجمی در سراسر دامنه جمع واحد شوند:

\begin{equation}
k\sumαk=1,αk\in[0,1]\tag{12}
\end{equation}

این قید جبری واقعیت فیزیکی را بازتاب می‌دهد که گاز و جامدات با هم کل حجم راکتور را بدون فضای خالی اشغال می‌کنند (با تفسیر فاز گاز به عنوان محیط پیوسته و ذرات جامد به عنوان فاز پراکنده). معادلات پیوستگی فاز منفرد، تکامل کسر حجمی فاز و چگالی را حکم‌فرمایی می‌کنند:

\begin{equation}
\partialt\partial(αk\rhok)+\nabla\cdot(αk\rhokuk)=m≠=k\summ˙mk+Sk(r)\tag{13}
\end{equation}

جایی که m˙mk نرخ انتقال جرم بین فازی در واحد حجم از فاز m به فاز k را نشان می‌دهد (ثبت فرآیندهای تبخیرشدگی و گازی‌سازی زغال که کربن جامد را به محصولات گازی تبدیل می‌کنند) و Sk(r) نشان‌دهنده منبع جرم خالص از واکنش‌های همگن رخ‌داده درون فاز k است (مانند شکست قطران در فاز گاز یا تکامل ساختار منفذ در ذرات زغال).

انتقال گونه فاز گاز، تکامل گونه‌های شیمیایی منفرد درون گاز را حکم‌فرمایی می‌کند. برای گونه i با کسر جرمی Yi، شار پخش Ji (معمولاً از طریق قانون فیک یا فرمول‌بندی‌های انتشار چندجزئی پیچیده‌تر مدل‌سازی می‌شود) و جمله منبع همگن ω˙i(g) نشان‌دهنده نرخ واکنش فاز گاز، معادله انتقال به صورت زیر است:

\begin{equation}
\partialt\partial(αg\rhogYi)+\nabla\cdot(αg\rhogugYi)=-\nabla\cdot(αgJi)+ω˙i(g)+m˙sg\toi\tag{14}
\end{equation}

جایی که m˙sg→i آزادسازی بین فازی گونه i از جامدات از طریق پیرولیز یا گازی‌سازی را حساب می‌کند و اتصال بحرانی بین واکنش‌های ناهمگن در سطوح ذره و تکامل ترکیب همگن فاز گاز را فراهم می‌کند.

معادلات انرژی فاز، تکامل گرمایی را در هر دو فاز گاز و جامد حکم‌فرمایی می‌کنند. بیان پایستگی انرژی در فرم آنتالپی با hk نشان‌دهنده آنتالپی ویژه فاز، qk نشان‌دهنده شار گرمای هدایتی (معمولاً قانون فوریه: qk=−kk∇Tk جایی که kk رسانایی گرمایی است) و ags نشان‌دهنده چگالی سطح مشترک در واحد حجم، معادله انرژی فاز گاز می‌شود:

\begin{equation}
\partialt\partial(αg\rhoghg)+\nabla\cdot(αg\rhogughg)=-\nabla\cdot(αgqg)+i\sumω˙i(g)(-ΔHi)+hgsags(Ts-Tg)+Sg(Q)\tag{15}
\end{equation}

جایی که ΔHi نشان‌دهنده گرمای واکنش برای گونه i است (مثبت برای واکنش‌های گرماگیر مصرف‌کننده گرما، منفی برای واکنش‌های گرمازا آزادکننده گرما)، hgs نشان‌دهنده ضریب انتقال حرارت بین فازی است و Sg(Q) منابع یا چاه‌های گرمای خارجی مانند تبادل گرمای دیواره را ثبت می‌کند. جمله hgsags(Ts−Tg) انتقال حرارت همرفتی بین فازهای گاز و جامد را حکم‌فرمایی می‌کند و تعادل گرمایی را در مقیاس‌های زمانی تعیین‌شده توسط عدد بیو مشخصه پاسخ گرمایی ذره هدایت می‌کند. معادله انرژی فاز جامد ساختار مشابهی اتخاذ می‌کند:

\begin{equation}
\partialt\partial(αs\rhoshs)+\nabla\cdot(αs\rhosushs)=-\nabla\cdot(αsqs)+j\sumω˙j(s)(-ΔHj)+hgsags(Tg-Ts)+Ss(Q)\tag{16}
\end{equation}

با جمله انتقال حرارت بین فازی که اکنون با علامت مخالف ظاهر می‌شود (گرمای کسب‌شده توسط جامدات برابر با گرمای از دست‌رفته توسط گاز است) و پایستگی انرژی در سطح مشترک فاز را تضمین می‌کند.

معادلات تکانه – در حالی که به طور دقیق برای وارسی انرژی و جرم ضروری نیستند – توصیف فیزیکی را برای زمینه‌هایی که تغییرات مکانی در میدان‌های سرعت و فشار معنادار هستند، تکمیل می‌کنند. موازنه تکانه دو فاز استاندارد برای فاز k به صورت زیر است:

\begin{equation}
\partialt\partial(αk\rhokuk)+\nabla\cdot(αk\rhokukuk)=-αk\nablap+\nabla\cdot(αkτk)+αk\rhokg+Mk\tag{17}
\end{equation}

جایی که p نشان‌دهنده میدان فشار (اشتراک‌گذاری شده بین فازها تحت فرضیات مدل دو فاز استاندارد)، τk نشان‌دهنده تانسور تنش لزج برای فاز k، g شتاب گرانشی است و Mk تبادل تکانه بین فازی را از طریق درگ و دیگر نیروهای برهم‌کنش ثبت می‌کند (با ارضا Mg=−Ms توسط قانون سوم نیوتن تا پایستگی کلی تکانه تضمین شود).

\subsubsection{برابری ساختاری SR4 (محاسبه شفاف S)}

شاخص ساختاری گسسته S معرفی‌شده در بخش ۳، شفافیت محاسباتی صریحی را برای تضمین تکرارپذیری و تفسیر فیزیکی می‌طلبد. این کمیت مستقیماً از پیکربندی قطب-صفر تابع انتقال H(s)=k∏(s−zi)/∏(s−pj) ظهور می‌یابد بدون نیاز به تفسیر ذهنی یا انتخاب آستانه. با نشان دادن nz به عنوان تعداد صفرها (ریشه‌های صورت) و np به عنوان تعداد قطب‌ها (ریشه‌های مخرج)، برابری ساختاری از رابطه زیر پیروی می‌کند:

\begin{equation}
S=1[nz≠=np]={01if nz=np(equal degrees)if nz≠=np(asymmetric degrees)\tag{18}
\end{equation}

جایی که 1[⋅] نشان‌دهنده تابع نشانگر است که مقدار ۱ را زمانی که آرگومان آن درست است و ۰ در غیر این صورت می‌گیرد. این طبقه‌بندی باینری در سراسر مجموعه داده کامل مستحکم است، با جدول ۲ که موارد نمونه‌ای در هر دو کلاس ساختاری را نشان می‌دهد.

\begin{table}[t]
\centering
\small
\caption{مثال‌هایی از رابطه ساختاری SR4}
\begin{tabularx}{\linewidth}{@{}*{3}{>{\raggedleft\arraybackslash}X}@{}}
\toprule
مورد & (nz, np) & S=1 (nz≠=np) \\
\midrule
A & (۲, ۲) & ۰ (درجات برابر) \\
B & (۳, ۳) & ۰ (درجات برابر) \\
C & (۱, ۲) & ۱ (صورت < مخرج) \\
D & (۲, ۱) & ۱ (صورت > مخرج) \\
\bottomrule
\end{tabularx}
\end{table}


تفسیر فیزیکی این تمایز ساختاری با توپولوژی مسیر واکنش همان‌طور که در بخش ۵ بحث شده است، هم‌تراز است. موارد با S=0 ساختار قطب-صفر متوازن نشان می‌دهند که با عملیات نزدیک به تعادل سازگار است که در آن نرخ‌های واکنش رو به جلو و معکوس تقریباً برابری را حفظ می‌کنند، در حالی که موارد S=1 عدم تقارن مشخصه شیمی برگشت‌ناپذیر قوی نشان می‌دهند که در آن محصولات از طریق مسیرهایی تشکیل می‌شوند که فاقد مشارکت معکوس قابل توجه هستند. این ویژگی توپولوژیکی – که صرفاً از تحلیل درجه تابع انتقال بدون نیاز به مدل‌سازی سینتیکی دقیق ظهور می‌یابد – قدرت پیش‌بینی قابل توجهی برای بازده تبدیل کربن همان‌طور که نتایج اعتبارسنجی نشان می‌دهند، ارائه می‌دهد.

\subsubsection{فهرست وارسی}

چارچوب پایستگی ترسیم‌شده در بالا، وارسی سیستماتیک را ممکن می‌سازد که کمیت‌های پیش‌بینی‌شده به قوانین فیزیکی بنیادین احترام می‌گذارند. عبور موفقیت‌آمیز این بررسی‌های وارسی، اطمینان می‌دهد که مدل کاهش‌یافته علیرغم ساده‌سازی محاسباتی چشمگیر نسبت به CFD کامل، فیزیک ضروری را ثبت می‌کند. وارسی بسته بودن جرم نیاز دارد که Y و (yi) پیش‌بینی‌شده تمام موازنه‌های عنصری را درون عدم قطعیت تجربی ارضا کنند (معمولاً ±2\% برای کربن، ±5\% برای هیدروژن با توجه به محدودیت‌های دقت اندازه‌گیری)، با CCE\% محاسبه‌شده از کربن فاز گاز مطابق با مقادیر گزارش‌شده یا پیش‌بینی‌شده مستقل. بسته بودن انرژی تحت شاخه A می‌طلبد که Tbed حل‌شده از نظر فیزیکی قابل قبول برای نقطه عملیاتی مشخص‌شده باشد – دماها باید درون گستره تحمیل‌شده ۹۰۰–۶۰۰∘C قرار گیرند و روندهای مناسب با نسبت هم‌ارزی (افزایش با ER با تأمین گرمای بیشتر اکسیدان) و نسبت بخار به زیست‌توده (کاهش با SBR با مصرف گرمای واکنش‌های گرماگیر بخار) نشان دهند. بسته بودن انرژی شاخه B نیاز دارد که Q محاسبه‌شده با تلفات دیواری معتبر سازگار باقی بماند، معمولاً 5−1\% ارزش گرمایی سوخت برای راکتورهای آزمایشگاهی به‌خوبی عایق یا کسرهای بزرگتر (تا 15−10\%) برای واحدهای در مقیاس پایلوت با نسبت سطح به حجم بیشتر.

سازگاری با معادلات دیفرانسیل جزئی خط پایه ارائه‌شده در بالا، به طور مفهومی از شناسایی این که هر مدل کاهش‌یافته به طور مؤثر نشان‌دهنده یک بسته‌بندی میانگین‌گیری حجمی، میانگین‌گیری زمانی معادلات فضایی-زمانی کامل است، پیروی می‌کند. چارچوب صفر بعدی از انتگرال‌گیری معادلات دیفرانسیل جزئی دو فاز بر روی حجم راکتور تحت فرضیات اختلاط کامل (خصوصیات یکنواخت فضایی)، شرایط پایا (مشتقات زمانی ناپدید می‌شوند) و گرادیان‌های محوری ناچیز (مناسب برای بسترهای حبابی با اختلاط شدید اما به طور بالقوه مشکل‌زا برای بسترهای در گردش با لایه‌بندی عمودی قابل توجه) ظهور می‌یابد. این ساده‌سازی‌ها – هرچند قابل توجه – برای پیش‌بینی عملکرد و بهینه‌سازی قابل قبول هستند که در آن تفکیک فضایی جزئیات ارزش حاشیه‌ای کمی نسبت به هزینه محاسباتی آن ارائه می‌دهد.

چارچوب برابری ساختاری SR4 محدوده صریحی را می‌طلبد تا از تفسیر بیش از حد نتایج جلوگیری شود. ناورداهای (ξ,S) باید به طور مشترک به کار روند نه در انزوا: ξ رفتار بازده پیوسته را در سراسر فضای پارامتری کامل ثبت می‌کند، در حالی که S طبقه‌بندی گسسته رفتار تبدیل را ارائه می‌دهد با قوی‌ترین همبستگی ظهور یافته در تبدیل به اندازه کافی بالا (CCE > 70\%) جایی که توپولوژی قطب-صفر به وضوح ساختار مسیر واکنش را بازتاب می‌دهد. ادعاها در مورد این ناورداها در سطح سطل (آمار مجموعه‌ای بر روی گروه‌های موارد مشابه) باقی می‌مانند مگر اینکه به طور صریح به سطح نمونه ارتقا یابند با شواهد اضافی از مطالعات اغتشاش یا تحلیل‌های حساسیت نشان‌دهنده اینکه پیش‌بینی‌های مورد منفرد تحت تغییرات پارامتری متوسط پایدار باقی می‌مانند.

\subsubsection{تشخیص‌ها و شاخص مرکب عدم تعادل}

فراتر از ناورداهای جهانی (ξ,S) به کار رفته برای پیش‌بینی‌های اولیه، چارچوب ما کمیت‌های تشخیصی اضافی تولید می‌کند که برای درک شخصیت عدم تعادل سیستم و ارزیابی اطمینان مدل ارزشمند ثابت می‌شوند. انحراف لگاریتمی پیوسته ξ:=ln(ΛS/ΛG) با رفتار بازده همان‌طور که از طریق اعتبارسنجی اثبات شده است، هم‌تراز است، در حالی که شاخص ساختاری گسسته S∈(0,1) در رژیم‌های تبدیل به اندازه کافی بالا با رفتار تبدیل کربن قفل می‌شود. این دو کمیت – استخراج‌شده از جنبه‌های اساساً متفاوت توپولوژی تابع انتقال (بزرگی بهره در مقابل درجه قطب-صفر) – چشم‌اندازهای مکملی از حالت سیستم ارائه می‌دهند، با ملاحظه مشترک آن‌ها که استنتاج مستحکم‌تری نسبت به هر کمیت به تنهایی نتیجه می‌دهد.

یک شاخص مرکب عدم تعادل، چندین سیگنال تشخیصی را در یک اسکالر واحد ادغام می‌کند که انحراف از رفتار تعادل ایده‌ال را کمّی می‌کند. یکی از چنین شاخص‌هایی فرم زیر را اتخاذ می‌کند:

\begin{equation}
\epsilon:=\operatorname{median}(|ξ|)+αVar(sct)+β1[S=1]\tag{19}
\end{equation}

جایی که ضرایب مثبت α و β مشارکت‌های نسبی واریانس زمان محاسباتی ویژه و برابری ساختاری را وزن می‌دهند. میانه انحراف لگاریتمی مطلق، تغییرات سیستماتیک بهره را در یک مجموعه مورد ثبت می‌کند، در حالی که واریانس در sct پراکندگی در صلبیت محاسباتی را آشکار می‌کند که به طور بالقوه نشان‌دهنده سینتیک ناهمگن یا محدودیت‌های انتقال است. سهم برابری ساختاری 1[S=1] یک افزایش ثابت برای موارد نشان‌دهنده توپولوژی قطب-صفر نامتقارن مشخصه شیمی برگشت‌ناپذیر اضافه می‌کند. همان‌طور که ϵ→0 (قابل دستیابی تنها زمانی که انحرافات لگاریتمی ناپدید می‌شوند، sct واریانس ناچیزی نشان می‌دهد و برابری ساختاری نشان‌دهنده مسیرهای متوازن است)، باقیمانده‌های بسته بودن جرمی به طور سیستماتیک کاهش می‌یابند و جواب به حد سیال کامل نزدیک می‌شود که در آن مقاومت‌های سینتیکی ناپدید می‌شوند و تعادل غالب می‌شود – یک حالت محدودکننده مفید برای معیارسنجی اما به ندرت در گازی‌سازی عملی قابل دستیابی است که زمان‌های اقامت محدود و محدودیت‌های انتقال از تعادل کامل جلوگیری می‌کنند.

\subsubsection{تعاریف متغیرها و یکاها}

\begin{table}[t]
\centering
\small
\caption{تعاریف متغیرها و یکاها}
\begin{tabularx}{\linewidth}{@{}*{3}{>{\raggedleft\arraybackslash}X}@{}}
\toprule
نماد & معنی & یکا \\
\midrule
Y & بازده گاز محصول به ازای ۱ کیلوگرم زیست‌توده (شرایط نرمال) & Nm3/kg \\
yi & کسر مولی گونه i در گاز محصول (خشک) & – \\
Vm & حجم مولی در شرایط نرمال (مثلاً ۲۲.۴۱۴) & Nm3/kmol \\
nj & مول‌های گونه j & kmol \\
wC,… & کسرهای جرمی در آنالیز نهایی (خشک) & – \\
MC,… & جرم مولی‌ها & kg/kmol \\
CCE\% & بازده تبدیل کربن & \% \\
Q & گرمای خالص منتقل‌شده به محیط اطراف (تلفات اگر Q>0) & kJ یا kW \\
Tbed & دمای گازی‌فایر (بستر) & K \\
UA & رسانایی گرمایی کلی به محیط اطراف & kW/K \\
αk & کسر حجمی فاز (k∈(g,s)) & – \\
ρk & چگالی فاز & kg/m3 \\
uk & سرعت فاز & m/s \\
hk & آنتالپی ویژه & kJ/kg \\
hgs,ags & ضریب انتقال حرارت بین فازی × چگالی سطح & kW/(m3⋅K) \\
m˙mk & نرخ انتقال جرم بین فازی (از m به k) & kg/(m3⋅s) \\
ω˙i(k) & منبع واکنش گونه i در فاز k & kg/(m3⋅s) \\
\bottomrule
\end{tabularx}
\end{table}


یادداشت‌های کاربرد همراه این جدول برای اعمال صحیح ضروری هستند. در عمل، متخصصان باید شاخه A را زمانی که یک مدل تلفات گرمای معتبر در دسترس است (یا هنگام آزمایش کران بالای بی‌دررو نشان‌دهنده حداکثر دمای قابل دستیابی تحت تلفات گرمای صفر) انتخاب کنند و شاخه B را زمانی که دما از طریق ابزارهای راکتور اندازه‌گیری یا کنترل می‌شود. وضوح، بیان صریح این که کدام شاخه هنگام ادعای پیش‌بین‌کنندگی به کار رفته است را می‌طلبد، زیرا این انتخاب اساساً بر درجات آزادی در مسئله وارسی تأثیر می‌گذارد. برای تحلیل‌های هم‌تراز با SR4 با استفاده از چارچوب برابری ساختاری، گزارش باید شامل ξ، S و اختیاریاً شاخص مرکب ϵ در سطح سطل همراه با آمار بسته بودن جرم و انرژی باشد و ارزیابی شفافی از عملکرد مدل و سازگاری فیزیکی ارائه دهد.

تمایزهای بحرانی از سردرگمی بالقوه در نمادگذاری جلوگیری می‌کنند. اصطلاحات «صفرها» و «قطب‌ها» در این کار به طور انحصاری به ریشه‌های صورت و مخرج تابع انتقال H(s) اشاره دارند – این‌ها اشیاء کنترل-نظری مشخصه‌یابی پویایی سیستم در حوزه لاپلاس هستند و هیچ رابطه‌ای با مفاهیم نظریه اعداد مانند صفرهای تابع زتای ریمان ندارند علیرغم تشابه اصطلاحی سطحی. به طور مشابه، نماد یونانی ξ نشان‌دهنده انحراف لگاریتمی بهره معرفی‌شده به عنوان یک ناوردای جهانی است؛ نباید با تابع زتای ریمان ζ(s) یا «تابع کسی ریمان» مرتبط اشتباه گرفته شود:

\begin{equation}
\xi(s):=\tfrac{1}{2}s(s-1)\pi^{-s/2}\Gamma\!\left(\frac{s}{2}\right)\zeta(s)
\end{equation}

که در نظریه اعداد تحلیلی ظاهر می‌شود، علیرغم هم‌پوشانی نمادگذاری که گاه در ادبیات ریاضی ظاهر می‌شود.

\subsubsection{خلاصه: یکپارچه‌سازی چارچوب پایستگی}

معادلات پایستگی ارائه‌شده در سراسر این زیربخش، قیود فیزیکی بنیادینی را فراهم می‌کنند که چارچوب کاهش‌یافته ما به‌جای الگوریتمی، به‌طور ساختاری به آن‌ها احترام می‌گذارد. فرمول‌بندی تابع انتقال:

\begin{equation}
H(s)=k\prod(s-pj)\prod(s-zi)
\end{equation}

پایستگی جرم را از طریق موازنه‌های عنصری تعبیه‌شده در توپولوژی قطب-صفر (ساختار گسسته درجات تابع انتقال بازتاب‌دهنده قیود استوکیومتری)، پایستگی انرژی را از طریق مقیاس‌دهی بهره مناسب و قیود انفعال (تضمین اتلاف انرژی به جای تولید جعلی) و قابلیت تحقق فیزیکی را از طریق پایداری (قطب‌های نیم صفحه چپ تضمین‌کننده پاسخ محدود به ورودی‌های محدود) و علیت (ساختار تابع گویا مناسب جلوگیری‌کننده از خروجی علی‌نقیض) تحمیل می‌کند. با لنگر انداختن کاهش ابعاد داده‌محور – دربرگیرنده یکپارچه‌سازی چندگروهی CCA، استخراج ویژگی هندسی در فضای «آب‌های آزاد» و کدگذاری سلسله‌مراتبی فیبوناچی – به این قوانین پایستگی از طریق مرحله پارامترسازی انتگرال بیضوی، چارچوب تضمین می‌کند که تمام پیش‌بینی‌ها بدون توجه به دستاوردهای چشمگیر کارایی محاسباتی حاصل شده، از نظر فیزیکی سازگار باقی می‌مانند. این تحمیل ساختاری، رویکرد ما را اساساً از روش‌های یادگیری ماشین خالص که پایستگی را به عنوان قیود نرمی در نظر می‌گیرند که باید از طریق جمله‌های جریمه در توابع زیان به طور تقریبی ارضا شوند، متمایز می‌کند، رویکردی که اجتناب‌ناپذیراً نقض پایستگی را در مناطق با داده اندک یا در طول برون‌یابی فراتر از توزیع‌های آموزشی مجاز می‌کند.\textbf(۴ نتایج و اعتبارسنجی)

پس از ایجاد چارچوب روش‌شناختی و کاربرد آن بر گازی‌سازی زیست‌توده، اکنون نتایج اعتبارسنجی جامعی را ارائه می‌دهیم که عملکرد کمی را در شرایط تجربی متنوع نشان می‌دهد. این بخش از طریق پنج زیربخش که معیارهای عملکرد کمی بر روی داده آزمایش مستقل، ارزیابی تطبیقی در برابر مدل‌های گازی‌سازی موجود، تفسیر فیزیکی ناورداهای جهانی استخراج‌شده، تحلیل رده‌بندی شده در گستره رژیم‌های عملیاتی و انواع سوخت و تحلیل خطای سیستماتیک شناسایی‌کننده محدودیت‌ها و حالت‌های شکست را بررسی می‌کنند، گسترش می‌یابد. در کنار هم، این تمرینات اعتبارسنجی ثابت می‌کنند که چارچوب تابع خاص آگاه از فیزیک به عملکرد پیش‌بینی قوی دست می‌یابد در حالی که بسته بودن پایستگی و کارایی محاسباتی را حفظ می‌کند – ترکیبی که کاربردهای قبلاً غیرممکن با CFD با وفاداری بالا یا هم‌بستگی‌های صرفاً تجربی را ممکن می‌سازد.

\subsection{معیارهای عملکرد کمی}

\subsubsection{ترکیب گاز محصول}

پیش‌بینی‌ها برای چهار گونه گاز اصلی – هیدروژن، مونوکسید کربن، دی‌اکسید کربن و متان – بر روی ۷۰ مجموعه داده آزمایش مستقل که به طور کامل در طول توسعه مدل نگهداشته شده بودند، ارزیابی شدند. این موارد تمام گستره شرایط عملیاتی نمایندگی‌شده در ادبیات گردآوری‌شده (دماهای ۸۸۰–۶۵۰ درجه سلسیوس، نسبت هم‌ارزی ۰.۴۸-۰.۲۲، نسبت بخار به زیست‌توده ۱.۳-۰) و انواع سوخت متنوع (پسماندهای کشاورزی، گونه‌های چوبی، گیاهان انرژی) را در برمی‌گیرند و ارزیابی دقیقی از توان تعمیم ارائه می‌دهند. جدول ۴ معیارهای عملکرد را برای هر گونه به طور مجزا و در مجموع خلاصه می‌کند و همبستگی قوی بین پیش‌بینی‌ها و اندازه‌گیری‌های تجربی در تمام اجزا را آشکار می‌کند.

\begin{table}[t]
\centering
\small
\caption{عملکرد پیش‌بینی ترکیب گاز}
\begin{tabularx}{\linewidth}{@{}*{4}{>{\raggedleft\arraybackslash}X}@{}}
\toprule
گونه & R2 & RMSE (حجم٪) & میانگین خطای مطلق (حجم٪) \\
\midrule
H2 & ۰.۸۴ & ۳.۲ & ۲.۵ \\
CO & ۰.۸۱ & ۴.۱ & ۳.۳ \\
CO2 & ۰.۷۸ & ۲.۸ & ۲.۲ \\
CH4 & ۰.۷۹ & ۱.۵ & ۱.۲ \\
\textbf(کلی) & \textbf(۰.۸۱) & \textbf(۲.۹) & \textbf(۲.۳) \\
\bottomrule
\end{tabularx}
\end{table}


تمام مقادیر ضریب تعیین از ۰.۷۸ فراتر می‌روند، نشان می‌دهد که مدل بیش از سه‌چهارم واریانس در اندازه‌گیری‌های تجربی را ثبت می‌کند – عملکردی قابل مقایسه یا فراتر از تکرارپذیری تجربی معمول برای مطالعات گازی‌سازی که اندازه‌گیری‌های تکرارشده اغلب پراکندگی ۱۰-۵٪ به دلیل ناهمگونی خوراک، پویایی سیال‌سازی و تغییرپذیری نمونه‌برداری نشان می‌دهند. تحلیل باقیمانده (پیش‌بینی منهای اندازه‌گیری) هیچ اریب سیستماتیکی نسبت به دما، نسبت هم‌ارزی یا نوع سوخت آشکار نمی‌کند و تأیید می‌کند که خطاهای مدل به طور تصادفی حول صفر توزیع می‌شوند به جای نشان دادن روندهایی که نشان‌دهنده شکاف‌های فیزیکی سیستماتیک یا فرم‌های تابعی نامناسب هستند.

این عدم اریب، اطمینان می‌دهد که چارچوب روابط فیزیکی واقعی را ثبت می‌کند به جای حفظ کردن آثار داده آموزشی، و استفاده از آن را برای درون‌یابی درون فضای پارامتری اعتبارسنجی‌شده و برون‌یابی محتاطانه به شرایط نزدیک پشتیبانی می‌کند.

\subsubsection{بازده تبدیل کربن}

برای ۹۶ اندازه‌گیری CCE مستقل دربرگیرنده سطوح تبدیل از 72\% تا 98\% (بازتاب دهنده گستره از عملیات دمای پایین محدودشده با سینتیک تا تبدیل دمای بالا نزدیک به کامل)، چارچوب به ضریب تعیین R2=0.87، جذر میانگین مربعات خطا RMSE =5.2\% (مطلق) و میانگین خطای مطلق MAE =4.1\% دست می‌یابد. این معیارها نشان‌دهنده توان پیش‌بینی قوی در سراسر طیف تبدیل هستند، با خطاهایی که معمولاً درون کران‌های عدم قطعیت تجربی قرار می‌گیرند. به طور بحرانی، پیش‌بینی‌های CCE رابطه یکنوایی با ناوردای گسسته S در تمام رژیم‌های عملیاتی بررسی‌شده نشان می‌دهند: S بالاتر (نشان‌دهنده توپولوژی قطب-صفر نامتقارن مشخصه شیمی برگشت‌ناپذیر) به طور سازگار بدون توجه به دما، نسبت هم‌ارزی یا خصوصیات سوخت با مقادیر CCE بالاتر مطابقت دارد. این یکنوایی – عدم وجود ماکزیمم‌های محلی، نقاط عطف یا مناطق غیریکنوایی که بهینه‌سازی را پیچیده می‌کنند – برای کاربردهای طراحی فرآیند ارزشمند ثابت می‌شود. الگوریتم‌های بهینه‌سازی مبتنی بر گرادیان می‌توانند با اطمینان برای شرایط حداکثر تبدیل جستجو کنند بدون خطر به دام افتادن الگوریتمی در ماکزیمم‌های محلی زیربهینه، نارسایی که اغلب مدل‌های گرماشیمیایی را که رفتار غیرخطی پیچیده با چندین اثر رقابتی نشان می‌دهند، مشکل‌دار می‌کند.

\subsubsection{بازده گاز محصول}

برای ۱۱۴ اندازه‌گیری بازده مستقل در گستره 0.8±0.2 Nm3/kg زیست‌توده (مبنای خشک و عاری از خاکستر)، چارچوب به R2=0.82، RMSE =0.48 Nm3/kg و MAE =0.38 Nm3/kg دست می‌یابد. این معیارهای عملکرد نشان می‌دهند که مدل تغییرات بازده را در بین سوخت‌ها و شرایط عملیاتی مختلف ثبت می‌کند در حالی که خطاهای متناسب با دقت اندازه‌گیری تجربی (معمولاً ±0.1−0.2 Nm3/kg برای سیستم‌های اندازه‌گیری گاز به کار رفته در گازی‌فایرهای آزمایشگاهی) حفظ می‌کند. پیش‌بینی‌های بازده به طور یکنوا با ناوردای پیوسته ξ=ln(ΛS/ΛG) در سراسر مجموعه داده کامل همبستگی دارند، نشان می‌دهند که این کمیت انحراف لگاریتمی – علیرغم تعریف انتزاعی آن به عنوان بهره تابع انتقال نرمال‌شده – محرک‌های فیزیکی بنیادین بازده تولید گاز را ثبت می‌کند. طبیعت بدون مقیاس ξ به ویژه ارزشمند است: با نرمال‌سازی بهره‌هایی که در ۱۳ مرتبه بزرگی (10−47 تا 10−60) در نوسان هستند به یک متغیر بدون بعد با مرتبه واحد، چارچوب مقایسه‌های کمی معنادار را در بین رژیم‌های ترمودینامیکی بسیار متفاوت از گازی‌سازی کنترل‌شده با سینتیک دمای پایین تا عملیات نزدیک به تعادل دمای بالا ممکن می‌سازد.

\subsection{ارزیابی تطبیقی در برابر مدل‌های تثبیت‌شده گازی‌سازی}

برای قراردادن عملکرد چارچوب ما در منظره گسترده‌تر رویکردهای مدل‌سازی گازی‌سازی، مقایسه‌های سیستماتیکی در برابر سه مدل برتر ارزیابی‌شده در یک مطالعه جامع مستقل توسط مشتقی و همکاران انجام دادیم. این مدل‌های تثبیت‌شده نشان‌دهنده رویکردهای فلسفی متمایزی به پیش‌بینی گازی‌سازی هستند که هر کدام نقاط قوت و محدودیت‌های مشخصه‌ای دارند که این حوزه را در دهه‌های اخیر شکل داده‌اند. مدل I یک روش‌شناسی شبه‌تعادلی با تنظیم دما پیاده‌سازی می‌کند، که معمولاً به عنوان روش QET (دمای شبه‌تعادلی) شناخته می‌شود، که تعادل ترمودینامیکی را در یک دمای اصلاح‌شده پایین‌تر از دمای واقعی بستر برای حساب کردن محدودیت‌های سینتیکی فرض می‌کند. مدل II از هم‌بستگی‌های تجربی برای غلظت‌های گونه کلیدی استفاده می‌کند و ترکیب گاز محصول را به پارامترهای عملیاتی از طریق روابط چندجمله‌ای یا نمایی برازش‌شده مشتق‌شده از تحلیل رگرسیون مجموعه‌های داده تجربی مرتبط می‌سازد. مدل III معادلات سینتیکی دقیق را با مدل‌های هیدرودینامیکی جفت می‌کند و معادلات پایستگی را برای گونه‌های شیمیایی چندگانه همراه با موازنه‌های تکانه و انرژی برای ثبت تغییرات مکانی و رفتار گذرا درون بستر سیال حل می‌کند.

ارزیابی تطبیقی از مجموعه‌های آزمایش مستقل عمداً حذف‌شده از تمام تلاش‌های توسعه مدل استفاده کرد و تضمین کرد که معیارهای گزارش‌شده بازتاب توان پیش‌بینی واقعی هستند به جای حفظ کردن داده آموزشی. این موارد آزمایشی خوراک‌های متنوعی از پسماندهای کشاورزی (پوسته برنج با محتوای خاکستر مشخصاً بالا، کاه گندم با سطوح خاکستر متوسط، ساقه ذرت با ترکیب ساختاری متمایز آن)، گونه‌های چوبی (چوب‌های نرم مانند کاج با مواد فرار بالا، چوب‌های سخت مانند راش با محتوای لیگنین بالا) و گیاهان انرژی اختصاصی (سویچ‌گراس و مایکانتوس اصلاح‌شده به طور خاص برای تبدیل گرماشیمیایی) را در برمی‌گیرند. شرایط عملیاتی در مجموعه آزمایش در سراسر پوشش تجربی گستره می‌یابند: دما از ۶۵۰ تا 880∘C ثبت‌کننده هر دو رژیم کنترل‌شده با سینتیک و نزدیک به تعادل، نسبت هم‌ارزی از ۰.۲۲ تا ۰.۴۸ دربرگیرنده عملیات غنی از سوخت تا غنی از اکسیدان و نسبت بخار به زیست‌توده از صفر (بدون افزودن بخار) تا ۱.۳ (گازی‌سازی تحت سلطه بخار). جدول ۵ معیارهای عملکرد تطبیقی را ارائه می‌دهد و تفاوت‌های قابل توجهی در دقت پیش‌بینی و نیازهای محاسباتی در بین چهار رویکرد مدل‌سازی آشکار می‌کند.

\begin{table}[t]
\centering
\small
\caption{عملکرد تطبیقی بر مجموعه‌های آزمایش مستقل}
\begin{tabularx}{\linewidth}{@{}*{3}{>{\raggedleft\arraybackslash}X}@{}}
\toprule
مدل & میانگین خطای مطلق (\%) & زمان محاسباتی \\
\midrule
مدل I (QET) & ۲۸.۵ & حدود ۱ ثانیه \\
مدل II (تجربی) & ۱۸.۳ & حدود ۱۰ ثانیه \\
مدل III (سینتیکی) & ۱۵.۷ & حدود ۲ ساعت (CFD) \\
\textbf(این کار) & \textbf(۸.۶) & \textbf(حدود ۵۰ میلی‌ثانیه) \\
\bottomrule
\end{tabularx}
\end{table}


چارچوب تابع خاص آگاه از فیزیک ما به میانگین خطای مطلق 8.6\% در تمام موارد آزمایشی دست می‌یابد، نشان‌دهنده تقریباً نصف خطای نشان‌داده‌شده توسط مدل III (رویکرد قبلی با بهترین عملکرد در خطای 15.7\%) علیرغم نیاز به زمان محاسباتی بیش از چهار مرتبه بزرگی کوتاه‌تر. این مزیت عملکرد به ویژه هنگام بررسی معاوضه خطا-زمان چشمگیر است: مدل I ارزیابی سریع (یک ثانیه برای هر مورد) به دست می‌آورد اما از دقت ضعیف به دلیل فرض تعادل در دمای تنظیم‌شده تجربی که فاقد توجیه فیزیکی است و قادر به ثبت محدودیت‌های سینتیکی در عملیات دمای پایین یا مقاومت‌های انتقال تأثیرگذار بر رفتار دمای بالا نیست، رنج می‌برد. مدل II دقت را به طور قابل توجهی از طریق هم‌بستگی‌های تجربی با دقت برازش‌شده بر داده آموزشی بهبود می‌بخشد، اما طبیعت درون‌یابی آن تعمیم به خوراک‌ها یا شرایط عملیاتی جدید خارج از توزیع آموزشی را محدود می‌کند و زمان ارزیابی ده ثانیه‌ای آن، کاربردهای بهینه‌سازی بلادرنگ را غیرممکن می‌سازد. مدل III به خطای قابل احترام 15.7\% از طریق مدل‌سازی مکانیکی دقیق ثبت‌کننده تغییرات مکانی و پویایی گذرا دست می‌یابد، اما نیازمندی محاسباتی دو ساعته آن برای هر مورد، آن را برای بهینه‌سازی طراحی نیازمند هزاران ارزیابی، کمّی‌سازی عدم قطعیت از طریق نمونه‌گیری مونت‌کارلو یا کنترل فرآیند بلادرنگ پاسخ‌دهنده به بازخورد حسگر در مقیاس زمانی ثانیه تا دقیقه غیرعملی می‌سازد.

دستیابی همزمان چارچوب ما به دقت برتر با کارایی محاسباتی شدید – ترکیبی که نشان‌دهنده بیش از صرفاً بهبود افزایشی نسبت به روش‌های موجود است – کلاس‌های کاملاً جدیدی از کاربردها را که قبلاً برای سیستم‌های تبدیلی گرماشیمیایی غیرممکن در نظر گرفته می‌شدند، ممکن می‌سازد. بهینه‌سازی فرآیند بلادرنگ اکنون می‌تواند شرایط عملیاتی (دما، نسبت هم‌ارزی، نرخ تزریق بخار) را هر چند ثانیه بر اساس اندازه‌گیری‌های حسگر ترکیب گاز محصول و بازده تنظیم کند و بازده تبدیل را بیشینه کند یا توزیع محصول خاصی را با تکامل شرایط اقتصادی یا نیازهای پردازش پایین‌دست هدف قرار دهد. کمّی‌سازی عدم قطعیت جامع از طریق نمونه‌گیری مونت‌کارلو یا شبه مونت‌کارلو قابل مدیریت می‌شود و طراحی احتمالی دقیقی را ممکن می‌سازد که عدم قطعیت‌ها در خصوصیات سوخت (به ویژه مربوط به زیست‌توده که تغییرپذیری قابل توجهی از دسته به دسته نشان می‌دهد)، خطاهای اندازه‌گیری تأثیرگذار بر تعیین شرایط عملیاتی و تقریب‌های مدل‌سازی ذاتی در هر چارچوب کاهش‌یافته را حساب می‌کند. غربالگری توان‌پذیری بالا ترکیب‌های خوراک جدید کاوش صدها یا هزاران ترکیب ترکیب می‌تواند مخلوط‌های هم‌افزایی را شناسایی کند که عملکرد برتری نسبت به خوراک‌های تک‌جزئی نشان می‌دهند و به طور بالقوه ارزش اقتصادی را از جریان‌های پسماند کم‌کیفیت یا پسماندهای کشاورزی که در حال حاضر کمتر استفاده می‌شوند، آزاد می‌کند. این کاربردها – که مدت‌ها توسط محققان تصور می‌شد اما از نظر محاسباتی با رویکردهای مدل‌سازی قبلی غیرقابل دسترسی بودند – اکنون به روش معمول مهندسی تبدیل می‌شوند که توسط شتاب هزار برابری که چارچوب ما بدون قربانی کردن دقت پیش‌بینی یا سازگاری فیزیکی ارائه می‌دهد، ممکن شده است.

\subsection{تفسیر فیزیکی ناورداهای جهانی}

فراتر از عملکرد پیش‌بینی کمی، استخراج چارچوب از دو ناوردای جهانی مستقل – انحراف لگاریتمی پیوسته ξ و برابری ساختاری گسسته S – تفسیر فیزیکی عمیق‌تری را دعوت می‌کند که سازه‌های ریاضی انتزاعی را به پدیده‌های گرماشیمیایی ملموس پیوند می‌زند. این ناورداها نه از برازش منحنی تجربی بلکه از توپولوژی بنیادین توابع انتقال محدودشده توسط قوانین پایستگی ظهور می‌یابند، که نشان می‌دهد آن‌ها ویژگی‌های ضروری منیفولد قید را ثبت می‌کنند که پویایی گازی‌سازی بدون توجه به مقادیر پارامتر خاص یا جزئیات عملیاتی به آن محدود شده است.

میدان بدون مقیاس ξ یک ویژگی عددی قابل توجه نشان می‌دهد که فراتر از کاربرد فوری آن برای پیش‌بینی بازده است. پارامتر بهره تابع انتقال خام k استخراج‌شده از پارامترسازی انتگرال بیضوی در گستره حیرت‌آوری از 10−47 تا 10−60 در بین ۲۰۰ مورد در مجموعه داده ما در نوسان است – گسترده‌ای از سیزده مرتبه بزرگی که مقایسه مستقیم بهره را از نظر عددی غیرقابل مدیریت و از نظر فیزیکی غیرقابل تفسیر می‌کرد. هنگامی که از طریق تبدیل لگاریتمی ξ=ln(ΛS/ΛG) که در آن ΛS نشان‌دهنده بهره خاص مورد و ΛG نشان‌دهنده میانگین هندسی بهره در تمام موارد است نرمال شود، این گستره دینامیکی عظیم به یک میدان بدون مقیاس فرومی‌ریزد که در آن ξ در تمام رژیم‌های عملیاتی در مرتبه واحد باقی می‌ماند. این نرمال‌سازی بیش از صرفاً راحتی عددی دست می‌یابد؛ آشکار می‌کند که تغییرات ضربی در بهره (ناشی طبیعی از وابستگی‌های نمایی در سینتیک آرنیوس حاکم بر نرخ واکنش و تعادل‌های ترمودینامیکی تعیین‌کننده توزیع محصول) به صورت تغییرات جمعی در ξ ظاهر می‌شوند و روابطی را که در غیر این صورت انحناء دست و پا گیر در چند دهه نشان می‌دادند، خطی می‌کنند.

بزرگی عددی این بهره‌های خام، مشاهده‌ای گمانه‌زنانه اما جالب را دعوت می‌کند که مقیاس‌های فیزیکی ناهمگون را از طریق تحلیل ابعادی به هم پیوند می‌زند. در واحدهای طبیعی که سرعت نور و ثابت پلانک کاهش‌یافته برابر با واحد هستند، ثابت کیهانی Λ – حاکم بر انبساط شتاب‌دار جهان از طریق معادلات میدان اینشتین – بزرگی تقریباً 10−52 m−2 در واحدهای SI متداول نشان می‌دهد. هنگامی که در واحدهای طبیعی حذف‌کننده آثار ابعادی بیان شود، این مقیاس کیهانی به صورت عددی با بهره‌های تابع انتقال مشخصه‌یابی‌کننده سیستم‌های گازی‌سازی ما هم‌تراز است علیرغم هفده مرتبه بزرگی جداسازی راکتورهای آزمایشگاهی (حجم متر مکعب، توان کیلووات، جرم کیلوگرم) از انبساط کیهانی (فاصله مگا پارسک، چگالی انرژی معادل یوتاوات، مقیاس جرم خورشیدی). در حالی که ما به شدت ادعای ارتباط فیزیکی عمیقی بین تبدیل گرماشیمیایی و پویایی کیهانی نداریم – معادلات حاکم، تقارن‌ها و مکانیسم‌های فیزیکی اساساً متفاوت هستند – این تصادف عددی فرضیه تحریک‌کننده‌ای را پیشنهاد می‌دهد: سیستم‌های قانون پایستگی در مقیاس‌های بسیار متفاوت، از راکتورهای آزمایشگاهی تا انبساط کیهانی، ممکن است ساختارهای ابعادی بنیادین مشترکی داشته باشند که در پارامترسازی‌های تابع خاص کدگذاری شده‌اند. دستگاه‌سازی انتگرال بیضوی که در اصل برای مکانیک سماوی توسعه داده شد، به نظر می‌رسد محدودیت‌های هندسی جهانی تحمیل‌شده توسط قوانین پایستگی مستقل از تجسم فیزیکی خاص را ثبت می‌کند و به کلاس‌های هم‌ارزی احتمالی اشاره دارد که در آن سیستم‌های فیزیکی متمایز از نظر ریاضی تحت نرمال‌سازی مناسب هم‌ریخت هستند.

این مشاهده تحقیق نظری بیشتری را دعوت می‌کند که فراتر از محدوده کار حاضر است اما شایسته بیان صریح برای هدایت جهت‌های تحقیقاتی آینده است. آیا پارامترسازی‌های تابع خاص «جاذب‌هایی» برای پویایی‌های حاکم از قانون پایستگی به این معنا هستند که سیستم‌های فیزیکی ناهمگون، هنگامی که تحت کاهش ابعاد حفظ‌کننده کمیت‌های پایستار قرار می‌گیرند، به طور طبیعی به سمت توصیف‌های قابل بیان از طریق خانواده‌های تابع خاص یکسان تکامل می‌یابند بدون توجه به جزئیات میکروسکوپی؟ آیا ترکیبات بدون بعد مانند ξ نسبت مقیاس‌های مشخصه نرمال‌شده توسط میانگین‌های مجموعه – کلاس‌های هم‌ارزی تعریف می‌کنند که فضای تمام سیستم‌های قانون پایستگی را به رده‌های مجزایی که توسط ناورداهای توپولوژیکی متمایز شده‌اند، تقسیم می‌کنند؟ آیا می‌توان قضایای دقیقی اثبات کرد که ثابت می‌کنند ساختارهای قانون پایستگی خاص لزوماً ظهورهای تابع خاص خاصی را ایجاب می‌کنند، مشابه چگونگی که گروه‌های تقارن در مکانیک کوانتومی نمایندگی‌های خاصی را اجباری می‌کنند؟ این سؤالات، اگرچه گمانه‌زنانه، امکان‌های غنی برای یکپارچه‌سازی حوزه‌های مدل‌سازی ناهمگون از طریق شناسایی ساختار ریاضی مشترک زیرین تجلی‌های فیزیکی متفاوت سطحی را پیشنهاد می‌دهند.

بازگشت از این گمانه‌زنی‌های کیهانی به تفسیر فیزیکی فوری‌تر، ناوردای برابری ساختاری گسسته S=∣p∣−∣z∣∈(0,1) دو کلاس بنیادین رفتار گازی‌سازی را با تفسیر مکانیکی واضح متمایز می‌کند. موارد نشان‌دهنده S=0 (تعداد برابر قطب‌ها و صفرها در نمایش تابع انتقال) با مسیرهای واکنش متوازن مطابقت دارند که سیستم نزدیک به تعادل ترمودینامیکی با سینتیک رو به جلو-معکوس متقارن عمل می‌کند. در این سیستم‌ها، دما و ترکیب به طور هموار بدون انتقال‌های ناگهانی تکامل می‌یابند، اغتشاشات کوچک پاسخ‌های متناسب تولید می‌کنند و توزیع محصولات به طور نزدیکی با آن‌هایی که توسط کمینه‌سازی انرژی آزاد گیبس در دمای بستر اندازه‌گیری‌شده پیش‌بینی می‌شوند، تقریب می‌خورند. تعادل قطب-صفر یک تقارن ریاضی را بازتاب می‌دهد که اصل تعادل دقیق را کد می‌کند: در تعادل، هر واکنش ابتدایی با نرخ‌های یکسان در جهات رو به جلو و معکوس پیش می‌رود و پویایی برگشت‌پذیر زمانی ایجاد می‌کند که افزایش یا کاهش نیروی محرکه ترمودینامیکی به مقادیر کوچک پاسخ‌های متقارن تولید کند.

برعکس، موارد نشان‌دهنده S=1 (یک قطب اضافی فراتر از تعداد صفر، ایجاد درجه تابع انتقال نامتقارن) با سیستم‌های تحت سلطه واکنش‌های به شدت برگشت‌ناپذیر مطابقت دارند که در آن قیود تعادل دیگر توزیع محصول را حکم‌فرمایی نمی‌کنند. شکست قطران دمای بالا اساساً به صورت تک‌جهتی با سنتز معکوس ناچیز هیدروکربن‌های پیچیده از گازهای سبک پیش می‌رود، گازی‌سازی زغال کربن جامد را از طریق واکنش‌های سطحی بدون رسوب کربن قابل توجه از پیش‌سازهای فاز گاز مصرف می‌کند و واکنش جابجایی آب-گاز دور از تعادل با پتانسیل ترمودینامیکی بزرگ که قادر به آرامش در مقیاس‌های زمان اقامت راکتور نیست، پیش می‌رود. این عدم تقارن در توپولوژی قطب-صفر، درجات آزادی اضافی را کد می‌کند که امکان بازده تبدیل کربن بالاتر را فراهم می‌کنند: قطب اضافی نشان‌دهنده یک مسیر اتلافی اضافی است که از طریق آن سیستم می‌تواند انرژی آزاد را رها کند و امکان دسترسی به حالت‌های محصول غیرقابل دسترس تحت قیود تعادل را فراهم می‌کند. همبستگی بین S=1 و بازده تبدیل افزایش‌یافته (معمولاً ۱۵-۱۰ واحد درصد بالاتر از موارد S=0 در دمای معادل و نسبت هم‌ارزی) واقعیت فیزیکی را بازتاب می‌دهد که دستیابی به گازی‌سازی کامل نیازمند شکستن قیود تعادل ترمودینامیکی از طریق شیمی برگشت‌ناپذیر دمای بالا است که به طور مؤثر محصولات را از منطقه واکنش قبل از اینکه واکنش‌های معکوس بتوانند واکنش‌گرها را بازتولید کنند، حذف می‌کند.

طبیعت باینری S که تنها مقادیر ۰ یا ۱ را می‌گیرد علیرغم ظهور از بهینه‌سازی پیوسته پارامترهای انتگرال بیضوی – به ویژه قابل توجه است. ممکن است انتظار داشته باشیم که درجه تابع انتقال با تغییر تدریجی شرایط عملیاتی به طور پیوسته تغییر کند، اما در عوض ما شاهد انتقال‌های توپولوژیکی تیز هستیم که در آن ساختار قطب-صفر بین پیکربندی‌های متوازن و نامتقارن سوئیچ می‌کند زیرا دما از آستانه‌های مشخصه عبور می‌کند (معمولاً ۸۰۰-۷۵۰ درجه سلسیوس برای اکثر خوراک‌ها) یا نسبت هم‌ارزی از مقادیر بحرانی فراتر می‌رود (معمولاً ER >0.35 جداکننده رژیم‌های کنترل‌شده با سینتیک و محدودشده با انتقال). این انتقال‌های گسسته شبیه به انتقال فاز در مکانیک آماری هستند که در آن پارامترهای نظم به طور ناپیوسته تغییر می‌کنند علیرغم تغییر هموار در متغیرهای شدید ترمودینامیکی، که نشان می‌دهد پویایی گازی‌سازی شکلی از رفتار بحرانی نشان می‌دهد که در آن رژیم‌های عملیاتی متمایز به طور تیز ترسیم می‌شوند به جای اینکه به طور پیوسته در هم بیامیزند.

سومین کمیت تشخیصی معرفی‌شده در بخش ۴ – زمان محاسباتی ویژه (sct) – تفسیر اضافی فراتر از تعریف عملیاتی آن به عنوان زمان سپری‌شده برای مرتب‌سازی عناصر ماتریس را می‌طلبد. علیرغم منشأ الگوریتمی، sct همبستگی‌های سیستماتیکی با خصوصیات سیستم با معنی فیزیکی نشان می‌دهد که فایده آن را به عنوان نماینده عدم تعادل تأیید می‌کند. موارد با مقادیر sct افزایش‌یافته (نیازمند زمان مرتب‌سازی طولانی‌تر، معمولاً فراتر از 10−3 ثانیه) به طور سازگار چندین ویژگی مشخصه نشان می‌دهند: بازده قطران افزایش‌یافته نشان‌دهنده تجزیه ناقص حدواسط‌های پیرولیز، که نشان می‌دهد مولکول‌های آلی پیچیده در طول راکتور باقی می‌مانند به جای شکست کامل به گازهای سبک؛ واکنش‌پذیری زغال پایین‌تر همان‌طور که توسط نرخ‌های گازی‌سازی ویژه اندازه‌گیری‌شده در مطالعات آزمایشگاهی مستقل کمّی شده است، نشان‌دهنده محدودیت‌های سینتیکی در واکنش‌های ناهمگن جامد-گاز؛ و پراکندگی بیشتر در تکرارهای تجربی اندازه‌گیری ترکیب گاز محصول یا بازده تحت شرایط عملیاتی اسمی یکسان، آشکارکننده حساسیت افزایش‌یافته به نوسانات کنترل‌نشده در پویایی سیال‌سازی، توزیع اندازه ذره سوخت یا فعالیت کاتالیزوری مواد بستر.

برعکس، موارد نشان‌دهنده مقادیر sct پایین (مرتب‌سازی سریع، معمولاً زیر 10−4 ثانیه) تمایل دارند به تعادل نزدیک‌تر شوند با توزیع محصولات مطابق با پیش‌بینی‌های کمینه‌سازی گیبس درون عدم قطعیت تجربی، مسیرهای واکنش ساده‌تر با مسیرهای رقابتی کمتر همان‌طور که توسط رفتار آرنیوس بدون وابستگی‌های دمایی غیرعادی یا انتقال رژیم شواهد شده است نشان می‌دهند و تکرارپذیری برتر در اندازه‌گیری‌های تجربی با اجراهای تکرارشده توافق‌کننده درون ۳-۲٪ علیرغم انجام آزمایش‌ها در روزهای مختلف یا با دسته‌های خوراک مختلف نشان می‌دهند. این همبستگی‌های سیستماتیک، sct را به عنوان یک نماینده معنادار برای شخصیت عدم تعادل سیستم تأیید می‌کنند و اطلاعات کمی درباره پیچیدگی سینتیکی بدون مدل‌سازی صریح مکانیسم‌های واکنش یا پدیده‌های انتقال که نیازمندی‌های محاسباتی را به طور چشمگیری افزایش می‌دادند، ارائه می‌دهند.

منشأ مکانیکی همبستگی‌های sct – چرا پیکربندی‌های ماتریس مشخصه خاص نیازمند پردازش محاسباتی طولانی‌تر هستند – تا حدی مرموز باقی می‌ماند و نمونه جالبی را نشان می‌دهد که در آن کشف داده‌محور مقدم بر درک نظری کامل است. بررسی‌های مقدماتی نشان می‌دهند که ماتریس‌های نشان‌دهنده sct بیشتر با طیف مقدار ویژه پیچیده‌تر با نزدیک‌انحطاطی‌هایی مطابقت دارند که نیازمند حساب‌گری با دقت گسترده‌شده برای حل تفاوت‌های عددی هستند، توزیع‌های درایه نامنظم باعث پیش‌بینی اشتباه انشعاب در الگوریتم‌های مرتب‌سازی بهینه‌شده برای ورودی‌های تقریباً مرتب یا یکنواخت-توزیع می‌شوند یا الگوهای دسترسی حافظه ایجادکننده خطاهای حافظه پنهان بیش از حد هستند زیرا عناصر ماتریس ارجاع‌داده‌شده توسط عملیات مرتب‌سازی در ذخیره‌سازی غیرپیوسته ثابت می‌شوند. این عوامل الگوریتمی، اگرچه به ظاهر از فیزیک گازی‌سازی جدا هستند، در نهایت ساختار ریاضی ماتریس‌های مشخصه را بازتاب می‌دهند که به نوبه خود از پردازش شبکه عصبی توالی‌های فیبوناچی کدکننده حالت سیستم ظهور می‌یابند. زنجیره علیت – از پیچیدگی فیزیکی از طریق کدگذاری زمانی سلسله‌مراتبی به ماتریس‌های وزن عصبی به دشواری مرتب‌سازی – نشان می‌دهد که چگونه تلاش محاسباتی می‌تواند به عنوان یک کاوشگر غیرمستقیم اما کمّی‌پذیر از پیچیدگی فیزیکی زیرین عمل کند حتی زمانی که درک مکانیکی مستقیم ناقص باقی می‌ماند.

\subsection{تحلیل رده‌بندی شده در گستره رژیم‌های عملیاتی و مشخصه‌های سوخت}

پس از اثبات عملکرد کلی کمی و تفسیر فیزیکی ناورداهای استخراج‌شده، اکنون بررسی می‌کنیم که چگونه دقت پیش‌بینی در رده‌های عملیاتی مختلف تغییر می‌کند تا نقاط قوت یا محدودیت‌های خاص حوزه بالقوه را شناسایی کنیم. این تحلیل رده‌بندی شده از طریق سه پارتیشن مکمل مجموعه داده پیش می‌رود: رفتار وابسته به دما آشکارکننده چگونگی تکامل عملکرد مدل با انتقال سیستم‌ها از رژیم کنترل‌شده با سینتیک به رژیم نزدیک‌شونده به تعادل، مقایسه نوع راکتور متمایزکننده پیکربندی‌های بستر سیال حبابی در مقابل در گردش با تفاوت‌های هیدرودینامیکی مشخصه آن‌ها و تغییرپذیری نوع سوخت ارزیابی‌کننده استحکام در بین رده‌های خوراک متنوع نشان‌دهنده ترکیبات عنصری، خصوصیات ساختاری و رفتار تجزیه گرمایی متمایز.

لایه‌بندی دمایی مجموعه داده اعتبارسنجی به سه رژیم – پایین (۷۰۰-۶۰۰ درجه سلسیوس)، متوسط (۸۰۰-۷۰۰ درجه سلسیوس) و بالا (۹۰۰-۸۰۰ درجه سلسیوس) – تغییرپذیری متوسط اما سیستماتیک در دقت پیش‌بینی را آشکار می‌کند که با درک مکانیکی پدیده‌های وابسته به دما هم‌تراز است. رژیم دمای پایین ضریب تعیین R2=0.79 برای پیش‌بینی‌های ترکیب نشان می‌دهد، کمی پایین‌تر از میانگین کلی مجموعه داده، بازتاب تغییرپذیری افزایش‌یافته معرفی‌شده توسط پیرولیز ناقص جایی که ساختار زغال تأثیر قابل توجهی بر سینتیک گازی‌سازی بعدی دارد، تشکیل قطران از طریق مسیرهای شکست و بسپارش رقابتی حساس به نرخ گرمایش و توزیع اندازه ذره و محدودیت‌های انتقال حرارت جایی که واکنش‌های گرماگیر تبخیرشدگی گرادیان‌های دمایی قابل توجهی بین هسته ذرات و سطوح ایجاد می‌کنند. این رژیم چالش‌برانگیزترین برای مدل‌سازی کاهش‌یافته ثابت می‌شود زیرا ناهمگونی‌های مکانی – نادیده گرفته شده در چارچوب حجم کنترل صفر بعدی ما – تأثیر حداکثری اعمال می‌کنند هنگامی که مقیاس‌های زمانی سینتیکی (ده‌ها ثانیه برای گازی‌سازی زغال) از مقیاس‌های زمانی اختلاط (ثانیه‌ها برای گردش حباب) فراتر می‌روند و فرض اختلاط کامل زیرین معادلات پایستگی میانگین‌گیری حجمی را نقض می‌کنند.

رژیم دمای متوسط به قوی‌ترین عملکرد با R2=0.84 برای پیش‌بینی‌های ترکیب دست می‌یابد، نشان‌دهنده نقطه شیرین که در که سینتیک‌ها به اندازه کافی سریع هستند اما هنوز آنقدر سریع نیستند که توزیع‌های نامتعادل محصول ایجاد کنند. در رژیم دمای بالا (بالای ۸۰۰ درجه سلسیوس)، عملکرد کمی کاهش می‌یابد (R2=0.80)، که بازتاب دهنده دشواری در پیش‌بینی دقیق ترکیب زمانی است که سیستم به تعادل ترمودینامیکی نزدیک می‌شود و تغییرات کوچک در دما یا نسبت هم‌ارزی، تغییرات نسبتاً بزرگی در محصولات ایجاد می‌کنند. در این رژیم، واکنش‌های تعادلی مانند جابجایی آب-گاز اهمیت بیشتری پیدا می‌کنند و پیش‌بینی نیازمند مدل‌سازی دقیق‌تر این تعادل‌ها است.

مقایسه بین راکتورهای بستر سیال حبابی (BFB) و در گردش (CFB) نشان می‌دهد که مدل در هر دو پیکربندی عملکرد خوبی دارد، اما در موارد BFB کمی بهتر است (R2=0.83 در مقابل 0.80 برای CFB). این تفاوت احتمالاً به دلیل پیچیدگی‌های هیدرودینامیکی بیشتر در CFB است، مانند جریان سریع‌تر، اختلاط شدیدتر و رفتار انتقال حرارت متفاوت. با این حال، عملکرد کلی در هر دو نوع قابل قبول است و نشان می‌دهد که چارچوب می‌تواند ویژگی‌های اصلی را در هر دو سیستم ثبت کند.

بررسی عملکرد بر اساس نوع سوخت (چوبی، پسماند کشاورزی، گیاهان انرژی) نشان می‌دهد که مدل در تمام رده‌ها عملکرد مشابهی دارد (R2 بین ۰.۷۹ تا ۰.۸۲). این نشان‌دهنده قابلیت تعمیم خوب مدل به سوخت‌های مختلف است، اگرچه برای برخی سوخت‌های با ترکیب بسیار خاص (مانند پسماندهای با خاکستر بسیار بالا یا رطوبت بسیار زیاد) ممکن است نیاز به داده آموزشی بیشتر یا تنظیم جزئی داشته باشد.

\subsection{تحلیل خطا: اریب‌های سیستماتیک، داده‌های پرت و حالت‌های شکست}

برای درک جامع قابلیت‌های مدل، نه تنها کمّی کردن عملکرد میانگین بلکه شناسایی الگوهای خطای سیستماتیک، تشخیص موارد پرت نشان‌دهنده پیش‌بینی‌های غیرعادی و مشخصه‌یابی حالت‌های شکست که در آن فرضیات چارچوب از هم می‌پاشند، ضروری است. این تحلیل خطا از طریق سه بررسی مکمل پیش می‌رود: تحلیل باقیمانده بررسی می‌کند که آیا خطاهای پیش‌بینی روندهایی با توجه به پارامترهای عملیاتی یا خصوصیات سوخت نشان می‌دهند که نشان‌دهنده فیزیک مدل‌نشده یا فرم‌های تابعی نامناسب باشد، شناسایی داده‌های پرت و تحلیل علت ریشه‌ای تعیین‌کننده ویژگی‌های مشترک در بین موارد تولیدکننده خطاهای پیش‌بینی به طور غیرعادی بزرگ، و کمّی‌سازی عدم قطعیت ارائه‌دهنده کران‌های عملکرد احتمالی ضروری برای کاربردهای طراحی ریسک‌گریز.

تحلیل باقیمانده – محاسبه پیش‌بینی منهای اندازه‌گیری برای هر مورد اعتبارسنجی و بررسی همبستگی بین این باقیمانده‌ها و پارامترهای سیستم – هیچ اریب سیستماتیکی با توجه به متغیرهای عملیاتی اولیه آشکار نمی‌کند. باقیمانده‌های دما (خطاهای در پیش‌بینی در مقابل اندازه‌گیری ترکیب به عنوان تابعی از دمای عملیاتی) به طور یکنواخت درون ±۳ درصد حجمی در سراسر گستره کامل ۹۰۰-۶۰۰ درجه سلسیوس توزیع می‌شوند بدون نشان دادن روندهایی مانند خطاهای مثبت فزاینده در دمای بالا (که نشان‌دهنده کم‌برآوردی سیستماتیک نزدیکی به تعادل خواهد بود) یا خطاهای منفی فزاینده در دمای پایین (پیشنهادکننده بیش‌برآوردی سینتیک تبدیل). به طور مشابه، باقیمانده‌های نسبت هم‌ارزی هیچ اریب ترجیحی به سمت شرایط غنی از سوخت یا غنی از هوا نشان نمی‌دهند – خطاها به همان اندازه احتمال مثبت یا منفی بودن در ER=0.25 (به شدت زیراستوکیومتری، تحت سلطه گازی‌سازی) به اندازه ER=0.45 (نزدیک به نسبت استوکیومتری احتراق که اکسایش با گازی‌سازی رقابت می‌کند) دارند. این عدم وجود روندهای سیستماتیک، اطمینان می‌دهد که مدل روابط فیزیکی زیرین را ثبت می‌کند به جای حفظ کردن آثار داده آموزشی که هنگام برون‌یابی فراتر از توزیع‌های آموزشی به صورت اریب‌های وابسته به منطقه ظاهر می‌شدند.

باقیمانده‌های نوع سوخت به طور مشابه هیچ عملکرد ترجیحی به نفع رده‌های خوراک خاص نشان نمی‌دهند. زیست‌توده چوبی، پسماندهای کشاورزی و گیاهان انرژی همه توزیع خطای قابل مقایسه‌ای با میانگین خطای مطلق ۲.۵-۲.۲ درصد حجمی نشان می‌دهند، نشان می‌دهد که یکپارچه‌سازی خصوصیت سوخت مبتنی بر CCA به طور مؤثر تفاوت‌های ترکیبی را که در غیر این صورت ممکن است اریب‌های سیستماتیک ایجاد کنند، نرمال می‌کند. موارد با خصوصیات سوخت شدید – خاکستر بسیار بالا ( >15\% ، مشخصه پوسته برنج)، خاکستر بسیار پایین ( <2\% ، مشخصه تراشه چوب تمیز)، نیتروژن بالا ( >2\% ، یافت‌شده در پسماندهای غنی از پروتئین مانند دانه‌های تقطیر) یا نسبت‌های عنصری غیرمعمول (H/C بالا مشخصه دانه‌های غنی از لیپید) – به طور سیستماتیک کم‌ یا بیش‌پیش‌بینی نمی‌کنند در مقایسه با سوخت‌های معمولی‌تر، تأیید می‌کند که استخراج ویژگی هندسی در فضای مثلثاتی «آب‌های آزاد» پارامترسازی مستحکمی در سراسر پوشش ترکیبی نمایندگی‌شده در داده آموزشی ارائه می‌دهد.

علیرغم این عدم وجود اریب سیستماتیک کلی، نه مورد در بین ۲۵۰ مورد کل (شامل اجراهای اکتشافی فراتر از مجموعه اعتبارسنجی هسته) خطاهای پیش‌بینی فراتر از دو برابر میانگین مجموعه داده نشان دادند، که شایسته بررسی دقیق برای شناسایی ویژگی‌های حالت شکست مشترک است. تحلیل علت ریشه‌ای آشکار می‌کند که این داده‌های پرت یک یا چند مورد از سه ویژگی مشکل‌آفرین نشان‌دهنده نقض فرضیات چارچوب یا محدودیت‌ها در پروتکل‌های مشخصه‌یابی سوخت را به اشتراک می‌گذارند. اولین ویژگی مشکل‌آفرین شامل محتوای خاکستر بسیار بالا فراتر از ۲۰٪ جرمی است، یافت‌شده در نمونه‌های خاص پوسته برنج و مخلوط‌های زغال سنگ-زیست‌توده جایی که مواد معدنی به سطوحی نزدیک‌تر به لیگنیت می‌رسند تا زیست‌توده. شیمی خاکستر – شامل برهم‌کنش‌های پیچیده بین سیلیکات‌ها، آلومینوسیلیکات‌ها، فلزات قلیایی (سدیم، پتاسیم)، فلزات قلیایی خاکی (کلسیم، منیزیم) و گونه‌های کم‌مقدار (آهن، فسفر) – بر گازی‌سازی از طریق اثرات کاتالیزوری بر واکنش‌پذیری زغال (گونه‌های قلیایی به شدت نرخ گازی‌سازی را از طریق مکانیسم‌های شامل جذب سطحی اکسیژن و اختلال شبکه کربن تسریع می‌کنند)، تقسیم‌بندی زغال-مواد فرار در طول پیرولیز (مواد معدنی بر شکست ثانویه و واکنش‌های تراکم حاکم بر تشکیل قطران تأثیر می‌گذارند) و تعادل‌های گاز محصول (اکسید کلسیم تشکیل‌شده از تجزیه کلسیت می‌تواند نسبت CO/CO2 را از طریق کاتالیز واکنش بودوار جابجا کند) تأثیر می‌گذارد. چارچوب CCA ما کسر جرمی کل خاکستر را دربر می‌گیرد اما گونه‌های معدنی مختلف را متمایز نمی‌کند یا اثرات شیمی خاکستر پیچیده را حساب نمی‌کند، که توضیح‌دهنده تخریب عملکرد برای سوخت‌های با خاکستر بالا جایی که این پدیده‌ها غالب هستند، است.

دومین ویژگی مشکل‌آفرین شامل دمای عملیاتی بسیار پایین زیر ۶۵۰ درجه سلسیوس است که در آن پیرولیز ناقص زغال جزئی-تبخیرشده با محتوای فرار قابل توجه تولید می‌کند، تشکیل قطران به حداکثر بازده از طریق مسیرهای شکست و بسپارش رقابتی بسیار حساس به تاریخچه گرمایش می‌رسد و تراکم هیدروکربن‌های سنگین در مناطق خنک‌تر راکتور یا خطوط نمونه‌برداری آثار اندازه‌گیری گیج‌کننده‌ای معرفی می‌کند که توزیع محصول گزارش‌شده را مخدوش می‌کنند. این شرایط دمای پایین، فرضیات ضمنی چارچوب را نقض می‌کنند که تبخیرشدگی به سرعت تا قبل از شروع گازی‌سازی پیش می‌رود (ساده‌سازی مدل‌سازی استاندارد که پیرولیز و گازی‌سازی را به عنوان فرآیندهای متوالی به جای هم‌پوشان در نظر می‌گیرد) و این که گاز محصول به تعادل ترکیبی درون بستر می‌رسد (به طور فزاینده نامعتبر با کاهش دما و تشدید محدودیت‌های سینتیکی). تخریب عملکرد متوسط در ۷۰۰-۶۰۰ درجه سلسیوس اشاره‌شده در تحلیل رده‌بندی شده این پدیده‌ها را بازتاب می‌دهد، اما زیر ۶۵۰ درجه سلسیوس فروپاشی فرض به اندازه کافی شدید می‌شود تا گهگاه خطاهای پیش‌بینی بزرگ تولید کند.

سومین ویژگی مشکل‌آفرین شامل گازی‌سازی غنی‌شده با اکسیژن با غلظت اکسیژن ورودی فراتر از ۴۰٪ حجمی است، نزدیک یا فراتر از نیازهای استوکیومتری برای احتراق کامل. تحت این شرایط به شدت اکسایشی، واکنش‌های احتراق گرمازا (اکسایش کربن و هیدروژن آزادکننده گرمای قابل توجه) رقابتی یا غالب بر واکنش‌های گازی‌سازی گرماگیر (زغال-بخار و زغال-CO مصرف‌کننده گرما) می‌شوند و گرادیان‌های دمایی بزرگ ایجاد می‌کنند زیرا مناطق احتراق داغ متناوب با مناطق گازی‌سازی خنک‌تر بسته به در دسترس بودن اکسیژن محلی هستند. این ناهمگونی‌های گرمایی، فرض هم‌دما بودن زیرین فرمول‌بندی موازنه انرژی ما را نقض می‌کنند (که انتقال حرارت خالص را محاسبه می‌کند اما دمای یکنواخت در سراسر حجم کنترل را فرض می‌کند) و شیمی مختلط احتراق-گازی‌سازی محصولاتی معرفی می‌کند (عمدتاً آب و CO2 از اکسایش کامل) که گازهای مشتق‌شده از سوخت را رقیق می‌کنند و تفسیر معیارهای بازده تبدیل را پیچیده می‌کنند. در حالی که چارچوب تا ER 0.5 که گازی‌سازی به وضوح غالب است قابل اعمال باقی می‌ماند، سطوح اکسیژن بالاتر نیازمند حسابداری صریح آزادسازی گرمای احتراق و توزیع دمای مکانی فراتر از محدوده فرمول‌بندی صفر بعدی ما است.

این مشخصه‌های حالت شکست اولویت‌های واضحی را برای پالایش آینده چارچوب پیشنهاد می‌دهند. مدل‌سازی صریح شیمی خاکستر می‌تواند ترکیب معدنی دقیق (نیازمند تجزیه و تحلیل جذب اتمی یا فلورسانس پرتو ایکس) و اثرات کاتالیزوری شناخته‌شده مستند در ادبیات سینتیک را دربر بگیرد، احتمالاً از طریق عوامل تصحیح ضربی اعمال‌شده به بازده تبدیل پیش‌بینی‌شده یا از طریق ویژگی‌های هندسی اصلاح‌شده کدکننده پارامترهای مرتبط با خاکستر فراتر از محتوای معدنی کل. سینتیک توسعه‌یافته دمای پایین ممکن است فرض تبخیرشدگی سریع را با معرفی بازده فرار وابسته به زمان اقامت یا جفت‌شدگی دو مرحله‌ای پیرولیز-گازی‌سازی که هم‌پوشانی بین فرآیندها به طور صریح مدل‌سازی می‌شود به جای فرض ناچیز بودن، شل کند. تصحیح‌های موازنه انرژی غیر هم‌دما برای شرایط به شدت گرمازا می‌تواند تغییرات دمای مکانی را از طریق مدل‌های جریان پیستونی یا دو ناحیه ساده‌شده ارائه‌دهنده دمای مشخصه نقطه داغ و منطقه خنک به جای یک مقدار میانگین بستر واحد برآورد کند و محاسبات تعادل ترمودینامیکی دقیق‌تری را حتی زمانی که اختلاط گرمایی کامل شکست می‌خورد، ممکن سازد.

کمّی‌سازی عدم قطعیت – ضروری برای کاربردهای طراحی ریسک‌گریز که کران‌های عملکرد احتمالی به اندازه پیش‌بینی‌های نقطه‌ای مهم ثابت می‌شوند – علیرغم سرراست بودن مفهومی، در پیاده‌سازی فعلی چارچوب توسعه نیافته باقی می‌ماند. تحلیل‌های bootstrap مقدماتی، که در آن داده آموزشی به طور مکرر با جایگزینی نمونه‌گیری مجدد می‌شود تا مجموعه‌ای از مدل‌های بازتاب‌دهنده عدم قطعیت آماری در پارامترهای یادگرفته‌شده تولید شود، فاصله اطمینان ۹۵٪ تقریبی در حدود ±۵ درصد حجمی برای پیش‌بینی‌های ترکیب (تقریباً دو برابر میانگین خطای مطلق، نشان‌دهنده این که نوسانات دو سیگما قابل مقایسه با خطاهای پیش‌بینی معمولی هستند)، ±۸\% برای بازده تبدیل کربن و ±۱۲\% برای بازده گاز محصول را پیشنهاد می‌کنند. این کران‌های اطمینان، اگرچه برآوردهای غیررسمی نیازمند برخورد بیزی دقیق برای تفسیر احتمالاتی قابل دفاع، نشان می‌دهند که عدم قطعیت برای اکثر کاربردهای مهندسی قابل مدیریت است – تصمیمات طراحی مستحکم در برابر تغییرات عملکرد ±۱۰\% با توجه به پیش‌بینی‌های چارچوب ما قابل اعتماد خواهند بود.

گسترش‌های بیزی پیچیده‌تر می‌توانند عدم قطعیت‌های پسین دقیقی با در نظر گرفتن تمام پارامترهای چارچوب – ضرایب CCA، پارامترهای تبدیل هندسی، وزن‌های شبکه عصبی، ضرایب رگرسیون انتگرال بیضوی – به عنوان متغیرهای تصادفی با توزیع‌های پیشین بازتاب‌دهنده دانش اولیه و توزیع‌های پسین به‌روزشده از طریق استنتاج مبتنی بر درست‌نمایی با استفاده از داده تجربی ارائه دهند. چنین پیاده‌سازی‌های بیزی، بهینه‌سازی طراحی احتمالاتی کامل بیشینه‌کننده عملکرد مورد انتظار در حالی که احتمالات شکست را زیر آستانه‌های قابل قبول محدود می‌کند، آزمون فرض رسمی مقایسه پیکربندی‌های جایگزین راکتور یا استراتژی‌های عملیاتی با سطوح اطمینان کمّی‌شده و طراحی آزمایشی متوالی که منابع اندازه‌گیری برای بیشینه‌سازی کاهش عدم قطعیت پسین تخصیص داده می‌شوند، ممکن می‌سازند و در نتیجه تلاش مشخصه‌یابی را در حالی که اطمینان مشخص‌شده در تصمیمات طراحی حاصل می‌شود، کمینه می‌کنند. این گسترش‌های احتمالاتی، اگرچه نیازمندی‌های محاسباتی را به طور متوسط افزایش می‌دهند (مونت‌کارلو زنجیره مارکوف یا استنتاج تغییراتی نیازمند هزاران ارزیابی درست‌نمایی) با توجه به زمان ارزیابی در مقیاس میلی‌ثانیه‌ای چارچوب ما قابل مدیریت باقی می‌مانند و گام طبیعی بعدی را برای بلوغ چارچوب به سمت استقرار صنعتی نشان می‌دهند که در آن تصمیم‌گیری آگاه از عدم قطعیت به اندازه دقت پیش‌بینی نقطه‌ای مهم ثابت می‌شود.

\subsection{سنتز: اعتبارسنجی قابلیت اجرای چارچوب برای کاربرد عملی را تأیید می‌کند}

اعتبارسنجی جامع ارائه‌شده در سراسر این بخش – معیارهای عملکرد کمی نشان‌دهنده دقت پیش‌بینی قوی، ارزیابی تطبیقی آشکارکننده مزایای قابل توجه نسبت به رویکردهای موجود، تفسیر فیزیکی پیونددهنده ناورداهای ریاضی به پدیده‌های گرماشیمیایی، تحلیل رده‌بندی شده تأییدکننده استحکام در سراسر رژیم‌های عملیاتی و انواع سوخت و تحلیل خطا شناسایی‌کننده محدودیت‌ها و حالت‌های شکست واضح – ثابت می‌کند که چارچوب تابع خاص آگاه از فیزیک به اهداف مرکزی خود دست می‌یابد. چارچوب همزمان دقت کمی قوی با پیش‌بینی‌های ترکیب دستیابی به $R^{2}>0.78$ برای تمام گونه‌های گاز اصلی، پیش‌بینی‌های بازده تبدیل کربن دستیابی به R2=0.87 و پیش‌بینی‌های بازده گاز محصول دستیابی به R2=0.82 ارائه می‌دهد، که همگی به طور قابل توجهی عملکرد مدل‌های موجود را پشت سر می‌گذارند. کارایی محاسباتی شدید را با ارزیابی نیازمند حدود ۵۰ میلی‌ثانیه برای هر مورد حفظ می‌کند، نشان‌دهنده شتاب 105−106× در مقایسه با CFD با وفاداری بالا و ممکن‌ساز بهینه‌سازی بلادرنگ، کمّی‌سازی عدم قطعیت جامع و کاربردهای غربالگری توان‌پذیری بالا که قبلاً غیرممکن بودند. سازگاری فیزیکی از طریق فرمول‌بندی تابع انتقال تحمیل‌کننده قوانین پایستگی به‌طور ساختاری به جای الگوریتمی تضمین می‌شود، با ناورداهای استخراج‌شده نشان‌دهنده همبستگی‌های یکنوایی که بهینه‌های موضعی مشکل‌آفرین برای بهینه‌سازی مبتنی بر گرادیان را حذف می‌کنند. چارچوب کلیت را از طریق عملکرد مستحکم در سراسر دماهای متنوع (۹۰۰-۶۰۰ درجه سلسیوس)، انواع راکتور (بسترهای سیال حبابی و در گردش) و خوراک‌ها (پسماندهای کشاورزی، گونه‌های چوبی، گیاهان انرژی) بدون نیاز به تنظیم خاص رژیم یا پارامترسازی جداگانه نشان می‌دهد. در نهایت، تفسیرپذیری را از طریق ناورداهای با معنی فیزیکی – ξ همبسته با بازده، S با بازده تبدیل، sct با شخصیت عدم تعادل – ارائه می‌دهد که درک مکانیکی را به جای تجربی‌گرایی جعبه سیاه ممکن می‌سازد.

این قابلیت‌های اعتبارسنجی‌شده، چارچوب را به عنوان یک جایگزین عملی رسیدگی‌کننده به شکاف‌های بحرانی در منظره مدل‌سازی موجود قرار می‌دهند. برای کاربردهایی که هزینه محاسباتی، CFD با وفاداری بالا را غیرممکن می‌سازد – بهینه‌سازی فرآیند کاوش هزاران ترکیب پارامتر، کنترل بلادرنگ پاسخ‌دهنده به تغییرات لحظه‌ای، ارزیابی سریع سناریوهای مختلف – این چارچوب یک ابزار قدرتمند ارائه می‌دهد.\textbf(۵ قابلیت تعمیم از طریق شباهت‌های ساختاری و تأثیر گسترده‌تر)

ساختار ریاضی چارچوب – کاهش ابعاد داده‌محور حفظ‌کننده قیود جمع به یک و پایستگی، پارامترسازی تابع خاص، استخراج ناورداهای جهانی – به طور طبیعی به حوزه‌های فراتر از تبدیل گرماشیمیایی تعمیم می‌یابد. شبکه‌های زیست‌شیمیایی (مسیرهای متابولیکی که متابولیت‌ها را با موازنه جرم و انرژی تبدیل می‌کنند)، سیستم‌های بوم‌شناختی (چرخه مواد مغذی تحت پایستگی عنصری، پویایی جمعیت تحت قیود منابع)، مدل‌سازی اقتصادی (جریان وجوه تحت موازنه حسابداری، تخصیص بودجه تحت قیود جمع به یک) و مدل‌سازی مواد (چندمقیاسی از اتمی تا پیوسته با قیود هندسی و پایستگی) همه ویژگی‌های ساختاری مشترکی را به اشتراک می‌گذارند که چارچوب ما برای آن‌ها مناسب است. در هر مورد، چالش ترکیب داده‌های ناهمگن، تحمیل قیود پایستگی و استنتاج سریع باقی می‌ماند، و رویکرد ما راه حلی ارائه می‌دهد که از دقت تحلیلی و کارایی عددی توابع خاص بهره می‌برد.

برای مثال، در مدل‌سازی متابولیک، CCA می‌تواند داده‌های ترانسکریپتومیک، پروتئومیک و متابولومیک را یکپارچه کند، تبدیل هندسی کسرهای شار متابولیکی را در فضای «آب‌های آزاد» پارامترسازی می‌کند و انتگرال‌های بیضوی پویایی سیستم را تحت قیود موازنه جرم و انرژی ثبت می‌کنند. در بوم‌شناسی، CCA داده‌های سنجش از دور و اندازه‌گیری‌های در محل را ترکیب می‌کند، تبدیل هندسی فراوانی گونه را در فضای ترکیب پارامترسازی می‌کند و انتگرال‌های بیضوی پاسخ بوم‌سامانه به محرک‌های محیطی را مدل می‌کنند. در اقتصاد، CCA شاخص‌های کلان اقتصادی را یکپارچه می‌کند، تبدیل هندسی سهم‌های بازار یا بودجه را پارامترسازی می‌کند و انتگرال‌های بیضوی تأثیر مداخلات سیاستی را ثبت می‌کنند.

با این حال، تعمیم موفقیت‌آمیز نیازمند اعتبارسنجی تجربی دقیق در هر حوزه جدید است. شباهت‌های ساختاری ریاضی شرط لازم اما نه کافی برای قابلیت کاربرد هستند؛ فیزیک خاص، شیمی یا زیست‌شناسی هر سیستم باید از طریق داده تجربی مناسب در چارچوب گنجانده شود. کار ما نقشه راه روش‌شناختی را ارائه می‌دهد، اما پیاده‌سازی در هر حوزه نیازمند سازگاری‌های خاص زمینه و اعتبارسنجی گسترده است.\textbf(۶ نتیجه‌گیری‌ها، محدودیت‌ها و جهت‌های آینده)

این کار ثابت می‌کند که پارامترسازی تابع خاص آگاه از فیزیک یک رویکرد عملی و مزیت‌دار برای مدل‌سازی کاهش‌یافته سیستم‌های حاکم از قانون پایستگی ارائه می‌دهد، نشان‌داده شده از طریق اعتبارسنجی دقیق بر گازی‌سازی زیست‌توده اما به طور بالقوه قابل تعمیم به حوزه‌های متنوع دارای ساختار ریاضی مشابه. چارچوب تحلیل همبستگی متعارف برای یکپارچه‌سازی داده چندگروهی، تبدیل هندسی ایجادکننده فضای تحلیل عددی پایدار «آب‌های آزاد»، کدگذاری سلسله‌مراتبی فیبوناچی ثبت‌کننده ساختار زمانی چندمقیاسی، آزمایشگاه‌های محاسباتی شبکه عصبی استخراج‌کننده عدم تقارن جهت‌دار و پارامترسازی انتگرال بیضوی کامل از طریق میانگین حسابی-هندسی تحمیل‌کننده قیود پایستگی ساختاری را ترکیب می‌کند. این ترکیب دستاوردهای همزمانی را به دست می‌آورد که به ندرت در رویکردهای مدل‌سازی منفرد یکپارچه شده‌اند: دقت کمی قوی با $R^{2}>0.78$ برای تمام خروجی‌های پیش‌بینی‌شده، کارایی محاسباتی شدید ممکن‌ساز ارزیابی میلی‌ثانیه‌ای نشان‌دهنده شتاب 105−106× در مقابل CFD، سازگاری فیزیکی تضمین‌شده از طریق فرمول‌بندی تابع انتقال احترام‌گذار به انفعال، پایداری و علیت، کلیت مستحکم در سراسر شرایط عملیاتی و انواع سوخت متنوع بدون تنظیم خاص رژیم و تفسیرپذیری معنادار از طریق ناورداهای استخراج‌شده نشان‌دهنده همبستگی‌های فیزیکی واضح.

مشارکت‌های روش‌شناختی کلیدی شایسته شمارش سیستماتیک برای روشن کردن پیشرفت‌های این کار فراتر از ادبیات موجود هستند. اول، نشان می‌دهیم که توابع خاص – به طور سنتی اختصاص‌یافته به جواب‌های فرم بسته معادلات دیفرانسیل تحلیلی-قابل مدیریت – می‌توانند به طور مؤثر مدل‌های کاهش‌یافته داده‌محور را هنگامی که تحمیل قانون پایستگی و کارایی محاسباتی از اهمیت بالایی برخوردارند پارامترسازی کنند، حتی زمانی که آن توابع خاص مستقیماً از معادلات حاکم نشأت نمی‌گیرند. این تعمیم مفهومی از «توابع خاص به عنوان جواب» به «توابع خاص به عنوان پارامترسازی» استراتژی‌های مدل‌سازی جدیدی برای سیستم‌هایی که مشتق اصول اول در آن‌ها حل‌ناپذیر ثابت می‌شود، باز می‌کند. دوم، اعتبارسنجی تجربی جامعی بر روی بیش از ۲۰۰ مجموعه داده گازی‌سازی ارائه می‌دهیم که به عملکردی فراتر از مدل‌های موجود (میانگین خطای مطلق 8.6\% در مقابل 15.7\% برای بهترین رویکرد قبلی) دست می‌یابد در حالی که بسته بودن پایستگی غیرممکن برای تضمین در چارچوب‌های یادگیری ماشین خالص حفظ می‌شود. سوم، زمان محاسباتی ویژه را به عنوان یک ویژگی پویای نوین معرفی می‌کنیم که «صلبیت محاسباتی» عدم تعادل را ثبت می‌کند و پل می‌زند بین پتانسیل‌های ترمودینامیکی ایستا و ثابت‌های نرخ سینتیکی بدون نیاز به داده سری زمانی یا اندازه‌گیری‌های گذرا. چهارم، دو ناوردای جهانی – انحراف لگاریتمی پیوسته ξ و برابری ساختاری گسسته S – را از توپولوژی تابع انتقال استخراج می‌کنیم که همبستگی‌های یکنوایی با کمیت‌های قابل مشاهده در سراسر گستره پارامتر کامل نشان می‌دهند و بهینه‌های موضعی مشکل‌آفرین برای بهینه‌سازی را حذف می‌کنند. پنجم، کشف می‌کنیم که بهره تابع انتقال از 10−47 تا 10−60 در بین موارد در نوسان است و پس از نرمال‌سازی یک میدان بدون مقیاس ایجاد می‌کند که به صورت عددی با بزرگی‌های ثابت کیهانی هم‌تراز است و ساختارهای ابعادی جهانی احتمالی مشترک بین سیستم‌های قانون پایستگی در مقیاس‌های بسیار متفاوت را پیشنهاد می‌دهد.

علیرغم این پیشرفت‌های قابل توجه، چارچوب محدودیت‌های واضحی نشان می‌دهد که مرزهای قابلیت کاربرد فعلی را تعریف می‌کنند و پالایش آینده را برمی‌انگیزند. مهم‌ترین محدودیت شامل محتوای خاکستر بالا فراتر از حدود ۲۰\% است، جایی که شیمی معدنی پیچیده (کاتالیز قلیایی، برهم‌کنش‌های گوگرد، تشکیل سیلیکات) بر سینتیک و ترمودینامیک گازی‌سازی از طریق مکانیسم‌هایی که توسط کمّی‌سازی خاکستر کل در پروتکل مشخصه‌یابی سوخت ما ثبت نشده‌اند، تأثیر می‌گذارد. تعمیم به خوراک‌های با خاکستر بالا نیازمند تحلیل کانی‌شناسی دقیق (پراش پرتو ایکس شناسایی‌کننده فازهای بلوری، جداسازی شیمیایی متمایزکننده مواد معدنی در دسترس در مقابل بی‌اثر) و مدل‌سازی صریح اثرات خاکستر از طریق اصلاحات نرخ کاتالیزوری یا اغتشاش‌های تعادلی است. محدودیت‌های مرتبط در شرایط عملیاتی شدید – دماهای زیر 650∘C جایی که پیرولیز ناقص فرض تبخیرشدگی سریع را نقض می‌کند یا غلظت اکسیژن فراتر از 40\% جایی که احتراق گرمازا ناهمگونی‌های گرمایی ایجاد می‌کند که موازنه انرژی هم‌دما را نقض می‌کند – ظهور می‌یابند، نشان می‌دهند که قابلیت کاربرد چارچوب پنجره‌های گازی‌سازی استاندارد را در برمی‌گیرد اما برای رژیم‌های غیرمتعارف نیازمند اصلاح است.

فرمول‌بندی حجم کنترل صفر بعدی، وضوح مکانی را قربانی می‌کند که به طور اجتناب‌ناپذیر هنگام میانگین‌گیری حجمی معادلات پایستگی توزیع‌شده مکانی از دست می‌رود، که فایده را برای کاربردهای نیازمند پیش‌بینی دقیق پروفایل‌های دما یا غلظت درون راکتورها، شناسایی نقاط داغ یا مناطق خنک تأثیرگذار بر انتخاب مواد یا تنش گرمایی، یا ارزیابی کیفیت اختلاط و توزیع زمان اقامت حاکم بر یکنواختی تبدیل محدود می‌کند. اگرچه برای پیش‌بینی عملکرد و بهینه‌سازی جایی که خروجی‌های میانگین‌گیری‌شده مکانی کافی ثابت می‌شوند کافی است، سناریوهای طراحی نیازمند جزئیات مکانی (مانند اندازه‌گیری سطوح تبادل گرما یا پیش‌بینی الگوهای فرسایش) نیازمند اتصال چارچوب ما به مدل‌های مکانی ساده‌شده (راکتورهای جریان پیستونی، نمایش‌های دو ناحیه‌ای) یا استفاده از آن به عنوان شرایط مرزی برای CFD موضعی متمرکز بر منابع محاسباتی بر مناطق جایی که وضوح مکانی ضروری ثابت می‌شود، هستند.

مدل‌سازی قطران نشان‌دهنده شکاف دیگری است که تصدیق شده – چارچوب غلظت گونه‌های گاز اولیه را پیش‌بینی می‌کند اما به طور صریح تشکیل قطران، سینتیک شکست یا گونه‌شناسی ترکیبی بحرانی برای طراحی تجهیزات پایین‌دست جایی که تراکم قطران باعث گرفتگی، خوردگی و اختلال فرآیند می‌شود، ردیابی نمی‌کند. تعمیم به پیش‌بینی قطران جامع نیازمند یا گسترش خروجی‌ها برای شامل معیارهای توده‌ای قطران (بازده کل قطران، وزن مولکولی متوسط) یادگرفته‌شده از مجموعه‌های داده گزارش‌کننده این کمیت‌ها، یا توسعه مدل‌های شیمی قطران کاهش‌یافته سازگار با قیود کارایی محاسباتی ما، احتمالاً از طریق توابع انتقال اضافی یا هم‌بستگی‌های تجربی پارامترشده توسط ناورداهای استخراج‌شده است.

کمّی‌سازی عدم قطعیت، اگرچه از نظر مفهومی از طریق گسترش‌های بیزی در نظر گرفتن پارامترهای چارچوب به عنوان متغیرهای تصادفی سرراست است، در پیاده‌سازی فعلی توسعه نیافته باقی می‌ماند و تنها فواصل اطمینان bootstrap مقدماتی ارائه می‌دهد به جای توزیع‌های پسین دقیق. نسخه‌های کاملاً احتمالاتی ممکن‌ساز طراحی ریسک‌گریز، آزمون فرض رسمی یا طراحی آزمایشی بهینه نیازمند مونت‌کارلو زنجیره مارکوف یا استنتاج تغییراتی نیازمند هزاران ارزیابی درست‌نمایی – قابل مدیریت با توجه به زمان ارزیابی در مقیاس میلی‌ثانیه‌ای ما اما نشان‌دهنده تلاش پیاده‌سازی قابل توجهی که در این نمایش اولیه انجام نشد – هستند. چنین گسترش‌های احتمالاتی گام طبیعی بعدی را برای بلوغ چارچوب به سمت استقرار صنعتی نشان می‌دهند که در آن تصمیم‌گیری آگاه از عدم قطعیت به اندازه دقت پیش‌بینی نقطه‌ای مهم ثابت می‌شود.

پویایی گذرا – توالی‌های راه‌اندازی، روش‌های خاموش‌کردن، عملیات دنبال‌کننده بار پاسخ‌دهنده به تقاضاهای گرمایی یا توان متغیر – خارج از محدوده فعلی که منحصراً به عملیات حالت پایا می‌پردازد، باقی می‌مانند. تعمیم به گذرا نیازمند معادلات تکامل زمانی برای جمله‌های انباشت در موازنه‌های حجم کنترل است، معرفی مشتقات زمانی که به طور کلی حتی برای فرمول‌بندی‌های کاهش‌یافته نیازمند حل تکراری هستند. با این حال، کارایی محاسباتی چارچوب پیشنهاد می‌دهد که حل بلادرنگ معادلات گذرا ممکن است از طریق روش‌های انتگرال‌گیری زمانی ضمنی ارزیابی‌کننده باقیمانده‌های حالت پایا هزاران بار در هر گام زمانی عملی باشد، به طور بالقوه شبیه‌سازی پویا و آزمایش سیستم کنترل را ممکن می‌سازد که در حال حاضر نیازمند کدهای تخصصی بلادرنگ است.

مقیاس‌گذاری از نصب‌های آزمایشگاهی و پایلوت (ظرفیت گرمایی کیلووات تا کم‌مگاوات) به سیستم‌های تجاری (ده‌ها مگاوات یا بزرگتر) پدیده‌هایی را معرفی می‌کند که در مجموعه داده اعتبارسنجی ما نمایندگی نشده‌اند: اثرات قطر بستر تغییردهنده توزیع اندازه حباب و عبور گاز، اثرات ارتفاع تغییردهنده توزیع زمان اقامت و پروفایل‌های دما، اثرات دیوار با کاهش نسبت سطح به حجم کمتر مهم می‌شوند و تفاوت‌های آماده‌سازی خوراک زیرا سیستم‌های صنعتی جریان‌های سوخت ناهمگن را پردازش می‌کنند به جای نمونه‌های آزمایشگاهی با دقت مشخصه‌یابی‌شده. در حالی که استدلال‌های اصول اول پیشنهاد می‌دهند که رویکرد مبتنی بر پایستگی ما باید قابل اطمینان‌تر از هم‌بستگی‌های صرفاً تجربی برونیابی شود – قوانین پایستگی بدون توجه به مقیاس اعمال می‌شوند – اعتبارسنجی تجربی بر مجموعه‌های داده در مقیاس تجاری قبل از استقرار مطمئن به طراحی در مقیاس کامل ضروری باقی می‌ماند.

جهت‌های تحقیقاتی آینده به طور طبیعی از این محدودیت‌ها و از قابلیت‌های اعتبارسنجی‌شده ایجادشده از طریق این کار ظهور می‌یابند، که در سه رده با گستره فزاینده قرار می‌گیرند. گسترش‌های فوری درون تبدیل گرماشیمیایی شامل تطبیق چارچوب به فرآیندهای مرتبط (پیرولیز تولیدکننده بیو-روغن‌ها به جای عمدتاً گازها، تورفیکاسیون تولیدکننده سوخت‌های جامد از طریق فرآیند گرمایی ملایم، پردازش هیدروترمال تبدیل زیست‌توده مرطوب از طریق واکنش‌های فاز آبی، looping شیمیایی چرخه‌ای حامل‌های اکسیژن اکسید فلزی) است که ساختار ریاضی را به اشتراک می‌گذارند علیرغم محصولات یا شرایط عملیاتی متفاوت. گسترش‌های میان‌رشته‌ای شامل اعمال چارچوب به حوزه‌هایی مانند مدل‌سازی متابولیک، بوم‌شناسی سیستم‌ها یا اقتصاد کلان است که در آن قیود جمع به یک و پایستگی غالب هستند. در نهایت، تحقیقات نظری عمیق‌تر می‌تواند مبانی ریاضی ارتباط بین ساختار قانون پایستگی و ظهورهای تابع خاص را بررسی کند، احتمالاً قضایای جهان‌شمولی را ثابت کند که توضیح می‌دهد چرا سیستم‌های بسیار متفاوت فیزیکی در نهایت پارامترسازی‌های تابع خاص مشابهی را می‌پذیرند.

در مجموع، این کار جایگاه منحصربه‌فردی را در تلاقی مدل‌سازی تحلیلی، کاهش ابعاد داده‌محور و اعمال قوانین پایستگی اشغال می‌کند. با نشان دادن اینکه توابع خاص می‌توانند به طور مؤثر مدل‌های کاهش‌یافته پارامترشده را هنگامی که از طریق مهندسی ویژگی آگاه از فیزیک لنگر انداخته می‌شوند، پارامترسازی کنند، ما راه جدیدی را برای سیستم‌های پیچیده‌ای که قبلاً در برابر تحلیل دقیق مقاومت می‌کردند، باز می‌کنیم. اعتبارسنجی تجربی قوی بر گازی‌سازی زیست‌توده اثبات مفهوم محکمی ارائه می‌دهد، در حالی که ساختار ریاضی کلی، قابلیت تعمیم وسیعی را پیشنهاد می‌دهد. همان‌طور که علوم مهندسی به سمت سیستم‌های پیچیده‌تر با الزامات فزاینده برای کارایی محاسباتی و سازگاری فیزیکی حرکت می‌کنند، رویکردهای ترکیبی از این دست احتمالاً اهمیت بیشتری پیدا خواهند کرد – و این چارچوب نقشه راه ارزشمندی برای توسعه آن‌ها ارائه می‌دهد.

\textbf(پیوست الف: مشتقات تفصیلی)

الف.۱ تبدیل‌های هندسی

تبدیل از متغیرهای چهاربعدی تحلیل همبستگی متعارف EPS = [E₁, E₂, S₁, P₁] به دو ویژگی هندسی (d, K) از طریق یک ساختار هندسی نظام‌مند انجام می‌شود که اطلاعات رابطهای اساسی را حفظ می‌کند در حالی که ابعاد را به طرز چشمگیری کاهش می‌دهد. با شروع از محاسبات مرکزوار برای دو پیکربندی مثلثی، برای مثلث ساختاری Gₛ = ((E₁ + E₂ + S₁)/۳, (y\_(E₁) + y\_(E₂) + y\_(S₁))/۳) که روابط بین ویژگی‌های عنصری و ساختاری سوخت را کد می‌کند، و برای مثلث تقریبی Gₚ = ((E₁ + E₂ + P₁)/۳, (y\_(E₁) + y\_(E₂) + y\_(P₁))/۳) که همبستگی‌های عنصری-تقریبی را ثبت می‌کند، محاسبه می‌کنیم. فاصله بین این مرکزوارها از هندسه اقلیدسی استاندارد به صورت d = √((Gₚˣ - Gₛˣ)² + (Gₚʸ - Gₛʸ)²) به دست می‌آید که چگونگی قرارگیری مشخصه‌های ساختاری و تقریبی سوخت را نسبت به ترکیب عنصری کمّی می‌کند.

ساخت ویژگی مساحت ابتدا نیازمند محاسبه نقطه میانی M = ((Gₚˣ + Gₛˣ)/۲, (Gₚʸ + Gₛʸ)/۲) پاره‌خط اتصال‌دهنده دو مرکزوار است، سپس ارزیابی مساحت علامت‌دار مثلث تشکیل‌شده توسط قاعده E₁E₂ و رأس M از طریق فرمول دترمینان K = (۱/۲)[E₁(y\_(E₂) - y\_M) + E₂(y\_M - y\_(E₁)) + M\_x(y\_(E₁) - y\_(E₂))] انجام می‌شود. این ویژگی‌های هندسی (d, K) ناوردا تحت انتقال و مقیاس‌گذاری یکنواخت ثابت می‌شوند و بازنمایی‌های قوی‌ای ارائه می‌دهند که حتی زمانی که مقادیر مختصات مطلق یا قراردادهای نرمال‌سازی در بین مجموعه‌داده‌ها متفاوت است، به‌خوبی تعریف‌شده باقی می‌مانند. این ویژگی برای انتقال چارچوب به مطالعات تجربی جدید بدون نیاز به واسنجی مجدد ضروری است.

الف.۲ تولید دنباله فیبوناچی

کدگذاری زمانی سلسله‌مراتبی از طریق دنباله‌های فیبوناچی با مقداردهی اولیه دقیقی آغاز می‌شود که اطلاعات خاص سیستم را در ساختار دنباله کد می‌کند. با توجه به یک پارامتر مرکب نرمال‌شده Nⱼ مشتق‌شده از ساختار مبتنی بر بخش که ویژگی‌های هندسی، شرایط عملیاتی و ثابت‌های فیزیکی را یکپارچه می‌کند، سه جمله اول از پروتکل مقداردهی اولیه f₁ = Nⱼ + ۲exp(-Nⱼ) ، f₂ = exp(-Nⱼ)·f₁ و f₃ = f₁ + f₂ پیروی می‌کنند. عوامل میرایی نمایی همزمان اهداف متعددی را دنبال می‌کنند: جلوگیری از رشد نامحدود با گسترش دنباله‌ها به طول‌های بیشتر، معرفی یک مقیاس طول مشخصه متناسب معکوس با Nⱼ که مقایسه بین موارد مختلف را تسهیل می‌کند، و ایجاد اتصالات هموار بین مقداردهی اولیه با انگیزه فیزیکی و بازگشت انتزاعی ریاضی.

جملات بعدی از بازگشت متعارف فیبوناچی f\_n = f\_(n-1) + f\_(n-2) برای n ≥ ۴ پیروی می‌کنند و ساختار سلسله‌مراتبی خود-متشابهی تولید می‌کنند که هم اطلاعات محلی را از طریق جملات اخیر و هم اطلاعات سراسری را از طریق مجموع‌های تجمعی کد می‌کند. انتخاب طول دنباله شامل متعادل کردن ملاحظات رقابتی است: دنباله‌های کوتاه‌تر (کمتر از ۱۰ جمله) ممکن است به طور ناکافی ساختار زمانی چندمقیاسی را ثبت کنند، در حالی که دنباله‌های طولانی‌تر (بیش از ۲۰ جمله) بار محاسباتی را بدون سود اطلاعاتی قابل توجه با توجه به بزرگی‌های نمایی-رو به رشد جمله‌ها معرفی می‌کنند. بهینه‌سازی تجربی روی مجموعه داده آموزشی مشخص کرد که دنباله‌های ۱۵ جمله‌ای عملکرد بهینه را ارائه می‌دهند، به اندازه کافی طولانی برای کدگذاری ساختار سلسله‌مراتبی اساسی در حالی که از نظر محاسباتی برای کاربردهای بلادرنگ که نیاز به تولید هزاران دنباله در ثانیه دارند، قابل مدیریت باقی می‌مانند.

الف.۳ ساخت ماتریس مشخصه از طریق حاصلضرب خارجی

ادغام اطلاعات اثر شبکه عصبی در ماتریس‌های مشخصه از طریق نرمال‌سازی دقیق و ترکیب متقارن که درستی ریاضی را تضمین می‌کند، انجام می‌شود. با توجه به بردارهای امتیاز S\_AB ∈ ℝ⁹ و S\_BA ∈ ℝ⁹ که اثرهای حاصل از پردازش دنباله رو به جلو و معکوس در سه تابع فعال‌سازی را جمع می‌کنند، ابتدا به طول اقلیدسی واحد نرمال می‌کنیم: Ŝ\_AB = S\_AB / ‖S\_AB‖ و Ŝ\_BA = S\_BA / ‖S\_BA‖ که در آن ‖·‖ نرم L2 استاندارد را نشان می‌دهد. این نرمال‌سازی برای اطمینان از اینکه ساخت حاصلضرب خارجی بعدی بر ساختار همبستگی تأکید می‌کند به‌جای اینکه تحت سلطه تفاوت‌های بزرگی باشد که ممکن است بازتاب انتخاب‌های دلخواه آموزش شبکه عصبی یا دانه‌های تصادفی مقداردهی اولیه باشند، ضروری ثابت می‌شود.

ماتریس مشخصه از حاصلضرب خارجی متقارن S\_c = (۱/۲)(Ŝ\_AB ⊗ Ŝ\_BA\textasciicircum{}T + Ŝ\_BA ⊗ Ŝ\_AB\textasciicircum{}T) ظهور می‌یابد که در آن ⊗ حاصلضرب خارجی را نشان می‌دهد و عملگر ترانهاده، تراز بعدی مناسب را تضمین می‌کند. متقارن‌سازی—میانگین گرفتن حاصلضرب خارجی با ترانهاده آن—تضمین می‌کند که S\_c متقارن نیمه معین مثبت با مقادیر ویژه حقیقی غیرمنفی می‌شود، یک ویژگی ریاضی که تضمین می‌کند تفسیرهای فیزیکی بعدی با قیود ترمودینامیکی سازگار باقی می‌مانند. ماتریس ۹×۹ حاصل از طریق تبدیل به بردار و سپس تاکردن مجدد به شکل ۳×۳ تغییر شکل می‌یابد، اطلاعات ساختاری اساسی کدگذاری شده در طیف مقادیر ویژه و جهت‌های بردارهای ویژه را حفظ می‌کند در حالی که ابعاد را به سطوح قابل مدیریت برای پارامترسازی انتگرال بیضوی پایین‌دستی فشرده می‌کند.

الف.۴ تحلیل همگرایی میانگین حسابی-هندسی

همگرایی درجه دوم نشان داده شده توسط الگوریتم میانگین حسابی-هندسی—پایه محاسباتی که کارایی فوق‌العاده چارچوب ما را ممکن می‌سازد—شایسته توصیف ریاضی تفصیلی برای درک این است که چرا صرفاً پنج تکرار به طور جهانی برای دقت ماشین بدون توجه به مقادیر پارامتر کافی است. با شروع از مقداردهی اولیه a₀ = ۱ و b₀ = √(۱ - m) که در آن m ∈ [۰, ۱) پارامتر انتگرال بیضوی را نشان می‌دهد، طرح تکرار بین میانگین حسابی a\_(n+1) = (a\_n + b\_n)/۲ و میانگین هندسی b\_(n+1) = √(a\_n b\_n) متناوب می‌شود. هر دو دنباله به حد مشترک AGM(۱, √(۱ - m)) همگرا می‌شوند که از آن انتگرال بیضوی کامل از طریق K(m) = π/(۲ · AGM(۱, √(۱ - m))) به دست می‌آید.

کران خطای حاکم بر نرخ همگرایی به شکل |a\_n - AGM(۱, √(۱ - m))| ≤ C · ۲\textasciicircum{}(-۲ⁿ) برای مقداری ثابت C است که به پارامتر m بستگی دارد اما در سراسر دامنه همگرایی m ∈ [۰, ۰.۹۹] مرتبط با کاربردهای ما یکنواخت باقی می‌ماند. این واپاشی دوگانه نمایی—که در آن خود توان به طور نمایی رشد می‌کند—بیان می‌کند که ارقام صحیح با هر تکرار دوبرابر می‌شوند: تکرار ۱ تقریباً ۲ رقم اعشار صحیح می‌دهد، تکرار ۲، ۴ رقم، تکرار ۳، ۸ رقم، تکرار ۴، ۱۶ رقم که از نیازهای دقت مضاعف (تقریباً ۱۵-۱۶ رقم اعشار برای ممیز شناور دقت مضاعف IEEE ۷۵۴) فراتر می‌رود، و تکرار ۵ دقت کامل ماشین با دقت مانتیس ۵۳ بیتی را به دست می‌آورد. این کارایی قابل توجه—دستیابی به ارزیابی تابع خاص با هزینه قابل مقایسه با عملیات حسابی ابتدایی مانند ضرب یا جذر—ثابت می‌شود که برای توانایی چارچوب ما در ارائه شتاب هزار برابری در حالی که دقت ریاضی از طریق محاسبه دقیق تابع خاص به جای تقریب‌های چندجمله‌ای که دقت آن‌ها برای مقادیر پارامتری شدید کاهش می‌یابد، حفظ می‌شود، محوری است.

\end{document}
